% Guía del segundo parcial — Programación 3
% Compilar con: pdflatex guia-segundo-parcial.tex (dos veces)
% =====================================================================
%  Estilo de la guía del segundo parcial — Programación 3
%  Tokens de diseño:
%    tinta      #1E2A4A   texto y títulos
%    gris       #5B6375   texto secundario, notas al margen
%    papel      #FFFFFF
%    acento por parte (identifica la parte, no decora):
%      dyc #2F6F8F  pd #7A4E9C  flujo #2E7D5B  np #A33B3B
%  Tipografía: Libertinus Serif (cuerpo + matemática), Libertinus Sans (títulos)
%  Layout: bloque de texto a la izquierda, margen derecho ancho para
%          referencias a parciales/exámenes ("visto en").
% =====================================================================
\documentclass[11pt,a4paper]{article}

\usepackage[T1]{fontenc}
\usepackage[utf8]{inputenc}
\usepackage[spanish,es-noshorthands,es-nodecimaldot]{babel}
\usepackage{libertinus}
\usepackage{amsmath,amssymb,amsthm,mathtools}
\usepackage{libertinust1math}
\usepackage[a4paper,left=2.2cm,right=6.0cm,top=2.4cm,bottom=2.6cm,
            marginparwidth=4.1cm,marginparsep=0.7cm]{geometry}
\usepackage[dvipsnames]{xcolor}
\usepackage{microtype}
\usepackage{enumitem}
\usepackage{booktabs,tabularx,array}
\usepackage{tikz}
\usetikzlibrary{arrows.meta,positioning,shapes,calc,fit,backgrounds}
\usepackage[noend]{algpseudocode}
\usepackage{algorithm}
\usepackage[most]{tcolorbox}
\usepackage{titlesec}
\usepackage{fancyhdr}
\usepackage{marginnote}
\usepackage[hidelinks]{hyperref}
\usepackage{etoolbox}

% ---------- color ----------
\definecolor{tinta}{HTML}{1E2A4A}
\definecolor{gris}{HTML}{5B6375}
\definecolor{grisclaro}{HTML}{E6E8EE}
\definecolor{grafos}{HTML}{9A6A0F}
\definecolor{dyc}{HTML}{2F6F8F}
\definecolor{pd}{HTML}{7A4E9C}
\definecolor{flujo}{HTML}{2E7D5B}
\definecolor{np}{HTML}{A33B3B}
\colorlet{acento}{tinta}           % se redefine al abrir cada parte
\color{tinta}

% ---------- tipografía ----------
\linespread{1.12}
\setlength{\parindent}{0pt}
\setlength{\parskip}{0.45em}
\newcommand{\sans}{\sffamily}

% ---------- títulos ----------
\titleformat{\section}{\sans\Large\bfseries\color{acento}}{\thesection}{0.8em}{}
\titleformat{\subsection}{\sans\large\bfseries\color{tinta}}{\thesubsection}{0.7em}{}
\titleformat{\subsubsection}{\sans\normalsize\bfseries\color{tinta}}{}{0em}{}
\titlespacing*{\section}{0pt}{1.6em}{0.6em}
\titlespacing*{\subsection}{0pt}{1.2em}{0.4em}
\titlespacing*{\subsubsection}{0pt}{0.9em}{0.25em}
\renewcommand{\thesection}{\arabic{section}}

% ---------- encabezado ----------
\pagestyle{fancy}
\fancyhf{}
\fancyheadoffset[R]{4.8cm}
\renewcommand{\headrulewidth}{0pt}
\newcommand{\parteactual}{}
\fancyhead[L]{\sans\footnotesize\color{gris}Programación 3 — guía del segundo parcial}
\fancyhead[R]{\sans\footnotesize\color{acento}\parteactual}
\fancyfoot[R]{\sans\footnotesize\color{gris}\thepage}
\fancyfootoffset[R]{4.8cm}

% ---------- portada de cada parte: la fórmula central es el protagonista ----------
% \parte{nombre-color}{Título}{fórmula (modo matemático display)}{bajada de una o dos líneas}
\newcommand{\parte}[4]{%
  \clearpage
  \colorlet{acento}{#1}%
  \renewcommand{\parteactual}{#2}%
  \thispagestyle{empty}%
  \begin{tikzpicture}[remember picture,overlay]
    \fill[#1] (current page.north west) rectangle ([xshift=0.9cm]current page.south west);
  \end{tikzpicture}%
  \vspace*{3.2cm}
  {\sans\Huge\bfseries\color{#1}#2\par}
  \vspace{2.2cm}
  \noindent\makebox[\textwidth][l]{\begin{minipage}{\dimexpr\textwidth+4.8cm\relax}
    \color{tinta}\fontsize{26}{34}\selectfont
    \[ #3 \]
  \end{minipage}}
  \vspace{1.8cm}

  \begin{minipage}{0.9\textwidth}\large\linespread{1.3}\selectfont\color{gris}#4\end{minipage}\par
  \addcontentsline{toc}{part}{#2}%
  \clearpage
}

% ---------- recuadros con significado ----------
% receta: cómo se redacta la respuesta en el parcial
\newtcolorbox{receta}[1][]{enhanced,breakable,sharp corners,boxrule=0pt,
  frame hidden,colback=acento!6,borderline west={3pt}{0pt}{acento},
  left=10pt,right=8pt,top=6pt,bottom=6pt,
  title={\sans\bfseries Receta para el parcial\ifstrempty{#1}{}{: #1}},
  coltitle=acento,colbacktitle=acento!6,attach title to upper={\par\smallskip}}
% trampa: errores que se corrigen siempre
\newtcolorbox{trampa}[1][]{enhanced,breakable,sharp corners,boxrule=0pt,
  frame hidden,colback=white,borderline west={1.2pt}{0pt}{np},
  left=10pt,right=8pt,top=4pt,bottom=4pt,
  title={\sans\bfseries\color{np}Errores que restan puntos\ifstrempty{#1}{}{: #1}},
  coltitle=np,colbacktitle=white,attach title to upper={\par\smallskip}}
% problema: enunciado de un problema clásico o de examen
\newtcolorbox{problema}[1]{enhanced,breakable,sharp corners,boxrule=0.4pt,
  colframe=grisclaro,colback=white,left=8pt,right=8pt,top=5pt,bottom=5pt,
  fontupper=\small,title={\sans\bfseries #1},coltitle=tinta,colbacktitle=grisclaro}
% resultado: enunciado clave del teórico (lo que se cita "por (x.y)")
\newtcolorbox{resultado}[1][]{enhanced,breakable,sharp corners,boxrule=0pt,frame hidden,
  colback=grisclaro!45,left=10pt,right=8pt,top=6pt,bottom=6pt,
  before upper={\ifstrempty{#1}{}{{\sans\bfseries #1}\par\smallskip}}}

% ---------- notas al margen: dónde se preguntó ----------
% \visto{2P 2019 Ej. 2: ladrón de casas contiguas}
\newcommand{\visto}[1]{\marginnote{\raggedright\sans\footnotesize\color{gris}%
  {\color{acento}\rule[0.15em]{0.9em}{1.4pt}}\hspace{0.35em}#1}}
% nota al margen libre
\newcommand{\margen}[1]{\marginnote{\raggedright\footnotesize\color{gris}#1}}

% ---------- teoremas ----------
\newtheoremstyle{guia}{0.6em}{0.6em}{\itshape}{}{\sans\bfseries}{.}{0.5em}{}
\theoremstyle{guia}
\newtheorem{teorema}{Teorema}[section]
\newtheorem{lema}[teorema]{Lema}
\newtheorem{proposicion}[teorema]{Proposición}
\theoremstyle{definition}
\newtheorem{definicion}[teorema]{Definición}
\newtheorem*{ejemplo}{Ejemplo}
\renewcommand{\qedsymbol}{\color{acento}$\blacksquare$}

% ---------- pseudocódigo en español ----------
\floatname{algorithm}{Algoritmo}
\renewcommand{\algorithmicrequire}{\textbf{Entrada:}}
\renewcommand{\algorithmicensure}{\textbf{Salida:}}
\algrenewcommand\algorithmicwhile{\textbf{mientras}}
\algrenewcommand\algorithmicdo{\textbf{hacer}}
\algrenewcommand\algorithmicfor{\textbf{para}}
\algrenewcommand\algorithmicforall{\textbf{para todo}}
\algrenewcommand\algorithmicif{\textbf{si}}
\algrenewcommand\algorithmicthen{\textbf{entonces}}
\algrenewcommand\algorithmicelse{\textbf{sino}}
\algrenewcommand\algorithmicreturn{\textbf{retornar}}
\algrenewcommand\algorithmicfunction{\textbf{función}}
\algrenewcommand\algorithmicprocedure{\textbf{procedimiento}}

% ---------- macros ----------
\newcommand{\Oh}{\mathcal{O}}
\newcommand{\OPT}{\mathrm{OPT}}
\newcommand{\NP}{\mathsf{NP}}
\newcommand{\Pclass}{\mathsf{P}}
\newcommand{\coNP}{\mathsf{co}\text{-}\mathsf{NP}}
\newcommand{\redp}{\le_{p}}
\newcommand{\prob}[1]{\textsc{#1}}

\setlist{itemsep=0.15em,topsep=0.3em,leftmargin=1.4em}
\setcounter{tocdepth}{1}


\titleformat{\part}[display]{\sans\bfseries}{}{0pt}{}
\addto\captionsspanish{\renewcommand{\contentsname}{Contenido}}

\begin{document}

% ---------------- portada ----------------
\thispagestyle{empty}
\begin{tikzpicture}[remember picture,overlay]
  \foreach \c [count=\i from 0] in {grafos,dyc,pd,flujo,np}
    \fill[\c] ([xshift=\i*0.34cm]current page.north west) rectangle
              ([xshift=\i*0.34cm+0.34cm]current page.south west);
\end{tikzpicture}
\vspace*{4.5cm}

{\sans\fontsize{40}{44}\selectfont\bfseries\color{tinta}Segundo parcial\par}
\vspace{0.6em}
{\sans\LARGE\color{gris}Programación 3 — guía de estudio\par}
\vspace{3cm}

\begin{minipage}{0.92\textwidth}
\large\linespread{1.35}\selectfont
Todo lo que se evalúa desde Dijkstra en adelante: greedy en grafos,
dividir y conquistar, programación dinámica, flujo en redes y reducciones. Cada tema está acá porque se
preguntó algo de él en algún parcial o examen. Al margen figura dónde, y al
final de cada parte hay un mapa con todas las apariciones.
\end{minipage}

\vfill
{\sans\small\color{gris}
\begin{tabular}{@{}l@{\hspace{0.8em}}l@{}}
\textcolor{grafos}{\rule{1.6em}{0.7em}} & Greedy en grafos\\[2pt]
\textcolor{dyc}{\rule{1.6em}{0.7em}} & Dividir y conquistar\\[2pt]
\textcolor{pd}{\rule{1.6em}{0.7em}} & Programación dinámica\\[2pt]
\textcolor{flujo}{\rule{1.6em}{0.7em}} & Flujo en redes\\[2pt]
\textcolor{np}{\rule{1.6em}{0.7em}} & Reducciones y NP-completitud
\end{tabular}}
\clearpage

% ---------------- cómo leer ----------------
\section*{Cómo usar esta guía}
Cada parte abre con la fórmula que la resume. Adentro, cada tema tiene la
teoría que hay que saber citar, los problemas que se tomaron y tres
tipos de recuadro:

\begin{resultado}[Resultado del teórico]
Lo que se puede citar en el parcial sin demostrar, con su número del libro de
Kleinberg y Tardos.
\end{resultado}

\begin{receta}
Cómo estructurar la respuesta para que tenga todo lo que piden los criterios
de corrección.
\end{receta}

\begin{trampa}
Errores que aparecen una y otra vez en las correcciones.
\end{trampa}

Las referencias al margen indican dónde se preguntó cada cosa\visto{2P 2019
Ej. 2: así se ve una referencia}. «2P» es segundo parcial y «Ex» o
«Examen» es examen, con mes y año.

\vspace{1.5em}
\tableofcontents

% =====================================================================
%  Parte 0 — Greedy en grafos (Kleinberg–Tardos, cap. 4, secciones 4.4–4.6)
%  Fragmento: se incluye con \input desde el documento maestro.
% =====================================================================

% ---------- macros y estilos locales (prefijo gr) ----------
\newcommand{\grpred}{\mathrm{pred}}
\newcommand{\grop}[1]{\textsc{#1}}
\tikzset{
  grv/.style={circle,draw=tinta,fill=white,inner sep=1pt,minimum size=5.6mm,font=\small},
  gra/.style={draw=gris,semithick},
  grp/.style={draw=grafos,line width=1.7pt},
  grx/.style={draw=gris!55,semithick,dashed},
  grd/.style={-{Stealth[length=2.1mm]},draw=gris,semithick},
  grdp/.style={-{Stealth[length=2.3mm]},draw=grafos,line width=1.3pt},
  grl/.style={font=\scriptsize,fill=white,inner sep=0.8pt,text=tinta},
  grt/.style={font=\footnotesize\sffamily,text=gris},
}

\parte{grafos}{Greedy en grafos}
  {d(v)\;=\;\min_{e=(u,v):\,u\in S}\; d(u)+\ell_e}
  {Se arma la solución de a un paso, agregando siempre lo más barato que cruza la frontera de lo que ya está construido. Dijkstra lo hace con distancias, Prim y Kruskal con aristas. Las demostraciones son por inducción sobre ese conjunto o por intercambio a través de un corte.}

% =====================================================================
\section{Qué se pregunta y cómo se contesta}
% =====================================================================

Estos temas son las secciones 4.4 a 4.6 del libro (KT cap.~4). Entre 2018 y 2025 se evaluaron en el \emph{primer} parcial, y antes de 2017 en el segundo parcial o en la parte teórica de los exámenes. Aparecen de tres maneras:
\begin{itemize}
  \item \textbf{Teórico.} Enunciar y demostrar la propiedad de corte o la de ciclo, la correctitud de Dijkstra (4.14) o la unicidad del MST. También aparecen verdaderos o falsos con contraejemplo.
  \item \textbf{Ejecución a mano.} Correr Dijkstra, Prim o Kruskal sobre un grafo chico. La pauta exige el orden exacto de elección y el costo en cada paso.
  \item \textbf{Modelado.} Un problema contado con robots, caminos o cables, que se reduce a caminos mínimos o a MST. Hay que reescribir el algoritmo del libro y probar una cota como $\Oh(m\log n)$.
\end{itemize}

En los exámenes viejos se pedía además la \emph{forma general} de la técnica greedy y su correspondencia con Prim, Kruskal o Dijkstra.\visto{2P 2012 Ej.~2(a): esquema greedy y Kruskal}

\begin{algorithmic}[1]
\Function{Greedy}{$C$} \Comment{$C$: conjunto de candidatos}
  \State $S\gets\emptyset$
  \While{$S$ no es solución \textbf{y} $C\neq\emptyset$}
    \State $x\gets$ \Call{Seleccionar}{$C$}; $C\gets C\setminus\{x\}$ \Comment{óptimo local}
    \If{$S\cup\{x\}$ es factible} \State $S\gets S\cup\{x\}$ \EndIf
  \EndWhile
  \State \Return $S$ (si es solución)
\EndFunction
\end{algorithmic}

\begin{center}
\small
\begin{tabularx}{\textwidth}{@{}>{\raggedright\arraybackslash}p{1.7cm}>{\raggedright\arraybackslash}X>{\raggedright\arraybackslash}X>{\raggedright\arraybackslash}p{2.6cm}@{}}
\toprule
 & Candidatos & Selección & Factible si\dots \\
\midrule
Kruskal & aristas & la de menor costo & no forma ciclo \\
Prim & aristas entre $U$ y $V\setminus U$ & la de menor costo & siempre \\
Dijkstra & nodos fuera de $S$ & el de menor $d'(v)$ & siempre \\
\bottomrule
\end{tabularx}
\end{center}

La solución se alcanza al tener $n-1$ aristas (Kruskal), al llegar a $U=V$ (Prim) o al llegar a $S=V$ (Dijkstra). En las preguntas del tipo ``elija tres algoritmos y diga qué resuelven y con qué técnica'' alcanza con una línea por algoritmo. Dijkstra resuelve caminos de costo mínimo desde un origen con costos no negativos, y es greedy. Prim y Kruskal dan un árbol de cubrimiento de costo mínimo de un grafo no dirigido con costos, y también son greedy. Floyd y Bellman–Ford, en cambio, son programación dinámica.\visto{Examen jul.~2016 Ej.~4 y feb.~2018 P.~4: qué resuelve}

% =====================================================================
\section{Caminos mínimos: el algoritmo de Dijkstra}
% =====================================================================

Sea $G=(V,E)$ un grafo dirigido (o no dirigido, con cada arista en los dos sentidos), con longitudes $\ell_e\ge 0$ y un origen $s$. Se busca, para cada $v$, la distancia $d(v)$, que es la longitud mínima de un camino de $s$ a $v$, y un camino que la realice.

\subsection{El algoritmo}

Dijkstra mantiene un conjunto $S$ de nodos ya \emph{resueltos}, cuya distancia $d(u)$ se conoce. En cada paso mira las aristas que salen de $S$ y elige el nodo $v\notin S$ que minimiza
\[ d'(v)=\min_{e=(u,v):\,u\in S}\; d(u)+\ell_e , \]
es decir, el nodo al que se llega más barato saliendo de $S$ por una sola arista.

\begin{algorithmic}[1]
\Function{Dijkstra}{$G,\ell,s$}
  \State $S\gets\{s\}$; $d(s)\gets 0$
  \While{$S\neq V$}
    \State elegir $v\notin S$, con alguna arista desde $S$, que minimice $d'(v)$
    \State $S\gets S\cup\{v\}$; $d(v)\gets d'(v)$; $\grpred(v)\gets$ el $u$ que realiza el mínimo
  \EndWhile
  \State \Return $d$, $\grpred$
\EndFunction
\end{algorithmic}

El camino $P_v$ se obtiene siguiendo $\grpred$ desde $v$ hasta $s$. Las aristas $(\grpred(v),v)$ forman un árbol con raíz $s$, el \emph{árbol de caminos mínimos}.

En el examen de febrero de 2026 preguntaron exactamente qué propiedad cumple el próximo vértice que entra a $S$. La respuesta es la de la línea 4: entre todas las aristas $(u,v)$ con $u\in S$, $v\in V\setminus S$, se elige $v$ para el que $d(u)+\ell_{(u,v)}$ es mínimo.\visto{Examen feb.~2026 Ej.~1(c): próximo vértice de $S$}

\subsection{Correctitud}

\begin{resultado}[(4.14)]
En cualquier momento de la ejecución, para cada $u\in S$, el camino $P_u$ es un camino mínimo de $s$ a $u$. En consecuencia, al terminar, $d(u)$ es la distancia de $s$ a $u$ para todo $u$.
\end{resultado}

El paso inductivo de esta demostración se pidió aislado en el examen de julio de 2019: dados $S$ y caminos mínimos $P_u$ para $u\in S$, probar que la arista que minimiza $d(u)+\ell_{(u,v)}$ da un camino mínimo a $v$.\visto{Examen jul.~2019 Ej.~1: el argumento central}

\begin{proof}
Por inducción en $|S|$. Si $|S|=1$, $S=\{s\}$ y $d(s)=0$ es la distancia, porque las longitudes son no negativas.

Supongamos que vale para $|S|=k\ge1$ y que el algoritmo agrega $v$ usando la arista $(u,v)$, con $u\in S$. El camino $P_v$ es $P_u$ seguido de $(u,v)$, y su longitud es $d(u)+\ell_{(u,v)}=d'(v)$. Falta ver que ningún camino de $s$ a $v$ es más corto.

Sea $P$ un camino cualquiera de $s$ a $v$. Como $s\in S$ y $v\notin S$, $P$ sale de $S$ en algún punto. Sea $y$ el primer nodo de $P$ que no está en $S$ y sea $x\in S$ el nodo anterior. Entonces
\begin{align*}
\ell(P) &\ge \ell(P[s,x]) + \ell_{(x,y)} && \text{el tramo de $y$ a $v$ tiene longitud $\ge 0$}\\
        &\ge d(x)+\ell_{(x,y)} && \text{hipótesis inductiva: $P_x$ es mínimo}\\
        &\ge d'(y) && \text{definición de $d'(y)$}\\
        &\ge d'(v) && \text{el algoritmo eligió $v$ como mínimo.}
\end{align*}
Por lo tanto $\ell(P)\ge d'(v)=\ell(P_v)$, y $P_v$ es mínimo.
\end{proof}

La no negatividad se usa en un solo lugar, la primera desigualdad, y la corrección pide señalarlo. Si hubiera aristas negativas, el tramo de $y$ a $v$ podría ``devolver'' longitud.

\subsection{Por qué falla con longitudes negativas}

Contraejemplo de la solución oficial de julio de 2019: aristas $(s,v)$ de longitud $1$, $(s,y)$ de longitud $2$ e $(y,v)$ de longitud $-2$. Con $S=\{s\}$, el mínimo de $d(u)+\ell_e$ es $1$ y se alcanza en $(s,v)$, así que Dijkstra fija $d(v)=1$. Sin embargo, el camino $s\,y\,v$ tiene longitud $0$.\visto{Examen jul.~2019 Ej.~1(b): con negativos no vale}

\begin{center}
\begin{tikzpicture}[scale=1]
  \node[grv] (s) at (0,0) {$s$};
  \node[grv] (y) at (1.6,1.0) {$y$};
  \node[grv] (v) at (3.2,0) {$v$};
  \draw[grdp] (s)--node[grl,below]{$1$}(v);
  \draw[grd] (s)--node[grl,above left]{$2$}(y);
  \draw[grd] (y)--node[grl,above right]{$-2$}(v);
  \node[grt] at (1.6,-0.65) {Dijkstra fija $d(v)=1$, pero $\ell(s\,y\,v)=0$};
\end{tikzpicture}
\end{center}

\begin{trampa}[longitudes negativas]
\begin{itemize}
  \item Para ``arreglar'' los negativos sumándole una constante $M$ a todas las aristas. No sirve: un camino de $k$ aristas aumenta $kM$, así que se favorecen los caminos con pocas aristas y el óptimo cambia. (En MST sí funciona, porque todos los árboles tienen $n-1$ aristas; ver sección~\ref{gr:transf}.)
  \item Decir que Dijkstra falla ``si hay ciclos negativos''. Falla con una sola arista negativa, aunque no haya ciclos. Para negativos sin ciclos negativos se usa Bellman–Ford (programación dinámica), y en un DAG alcanza con el orden topológico (sección~\ref{gr:dag}).
\end{itemize}
\end{trampa}

\subsection{Implementación y orden}

\begin{resultado}[(4.15)]
Con una cola de prioridad, Dijkstra se implementa en tiempo $\Oh(m)$ más el costo de $n$ operaciones \grop{ExtraerMin} y $m$ operaciones \grop{CambiarClave}. Con un heap binario, cada una cuesta $\Oh(\log n)$ y el total es $\Oh(m\log n)$.
\end{resultado}

\begin{algorithmic}[1]
\Function{Dijkstra}{$G,\ell,s$} \Comment{versión con heap}
  \ForAll{$v\in V$} \State $d[v]\gets\infty$; $\grpred[v]\gets\mathrm{nil}$ \EndFor
  \State $d[s]\gets 0$; $Q\gets$ heap de mínimo con todos los nodos, clave $d[v]$
  \While{$Q\neq\emptyset$}
    \State $u\gets$ \Call{ExtraerMin}{$Q$} \Comment{$u$ entra a $S$}
    \ForAll{arista $(u,v)$ con $v\in Q$}
      \If{$d[u]+\ell_{(u,v)}<d[v]$}
        \State $d[v]\gets d[u]+\ell_{(u,v)}$; $\grpred[v]\gets u$; \Call{CambiarClave}{$Q,v,d[v]$}
      \EndIf
    \EndFor
  \EndWhile
  \State \Return $d,\grpred$
\EndFunction
\end{algorithmic}

\textbf{Análisis.} Inicializar cuesta $\Oh(n)$. Cada nodo sale del heap una sola vez, lo que da $n$ \grop{ExtraerMin} a $\Oh(\log n)$ cada una. Cada arista $(u,v)$ se examina una sola vez, cuando $u$ sale del heap, y provoca a lo sumo un \grop{CambiarClave} de $\Oh(\log n)$. Para \grop{CambiarClave} hace falta un arreglo de posiciones en el heap. El total es $\Oh((n+m)\log n)$, que es $\Oh(m\log n)$ si todos los nodos son alcanzables desde $s$, porque entonces $m\ge n-1$.

En los exámenes de antes de 2017 la versión ``del teórico'' era la de matriz de costos: un arreglo $D$ que en cada una de las $n-1$ iteraciones se recorre entero para buscar el mínimo y actualizar. Eso da $\Theta(n^2)$, que es mejor que el heap en grafos densos ($m\approx n^2$).\visto{2P 2014 Ej.~2.5 y Examen jul.~2011 Ej.~3: escribir Dijkstra}

\begin{trampa}[implementación]
\begin{itemize}
  \item Decir ``$\Oh(m\log n)$ por el heap'' sin contar \emph{cuántas} operaciones se hacen. La pauta pide $n$ extracciones y a lo sumo $m$ cambios de clave, cada uno $\Oh(\log n)$, y justificar que $n=\Oh(m)$.
  \item Olvidar el arreglo de predecesores cuando se pide \emph{el camino} y no solo su costo.
  \item Actualizar nodos que ya están en $S$. Su distancia es definitiva por (4.14).
\end{itemize}
\end{trampa}

\subsection{Dijkstra a mano}

Grafo del examen de diciembre de 2014 (no dirigido), con origen $2$: $c(1,2)=5$, $c(1,3)=3$, $c(1,5)=1$, $c(2,3)=2$, $c(2,4)=3$, $c(3,4)=1$, $c(4,5)=5$. La solución oficial muestra en cada paso $S$, $V\setminus S$, el arreglo $D$ y el mínimo elegido.\visto{Examen dic.~2014 Ej.~3: Dijkstra desde el vértice 2}

\begin{center}
\begin{tikzpicture}[scale=1.05]
  \node[grv] (2) at (0,0) {$2$};
  \node[grv] (1) at (2.2,1.3) {$1$};
  \node[grv] (3) at (2.2,-1.1) {$3$};
  \node[grv] (4) at (0,-2.2) {$4$};
  \node[grv] (5) at (4.4,-0.4) {$5$};
  \draw[grp] (2)--node[grl,above left]{$5$}(1);
  \draw[grp] (2)--node[grl,above right]{$2$}(3);
  \draw[gra] (1)--node[grl,right]{$3$}(3);
  \draw[grp] (2)--node[grl,left]{$3$}(4);
  \draw[gra] (3)--node[grl,below right]{$1$}(4);
  \draw[grp] (1)--node[grl,above right]{$1$}(5);
  \draw[gra] (4) to[bend right=28] node[grl,below right]{$5$}(5);
  \node[grt,above=1pt of 1] {$d=5$};
  \node[grt,right=1pt of 3] {$d=2$};
  \node[grt,left=1pt of 4] {$d=3$};
  \node[grt,right=1pt of 5] {$d=6$};
  \node[grt,left=1pt of 2] {$d=0$};
\end{tikzpicture}
\end{center}

\begin{center}
\small
\begin{tabularx}{\textwidth}{@{}c>{\raggedright\arraybackslash}p{1.9cm}>{\raggedright\arraybackslash}X>{\raggedright\arraybackslash}p{2.5cm}@{}}
\toprule
Paso & Elegido & Actualización de $D$ & $D=[1,3,4,5]$ \\
\midrule
0 & $2$, $d=0$ & inicial: aristas que salen de $2$ & $[5,2,3,\infty]$ \\
1 & $3$, $d=2$ & $D[1]=\min\{5,2+3\}=5$, $D[4]=\min\{3,2+1\}=3$ & $[5,\cdot,3,\infty]$ \\
2 & $4$, $d=3$ & $D[5]=\min\{\infty,3+5\}=8$ & $[5,\cdot,\cdot,8]$ \\
3 & $1$, $d=5$ & $D[5]=\min\{8,5+1\}=6$, $\grpred(5)=1$ & $[\cdot,\cdot,\cdot,6]$ \\
4 & $5$, $d=6$ & $S=V$, termina & \\
\bottomrule
\end{tabularx}
\end{center}

En el paso 1 hay un empate en $D[4]$: $2+1=3$ no mejora estrictamente, así que $\grpred(4)$ sigue siendo $2$. Con las dos opciones el camino es mínimo. Las aristas destacadas forman el árbol de caminos mínimos.

\begin{receta}[ejecutar Dijkstra]
\begin{enumerate}
  \item Una fila por iteración, con el vértice elegido \emph{en el orden en que el algoritmo lo elige} y su costo. En el 1P 2025 la respuesta se consideró incorrecta si faltaba alguna de las dos cosas.
  \item Mostrá las actualizaciones con la cuenta explícita $\min\{D[v],\,d(u)+\ell_{(u,v)}\}$. Es la ``justificación de cada paso'' que piden.
  \item Si se piden caminos, llevá la columna de predecesores y al final dibujá el árbol.
\end{enumerate}
\end{receta}

Otros enunciados de ejecución: febrero de 2010 (``costos y predecesores de cada nodo en cada iteración'') y el 1P de 2025, donde la respuesta esperada era la secuencia $x=3$, $z=5$, $y=6$, $w=10$.\visto{Examen feb.~2010 Ej.~3; 1P 2025 Ej.~1(e)}

% =====================================================================
\section{Variantes y modelado con caminos mínimos}
% =====================================================================

Casi todos los ejercicios de desarrollo con Dijkstra siguen el mismo patrón: modelar como grafo, decidir qué es la longitud de una arista, correr Dijkstra (a veces dos veces, o en el grafo invertido) y leer la respuesta. La pauta pide \emph{reescribir} el algoritmo del libro aunque se use tal cual.

\subsection{Grillas y robots}

En los parciales de 2021 y de 2024 se planteó una grilla de $n$ celdas libres u ocupadas, y un robot que se mueve a las 8 vecinas. Moverse en diagonal cuesta $\sqrt2$ y en línea recta $1$. Se pide el camino más corto en $\Oh(n\log n)$.\visto{1P 2021 Ej.~2 y 1P 2024 Ej.~2: robot en grilla}

\begin{problema}{1P 2021, Ej.~2 (y 1P 2024, Ej.~2): camino más corto en una grilla}
Matriz de celdas libres/ocupadas, posición inicial y destino. Dar un algoritmo para un camino transitable más corto y probar que es $\Oh(n\log n)$, con $n$ la cantidad de entradas.

\smallskip
\textbf{Solución.} Un nodo por celda libre y una arista entre celdas libres vecinas, de costo $1$ o $\sqrt2$. Dijkstra desde la celda inicial con predecesores, y al final se reconstruye el camino siguiendo $\grpred$ desde el destino. \emph{Orden:} hay $\Oh(n)$ nodos y cada uno tiene a lo sumo $8$ vecinos, así que $m=\Oh(n)$. Son $n$ extracciones del heap a $\Oh(\log n)$ y a lo sumo $8$ actualizaciones de $\Oh(\log n)$ por extracción. Reconstruir el camino es $\Oh(n)$. El total es $\Oh(n\log n)$.
\end{problema}

Si todos los pasos cuestan lo mismo, no hace falta Dijkstra: BFS da las distancias en $\Oh(m+n)$. Así se resolvió el examen de julio de 2020 (el robot más cercano a la zona $z$ en una grilla de $M\times M$ en $\Oh(M^2)$, con un BFS desde $z$). En el viñedo del 1P 2025 la pauta aclara que Dijkstra ``no cumple el orden requerido'' $\Oh(f\times c)$. Lo que hacía falta allí era un BFS con todas las plantas infectadas como fuentes iniciales.\visto{Examen dic.~2015 Ej.~3 y 1P 2025 Ej.~2: BFS alcanza}

\subsection{Reconstruir el camino y cortar antes}

En diciembre de 2016, además de las distancias, se pedía poder imprimir el recorrido y detener el algoritmo apenas se tuviera lo necesario.\visto{Examen dic.~2016 Probl.~A: predecesores y corte}

\begin{problema}{Examen dic.~2016, Probl.~A: recorridos a las paradas}
Matriz $C$ de tiempos por cuadra ($\infty$ si no hay cuadra), casa $K$ y conjunto $P$ de paradas. (I) Construir una estructura $E$ que permita obtener el recorrido más rápido de $K$ a cualquier $p\in P$ en tiempo $\Oh(r)$, siendo $r$ la cantidad de cuadras, y terminar apenas se tenga la información necesaria. (II) Imprimir el recorrido.

\smallskip
\textbf{Solución.} Grafo dirigido, porque nada dice que $C[i,j]=C[j,i]$. Se corre Dijkstra desde $K$ guardando $E[v]=\grpred(v)$, y el ciclo principal se corta cuando $P\subseteq S$. Por (4.14), un nodo que entra a $S$ ya tiene su camino definitivo. Para (II), \textsc{Recorrido}$(E,p)$: si $p\neq K$, llamar \textsc{Recorrido}$(E,E[p])$ y después imprimir $p$. Son $r+1$ llamadas, así que es $\Oh(r)$. Con la matriz, Dijkstra es $\Theta(n^2)$.
\end{problema}

\subsection{Dos Dijkstra: el comodín}

\begin{problema}{1P 2019, Ej.~3: comodín que anula una arista}
Grafo dirigido fuertemente conexo con longitudes positivas. Se puede poner en cero la longitud de \emph{una} arista. Hallar un camino mínimo de $s$ a $t$ usando el comodín de la mejor manera, en $\Oh(m\log n)$.

\smallskip
\textbf{Solución.} (b) Si $P$ es mínimo entre los caminos de $s$ a $t$ que usan $(v_i,v_{i+1})$, su tramo $s\leadsto v_i$ es mínimo y su tramo $v_{i+1}\leadsto t$ también. Si no, se reemplaza el tramo y se obtiene un camino más corto que sigue usando la arista. (c) Se corre Dijkstra desde $s$ en $G$, que da $d_s(u)$ y $P_u$, y Dijkstra desde $t$ en el grafo invertido $G^R$, que da $d_t(v)$, la distancia de $v$ a $t$ en $G$. Se elige la arista $(u,v)$ que minimiza $d_s(u)+0+d_t(v)$ y se concatena $P_u$, la arista y el camino de $v$ a $t$. (d) Dos Dijkstra cuestan $\Oh(m\log n)$. $G^R$ no se construye: se usan las listas de entrantes. Recorrer las aristas es $\Oh(m)$, y como el grafo es fuertemente conexo, $n=\Oh(m)$.
\end{problema}

El truco de \emph{Dijkstra desde $t$ en el grafo invertido} vale para cualquier pregunta del tipo ``el mejor camino que pasa por una arista o por un nodo''. La parte (b) es el principio de optimalidad de los caminos mínimos, que también se pidió por separado en febrero de 2013 y en el 2P 2013.\visto{1P 2019 Ej.~3: comodín, dos Dijkstra}

\subsection{Camino que visita una secuencia}

En el 2P de 2009 se pidió un camino de costo mínimo de $v_1$ a $v_k$ que visite $v_1,\dots,v_k$ en ese orden, pudiendo repetir vértices. La respuesta es concatenar caminos mínimos $v_i\leadsto v_{i+1}$, con $k-1$ corridas de Dijkstra si no hay aristas negativas o una de Floyd si solo se excluyen los ciclos negativos. La correctitud es por absurdo: si $C'$ fuera mejor, algún tramo de $C'$ entre $v_i$ y $v_{i+1}$ sería más corto que el mínimo.\visto{2P 2009 Ej.~3.1: camino por una secuencia}

\subsection{Costos en los nodos}

Si atravesar un nodo cuesta (en julio de 2023, un tiempo de procesamiento $TP(u)$ en cada router), se pasa el costo a las aristas salientes: $\ell_{(u,v)}=TP(u)+TT(u,v)$. Ese examen pedía resolverlo con Bellman–Ford (ver la parte de programación dinámica), pero con costos positivos Dijkstra también funciona.\visto{Examen jul.~2023 Ej.~2: costo en nodos y aristas}

\subsection{Caminos mínimos en un DAG}\label{gr:dag}

Si el grafo no tiene ciclos, no hace falta el heap: alcanza con recorrer los nodos en orden topológico.\visto{1P 2022 Ej.~3: caminos mínimos en un DAG}

\begin{problema}{1P 2022, Ej.~3: DAG en $\Oh(m+n)$}
$G$ dirigido acíclico con costos positivos y $s$ que alcanza a todos. Calcular los costos de los caminos mínimos desde $s$ en $\Oh(m+n)$. Sugerencia: orden topológico.

\smallskip
\textbf{Solución.} La recurrencia es $d[s]=0$ y $d[v]=\min_{(u,v)\in E}\{d[u]+c(u,v)\}$. En orden topológico $v_1=s,\dots,v_n$, cada $d[u]$ que aparece en el mínimo ya está calculado. Primero se calcula el orden con el algoritmo de los grados de entrada, en $\Oh(m+n)$. Después se recorre $v_2,\dots,v_n$ mirando las aristas entrantes de cada uno, y cada arista se mira una vez, en $\Oh(m+n)$. No hay heap ni $\log n$, y funciona aunque haya costos negativos.
\end{problema}

\subsection{Contar los caminos mínimos}

El ejercicio del libro infantil ``Elige tu propia aventura'' cayó en diciembre de 2018 y en julio de 2019: calcular el largo mínimo de una historia y cuántas historias hay de ese largo, en $\Oh(n)$. Sin costos, el nivel BFS $L(w)$ es la distancia, y la cantidad de caminos mínimos cumple
\[ N(1)=1,\qquad N(w)=\sum_{t:\,(t,w)\in E,\ L(t)=L(w)-1} N(t), \]
porque los caminos mínimos a $w$ se parten según el último nodo antes de $w$. Se calcula durante el mismo BFS.\visto{Examen dic.~2018 Ej.~1 y jul.~2019 Ej.~2: contar historias}

\subsection{Camino de cuello de botella (carga máxima)}

En el 2P de 2011 se pidió, para cada par de ciudades, el camino que permite transportar la mayor carga, sabiendo que cada carretera admite un peso máximo $W(i,j)$. La capacidad de un camino es el \emph{mínimo} de sus aristas y se quiere \emph{maximizarla}. Se pidió con programación dinámica, al estilo de Floyd:\visto{2P 2011 Ej.~2: carga máxima entre ciudades}
\[ f_0(i,j)=W(i,j),\qquad f_k(i,j)=\max\bigl\{f_{k-1}(i,j),\ \min\{f_{k-1}(i,k),f_{k-1}(k,j)\}\bigr\}. \]

\subsection{Dijkstra no da un MST, y un MST no da caminos mínimos}

Esta pregunta se repitió muchas veces: ¿existe un origen $v$ para el que el árbol de Dijkstra sea un MST? ¿Kruskal o Prim encuentran el camino de costo mínimo entre $u$ y $v$? ¿El camino mínimo está contenido en algún MST? A las tres se responde que no, con un contraejemplo.\visto{Examen dic.~2009 Ej.~2, dic.~2010 Ej.~1.2, 2P 2013 Ej.~2(d)}

\begin{center}
\begin{tikzpicture}[scale=1.05]
  \node[grv] (u) at (0,1.6) {$u$}; \node[grv] (v) at (2.2,1.6) {$v$};
  \node[grv] (w) at (0,0) {$w$};   \node[grv] (z) at (2.2,0) {$z$};
  \draw[grp] (u)--node[grl,above]{$2$}(v);
  \draw[gra] (w)--node[grl,below]{$2$}(z);
  \draw[grp] (u)--node[grl,left]{$3$}(w);
  \draw[gra] (v)--node[grl,right]{$3$}(z);
  \draw[grp] (u)--node[grl,pos=0.3]{$2$}(z);
  \draw[gra] (w)--node[grl,pos=0.3]{$2$}(v);
  \node[grt] at (1.1,-0.75) {Dijkstra desde $u$: costo $7$};
\end{tikzpicture}\hspace{1.2cm}%
\begin{tikzpicture}[scale=1.05]
  \node[grv] (u) at (0,1.6) {$u$}; \node[grv] (v) at (2.2,1.6) {$v$};
  \node[grv] (w) at (0,0) {$w$};   \node[grv] (z) at (2.2,0) {$z$};
  \draw[grp] (u)--node[grl,above]{$2$}(v);
  \draw[grp] (w)--node[grl,below]{$2$}(z);
  \draw[gra] (u)--node[grl,left]{$3$}(w);
  \draw[gra] (v)--node[grl,right]{$3$}(z);
  \draw[grp] (u)--node[grl,pos=0.3]{$2$}(z);
  \draw[gra] (w)--node[grl,pos=0.3]{$2$}(v);
  \node[grt] at (1.1,-0.75) {un MST: costo $6$};
\end{tikzpicture}
\end{center}

Por simetría, desde cualquier origen Dijkstra devuelve las tres aristas que inciden en él, de costos $2,2,3$, y llega a $7$. Un MST cuesta $6$. Para el otro sentido sirve el triángulo $c(a,b)=c(a,c)=4$, $c(b,c)=5$: el único MST es $\{ab,ac\}$, en el que el camino de $b$ a $c$ cuesta $8$, mientras que la arista directa cuesta $5$.\visto{Examen dic.~2011 Ej.~1.1: V/F camino mínimo en MST}

\begin{trampa}[caminos mínimos]
\begin{itemize}
  \item Usar BFS con costos distintos. BFS minimiza la \emph{cantidad} de aristas.
  \item Modelar con un grafo no dirigido cuando la matriz no es simétrica (dic.~2016) o cuando las aristas tienen sentido.
  \item En los ejercicios de orden, no justificar $m=\Oh(n)$ (grilla) o $n=\Oh(m)$ (grafo conexo). Sin eso, $\Oh((n+m)\log n)$ no es la cota pedida.
\end{itemize}
\end{trampa}

% =====================================================================
\section{Árboles de cubrimiento mínimo}
% =====================================================================

Sea $G=(V,E)$ conexo, no dirigido, con costos $c_e$. Un \emph{árbol de cubrimiento} es un subgrafo $(V,T)$, con $T\subseteq E$, que es conexo y acíclico. Equivalentemente, es conexo con $n-1$ aristas. Es \emph{de costo mínimo} (MST) si $\sum_{e\in T}c_e$ es mínima entre todos los árboles de cubrimiento. Estas dos definiciones se pidieron en 2P 2014 y en julio de 2013.\visto{2P 2014 Ej.~2.1 y Examen jul.~2013 Ej.~4: definir MST}

\begin{resultado}[(4.16)]
Si los costos son positivos, todo conjunto $T$ de costo mínimo tal que $(V,T)$ es conexo es un árbol.
\end{resultado}
Si $(V,T)$ tuviera un ciclo, se podría sacar una arista del ciclo sin desconectar y bajaría el costo. Por eso, cuando un enunciado pide ``conectar todo gastando lo menos posible'', la respuesta es un MST aunque no diga ``árbol''. Así se resolvió el examen de diciembre de 2020 (repavimentar caminos).

\subsection{Propiedad de corte}

\begin{resultado}[(4.17) Propiedad de corte]
Supongamos que los costos son todos distintos. Sean $S$ un subconjunto de nodos con $\emptyset\neq S\neq V$ y $e=(v,w)$ la arista de menor costo con $v\in S$ y $w\in V\setminus S$. Entonces \emph{todo} MST contiene a $e$.
\end{resultado}

\begin{proof}[Demostración (intercambio)]
Sea $T$ un árbol de cubrimiento que no contiene a $e$. Vamos a encontrar otro más barato. Como $T$ es conexo, tiene un camino $P$ de $v$ a $w$. Este camino empieza en $S$ y termina fuera, así que tiene una arista $e'=(v',w')$ con $v'\in S$ y $w'\notin S$. Sea $T'=(T\setminus\{e'\})\cup\{e\}$.
\begin{itemize}
  \item \emph{$T'$ es conexo.} Todo camino de $T$ que usaba $e'$ se puede desviar: de $v'$ a $v$ por $P$, después $e$, después de $w$ a $w'$ por $P$.
  \item \emph{$T'$ es acíclico.} Es conexo y tiene $n-1$ aristas.
  \item \emph{$T'$ es más barato.} $c(T')=c(T)-c_{e'}+c_e<c(T)$, porque $e'$ también cruza el corte y $e$ es la más barata que lo cruza.
\end{itemize}
Entonces ningún árbol sin $e$ es de costo mínimo.
\end{proof}

La arista $e'$ hay que tomarla \emph{en el camino $P$}. Si se elige cualquier arista de $T$ que cruce el corte, al cambiarla por $e$ puede quedar algo disconexo o con un ciclo. Ese es el error más común en esta demostración.

En los exámenes viejos se le llamaba \emph{propiedad MST} a una versión de existencia, que vale también con empates: si $(u,v)$ es una arista de costo mínimo entre $U$ y $V\setminus U$, \emph{existe} un MST que la contiene. La demostración oficial es la misma. Se toma un MST $T$ sin $(u,v)$; al agregarla se forma un ciclo, que tiene otra arista $(u',v')$ entre $U$ y $V\setminus U$; y $T'=T-(u',v')+(u,v)$ cumple $c(T')\le c(T)$, así que es un MST que contiene a $(u,v)$.\visto{Examen dic.~2012, dic.~2013, jul.~2015: propiedad MST}

\begin{problema}{Examen feb.~2017, P.~2: un enunciado incorrecto}
Se propone: ``si $a$ cruza el corte y ninguna arista que lo cruza tiene costo \emph{menor estricto}, entonces $a$ está en todo MST''. Dar un contraejemplo y corregirlo.

\smallskip
\textbf{Solución.} Triángulo con $c(x,y)=c(x,z)=2$ y $c(y,z)=1$. Con $U=\{x\}$, la arista $(x,y)$ cumple la hipótesis, pero $\{(x,z),(y,z)\}$ es un MST que no la contiene. Hay dos versiones correctas: ``ninguna otra tiene costo \emph{menor o igual}, entonces está en todo MST'', o ``ninguna tiene costo menor, entonces \emph{existe} un MST que la contiene''.
\end{problema}

\subsection{Propiedad de ciclo}

\begin{resultado}[(4.20) Propiedad de ciclo]
Supongamos que los costos son todos distintos. Sean $C$ un ciclo de $G$ y $e=(v,w)$ la arista de mayor costo de $C$. Entonces $e$ no pertenece a \emph{ningún} MST.
\end{resultado}

Se pidió demostrarla en el 1P de 2017 y enunciarla en el de 2025. La demostración oficial es por absurdo y con intercambio, usando el corte que deja $T$ al sacar $e$.\visto{1P 2017 Ej.~2(a) y 1P 2025 Ej.~1(c): ciclo}

\begin{proof}
Sea $T$ un árbol de cubrimiento que contiene a $e$. Al sacar $e$, $T$ queda partido en dos componentes: $S$, con $v$, y $V\setminus S$, con $w$. El resto del ciclo, $C\setminus\{e\}$, es un camino de $v$ a $w$, así que tiene una arista $e'$ que cruza de $S$ a $V\setminus S$. Entonces $T'=(T\setminus\{e\})\cup\{e'\}$ es conexo (se reconectan las dos partes) y tiene $n-1$ aristas, o sea que es un árbol. Además $c(T')=c(T)-c_e+c_{e'}<c(T)$, porque $e$ es la más cara de $C$. Por lo tanto $T$ no es mínimo.
\end{proof}

\begin{receta}[``¿la arista $e$ está en algún MST, en todos o en ninguno?'']
\begin{itemize}
  \item \emph{Está en todos.} Exhibí un corte concreto del que $e$ es la más barata y citá (4.17). En el 1P 2025 el corte era $\{4\}$ contra $\{1,2,3\}$ para la arista $(3,4)$. Aislar uno de los extremos suele funcionar.
  \item \emph{No está en ninguno.} Exhibí un ciclo concreto del que $e$ es la más cara y citá (4.20).
  \item La pauta del 1P 2025 dice que \emph{no alcanza} con nombrar la propiedad sin mostrar su aplicación en el ejemplo, que tampoco alcanza con decir que Kruskal la incluye, y que mostrar un árbol de costo $13$ no prueba que sea mínimo si no se comparan todos.
\end{itemize}
\end{receta}

El examen de diciembre de 2013 pedía usar la propiedad MST para probar que $(2,3)$ está en un MST. La solución elige $U=V\setminus\{3\}$, observa que solo dos aristas cruzan y que $(2,3)$ es la más barata de ellas.\visto{Examen dic.~2013 Ej.~1(b); 1P 2025 Ej.~1(d)}

\subsection{Kruskal, Prim y borrado inverso}

\begin{algorithmic}[1]
\Function{Kruskal}{$G,c$}
  \State ordenar $E$ por costo creciente; $T\gets\emptyset$
  \ForAll{$e=(u,v)$ en ese orden}
    \If{$u$ y $v$ están en componentes distintas de $(V,T)$} \State $T\gets T\cup\{e\}$ \EndIf
  \EndFor
  \State \Return $T$
\EndFunction
\end{algorithmic}

\begin{algorithmic}[1]
\Function{Prim}{$G,c,r$}
  \State $U\gets\{r\}$; $T\gets\emptyset$
  \While{$U\neq V$}
    \State sea $(u,v)$ la arista de menor costo con $u\in U$, $v\in V\setminus U$
    \State $T\gets T\cup\{(u,v)\}$; $U\gets U\cup\{v\}$
  \EndWhile
  \State \Return $T$
\EndFunction
\end{algorithmic}

\begin{algorithmic}[1]
\Function{BorradoInverso}{$G,c$}
  \State $T\gets E$
  \ForAll{$e$ en orden de costo \emph{decreciente}}
    \If{$(V,T\setminus\{e\})$ es conexo} \State $T\gets T\setminus\{e\}$ \EndIf
  \EndFor
  \State \Return $T$
\EndFunction
\end{algorithmic}

\begin{resultado}[(4.18), (4.19) y (4.21)]
Kruskal, Prim y el borrado inverso (\emph{reverse-delete}) devuelven un MST de $G$.
\end{resultado}

En julio de 2015 se preguntó cómo se aplica la propiedad MST en estos algoritmos, y en 2P 2014 se pidió justificar por qué Prim puede usar la técnica greedy. La respuesta es la misma en los dos casos: cada decisión está respaldada por la propiedad de corte.\visto{Examen jul.~2015 Ej.~3.2 y 2P 2014 Ej.~2.3(c)}

\begin{proof}[Prim (4.19)]
En cada iteración, $U$ es un conjunto propio no vacío y la arista elegida es la más barata entre $U$ y $V\setminus U$. Por (4.17) está en todo MST, así que $T$ queda contenido en el MST. Termina con $n-1$ aristas que forman un árbol, porque cada una agrega un nodo nuevo. Ese árbol es el MST.
\end{proof}

\begin{proof}[Kruskal (4.18)]
Sea $e=(v,w)$ una arista que Kruskal agrega, y sea $S$ el conjunto de nodos conectados a $v$ en $(V,T)$ justo antes. Entonces $w\notin S$, porque si no $e$ cerraría un ciclo. Ninguna arista entre $S$ y $V\setminus S$ fue examinada antes: si lo hubiera sido, no cerraba ciclo y se habría agregado, con lo que $S$ sería más grande. Por lo tanto $e$ es la más barata del corte $(S,V\setminus S)$ y está en todo MST por (4.17). La salida no tiene ciclos por construcción, y es conexa porque si quedaran dos componentes, alguna arista de $G$ las uniría y se habría agregado al examinarla.
\end{proof}

\begin{proof}[Borrado inverso (4.21)]
Cuando se borra $e$, el grafo sigue conexo, así que $e$ está en un ciclo $C$ del grafo actual. Las aristas que se examinaron antes y no se borraron eran puentes, es decir que no están en ningún ciclo, y por eso no están en $C$. Las demás aristas de $C$ todavía no se examinaron y son más baratas que $e$. Entonces $e$ es la más cara de $C$ y no está en ningún MST por (4.20). La salida es conexa, porque nunca se borra un puente, y acíclica, porque la más cara de un ciclo que sobreviviera se habría borrado. Como contiene al MST y es un árbol, es el MST.
\end{proof}

\begin{problema}{Examen dic.~2017, Probl.~A: el algoritmo \textsc{HallarMST}}
Se ordenan las aristas en forma decreciente. Para cada una, si $G-\{a\}$ sigue siendo conexo (con un BFS), se borra. (I) Peor caso, (II) orden, (III) ¿da un MST?, (IV) ¿es greedy?

\smallskip
\textbf{Solución.} Es el borrado inverso (4.21). (I) El peor caso es el grafo completo, que maximiza la cantidad de aristas. (II) Hay $m$ iteraciones, cada una con un BFS de $\Oh(m+n)$, así que el orden es $\Oh(m(m+n))$. (III) Sí, por (4.21). (IV) Sí, porque toma en cada paso la decisión localmente óptima.
\end{problema}

En el examen de febrero de 2015 se pedía elegir entre Prim y Kruskal con una condición: cada tendido nuevo tiene que mantener \emph{una única componente conexa}. Prim cumple esa condición porque crece un solo árbol. Kruskal no, porque construye un bosque.\visto{Examen feb.~2015 Ej.~3: una sola componente, Prim}

\subsection{Prim y Kruskal a mano}

Grafo del 2P de 2014: $c(1,2)=16$, $c(1,5)=19$, $c(1,6)=21$, $c(2,3)=3$, $c(2,4)=6$, $c(2,6)=11$, $c(3,4)=10$, $c(4,5)=18$, $c(4,6)=14$, $c(5,6)=33$. Se pedía Prim paso a paso, y la solución oficial arranca en el vértice $1$.\visto{2P 2014 Ej.~2.4: Prim paso a paso}

\begin{center}
\begin{tikzpicture}[scale=1.15]
  \node[grv] (3) at (-2.4,0) {$3$};
  \node[grv] (2) at (-1.1,1.1) {$2$};
  \node[grv] (4) at (-1.1,-1.1) {$4$};
  \node[grv] (1) at (0.6,2.0) {$1$};
  \node[grv] (6) at (0.8,0) {$6$};
  \node[grv] (5) at (2.6,0) {$5$};
  \draw[grp] (3)--node[grl,above left]{$3$}(2);
  \draw[gra] (3)--node[grl,below left]{$10$}(4);
  \draw[grp] (2)--node[grl,left]{$6$}(4);
  \draw[grp] (2)--node[grl,above,pos=0.55]{$11$}(6);
  \draw[gra] (4)--node[grl,below,pos=0.55]{$14$}(6);
  \draw[grp] (2)--node[grl,above left]{$16$}(1);
  \draw[gra] (1)--node[grl,right,pos=0.4]{$21$}(6);
  \draw[gra] (1)--node[grl,above right]{$19$}(5);
  \draw[gra] (6)--node[grl,above]{$33$}(5);
  \draw[grp] (4) to[bend right=22] node[grl,below right]{$18$}(5);
\end{tikzpicture}
\end{center}

\begin{center}
\small
\begin{tabularx}{\textwidth}{@{}c>{\raggedright\arraybackslash}p{2.4cm}>{\raggedright\arraybackslash}X>{\raggedright\arraybackslash}p{1.9cm}@{}}
\toprule
Paso & $U$ & Aristas entre $U$ y $V\setminus U$ & Elegida \\
\midrule
1 & $\{1\}$ & $16,19,21$ & $(1,2)$, $16$ \\
2 & $\{1,2\}$ & $19,21,3,6,11$ & $(2,3)$, $3$ \\
3 & $\{1,2,3\}$ & $19,21,6,11,10$ & $(2,4)$, $6$ \\
4 & $\{1,2,3,4\}$ & $19,21,11,14,18$ & $(2,6)$, $11$ \\
5 & $\{1,2,3,4,6\}$ & $19,18,33$ & $(4,5)$, $18$ \\
\bottomrule
\end{tabularx}
\end{center}

Costo total $16+3+6+11+18=54$. Kruskal sobre el mismo grafo examina $3$ (agrega $(2,3)$), $6$ (agrega $(2,4)$), $10$ (descarta $(3,4)$ porque $3$ y $4$ ya están en $\{2,3,4\}$), $11$ (agrega $(2,6)$), $14$ (descarta $(4,6)$), $16$ (agrega $(1,2)$) y $18$ (agrega $(4,5)$). Con eso llega a $n-1=5$ aristas y termina. Da el mismo árbol, como tiene que ser: los costos son distintos y el MST es único. Los grafos de diciembre de 2011, julio de 2012 y julio de 2013 se resolvían de la misma manera.\visto{Examen dic.~2011 Ej.~2, jul.~2012 Ej.~3, jul.~2013 Ej.~4}

\begin{problema}{Examen dic.~2025, Ej.~1(c): Prim desde $s$}
Aristas $c(s,x)=2$, $c(s,t)=4$, $c(x,y)=3$, $c(x,t)=5$, $c(y,t)=1$. Ejecutar Prim desde $s$ indicando la arista de cada iteración.

\smallskip
\textbf{Solución.} Paso 1: entre $(s,x)=2$ y $(s,t)=4$, elige $(s,x)$. Paso 2: entre $(s,t)=4$, $(x,y)=3$ y $(x,t)=5$, elige $(x,y)$. Paso 3: entre $(s,t)=4$, $(x,t)=5$ e $(y,t)=1$, elige $(y,t)$. Ya están todos los vértices. El MST tiene costo $6$. La pauta muestra las candidatas de cada paso.
\end{problema}

\begin{receta}[ejecutar Prim o Kruskal]
\begin{itemize}
  \item \emph{Prim:} decí desde qué vértice arrancás. En cada paso escribí $U$, las aristas candidatas entre $U$ y $V\setminus U$ con su costo, y la elegida.
  \item \emph{Kruskal:} listá las aristas ordenadas y marcá cada una como agregada o descartada, diciendo qué ciclo cerraría o qué componentes une. Cortá al llegar a $n-1$ aristas.
  \item Terminá con el árbol dibujado y su costo total.
\end{itemize}
\end{receta}

\subsection{Implementaciones}

\begin{resultado}[(4.22) y (4.25)]
Con un heap, Prim se implementa en $\Oh(m\log n)$. Con union-find, Kruskal también se implementa en $\Oh(m\log n)$.
\end{resultado}

\emph{Prim} es igual a Dijkstra, cambiando la clave. Para cada $v\notin U$ se guarda $a(v)=\min_{u\in U}c(u,v)$, que es el costo de \emph{una} arista, en vez de $d(u)+\ell$. Se hacen $n$ extracciones del mínimo y a lo sumo $m$ cambios de clave, lo que da $\Oh(m\log n)$. \emph{Kruskal}: ordenar cuesta $\Oh(m\log m)=\Oh(m\log n)$, porque $m\le n^2$. Después vienen $2m$ \grop{Find} y $n-1$ \grop{Union}, y con union-find (unión por tamaño: \grop{Find} en $\Oh(\log n)$, \grop{Union} en $\Oh(1)$) eso es $\Oh(m\log n)$. En diciembre de 2020 había que \emph{repetir} este análisis.\visto{Examen dic.~2020 Ej.~1: MST en $\Oh(m\log n)$}

\begin{problema}{Examen dic.~2020, Ej.~1: repavimentar caminos}
$n$ centros urbanos y $m$ caminos con longitud. Elegir caminos a repavimentar de modo que se pueda ir entre cualquier par usando solo caminos repavimentados, con la menor cantidad total de km. Dar un algoritmo $\Oh(m\log n)$ y probar la cota.

\smallskip
\textbf{Solución.} Un nodo por centro y una arista por camino, con costo igual a la longitud. El grafo es conexo por hipótesis. Lo que se busca es $T$ con $(V,T)$ conexo y costo mínimo, que es la definición de MST. Construir las listas de adyacencia es $\Oh(m+n)=\Oh(m)$ porque $n\le m+1$, y Prim o Kruskal (reescrito) agrega $\Oh(m\log n)$.
\end{problema}

\subsection{Unicidad y empates}

En julio de 2018 se pidió enunciar la propiedad de corte y después responder dos preguntas: con costos distintos, ¿puede haber dos MST distintos? ¿Pueden Prim y Kruskal dar resultados distintos?\visto{Examen jul.~2018 P.~3: corte y unicidad}

\begin{resultado}[Unicidad (KT Ej.~4.8)]
Si los costos son todos distintos, el MST es único.
\end{resultado}

\begin{proof}
Sean $T$ un MST y $e=(u,w)\in T$. Al sacar $e$, $T$ queda partido en $S\ni u$ y $V\setminus S\ni w$. La arista $e$ es la más barata del corte $(S,V\setminus S)$: si otra $f$ fuera más barata, $T-e+f$ sería un árbol más barato. Por (4.17), $e$ está en todo MST. Como $e$ era cualquiera, $T\subseteq T^*$ para todo MST $T^*$, y como los dos tienen $n-1$ aristas, $T=T^*$.
\end{proof}

De la unicidad sale la otra respuesta: con costos distintos, Prim y Kruskal devuelven el mismo árbol. Con empates, en cambio, pueden diferir. En febrero de 2016 se pidió un ejemplo. En un triángulo con las tres aristas de costo $1$, Kruskal con el orden $xy,yz,xz$ da $\{xy,yz\}$, y Prim desde $z$ que desempata eligiendo $zx$ y después $xy$ da $\{zx,xy\}$.\visto{Examen feb.~2016 Ej.~3.2: MST distintos con empates}

\begin{trampa}[MST]
\begin{itemize}
  \item Usar (4.17) o (4.20) sin decir que los costos son distintos, o aplicarlas con empates. Con empates vale la versión de \emph{existencia}.
  \item En la demostración del corte, cambiar $e$ por \emph{cualquier} arista que cruce el corte. Tiene que ser una del camino de $v$ a $w$ en $T$.
  \item Confundir ``la más barata del corte'' con ``la más barata incidente a un nodo'' cuando el corte no es $\{v\}$.
\end{itemize}
\end{trampa}

% =====================================================================
\section{Modificar un MST}
% =====================================================================

Son ejercicios cortos que salen todos de las dos propiedades. Se suponen costos distintos.

\subsection{Agregar una arista}

\begin{problema}{1P 2017, Ej.~2(b) : agregar una arista}
$T$ es MST de $G$ y se agrega $e=(v,w)$ de costo $c$. Verificar eficientemente si $T$ sigue siendo MST.

\smallskip
\textbf{Solución.} Se busca el único camino $P$ de $v$ a $w$ en $T$, con un DFS o BFS sobre $T$ en $\Oh(n)$. Con $e$, ese camino forma un ciclo. Si todas las aristas de $P$ son más baratas que $c$, $e$ es la más cara del ciclo, no está en ningún MST por (4.20) y $T$ sigue siendo MST. Si alguna es más cara, la más cara del ciclo, $e_{\max}$, está en $T$ y no está en ningún MST, así que $T$ ya no es MST. El nuevo es $T'=T+e-e_{\max}$. Para probarlo, toda arista fuera de $T'$ es la más cara de algún ciclo: $e_{\max}$ lo es del ciclo $P+e$, y cada $f\notin T$ lo es de su ciclo fundamental en $T$, porque $T$ era MST. Entonces el MST nuevo está contenido en $T'$, y como los dos tienen $n-1$ aristas, son iguales.
\end{problema}

La solución oficial de 2017 escribe un DFS recursivo sobre $T$ que devuelve el par (¿sigue siendo MST?, ¿llegué a $w$?) y va comparando $c$ con cada arista del camino al volver de la recursión.\visto{1P 2017 Ej.~2(b): agregar arista al MST}

\subsection{Quitar una arista}

\begin{problema}{1P 2020, Ej.~3: quitar una arista del MST}
$T$ es MST de $G$, $(u,v)\in T$ y $G'=G-(u,v)$ es conexo. Sabiendo que $T\setminus\{(u,v)\}\subseteq T'$, el MST de $G'$: (1) $S\gets$ componente de $u$ en $T\setminus\{(u,v)\}$; (2) $(u',v')\gets$ la arista más barata de $G'$ entre $S$ y $V\setminus S$; (3) $T'\gets T\setminus\{(u,v)\}\cup\{(u',v')\}$. Probar que es correcto y que es $\Oh(m)$.

\smallskip
\textbf{Solución.} (a) Por (4.17) aplicada en $G'$, $(u',v')\in T'$. Por hipótesis, también $T\setminus\{(u,v)\}\subseteq T'$. Además $(u',v')\notin T$, porque une $S$ con $V\setminus S$ igual que $(u,v)$ y cerraría un ciclo. Entonces el conjunto del paso 3 tiene $n-1$ aristas y está contenido en $T'$, que también tiene $n-1$, así que coinciden. (b) Para $S$ se arman las listas de $T\setminus\{(u,v)\}$ y se corre un BFS en $\Oh(n)$, dejando un arreglo booleano. Para el paso 2 se recorren las $m$ aristas chequeando en $\Oh(1)$ si cruzan. El total es $\Oh(m+n)=\Oh(m)$.
\end{problema}

\subsection{¿Es la más cara de algún ciclo?}

\begin{problema}{1P 2020, Ej.~1: la arista más cara de algún ciclo}
Costos distintos. Determinar si $(u,v)$ es la arista de mayor costo de algún ciclo, en $\Oh(m+n)$.

\smallskip
\textbf{Solución.} Sea $G'$ el grafo que resulta de borrar las aristas más caras que $(u,v)$. Entonces $(u,v)$ es la más cara de un ciclo de $G$ si y solo si está en un ciclo de $G'$: ese ciclo no usa aristas borradas, y en $G'$ la más cara es $(u,v)$. A su vez, $(u,v)$ está en un ciclo de $G'$ si y solo si $u$ y $v$ quedan conectados en $G'-(u,v)$. Algoritmo: borrar las aristas de costo $\ge c(u,v)$ en $\Oh(m+n)$ y hacer un BFS desde $u$ en $\Oh(m+n)$ para ver si llega a $v$.
\end{problema}

Con costos distintos, ``$e$ es la más cara de algún ciclo'' equivale a ``$e$ no está en el MST''. Si es la más cara de un ciclo, (4.20) dice que no está. Si no está, al agregarla al MST se forma un ciclo del que es la más cara, por el argumento del ejercicio anterior. Así que este algoritmo decide en $\Oh(m+n)$ si una arista pertenece al MST, sin construirlo.\visto{1P 2020 Ej.~1: más cara de un ciclo}

\subsection{Pocos ciclos: \texorpdfstring{$m=n$}{m=n}}

En julio de 2015 (problema 2) se pedía el MST de un grafo conexo con $m=n$ \emph{sin} usar la propiedad MST. Un grafo así tiene exactamente un ciclo: se lo encuentra con un DFS que guarda los padres, se busca su arista más cara y se la saca. Cualquier arista del ciclo que se saque deja un árbol de cubrimiento, y el más barato de esos árboles es el que resulta de sacar la más cara.\visto{Examen jul.~2015 Probl.~2: MST con $m=n$}

\subsection{Transformar los costos}\label{gr:transf}

Kruskal solo compara costos. Por eso cualquier transformación estrictamente creciente de los costos ($c_e^2$ con costos positivos, $c_e+k$, $\log c_e$) deja el mismo MST. Una decreciente lo convierte en árbol de peso \emph{máximo}. Con caminos mínimos no pasa lo mismo, porque se \emph{suman} costos de caminos de distinto largo.

\begin{problema}{Examen feb.~2024, Ej.~1: red de cableado de largo máximo}
Distancias $d_{ij}$ entre $n$ proveedores. Se busca una red conexa sin ciclos que \emph{maximice} el largo total, en $\Oh(n^2\log n)$, y demostrar la correctitud.

\smallskip
\textbf{Solución.} Grafo completo con costo $c_{ij}=M-d_{ij}$, donde $M>\max d_{ij}$, y $T=\textsc{MST}(G)$. Para cualquier árbol $T'$:
\[ \textstyle\sum_{T}d_{ij}=\sum_{T}(M-c_{ij})=(n-1)M-\sum_T c_{ij}\ \ge\ (n-1)M-\sum_{T'}c_{ij}=\sum_{T'}d_{ij}. \]
Las igualdades valen porque todo árbol tiene $n-1$ aristas, y la desigualdad porque $T$ es MST. Con $m=\Theta(n^2)$, el MST cuesta $\Oh(n^2\log n)$.
\end{problema}

La pauta valora que se diga en qué paso se usa ``todo árbol tiene $n-1$ aristas''. Es justamente lo que falla en caminos mínimos.\visto{Examen feb.~2024 Ej.~1: árbol de peso máximo}

% =====================================================================
\section{Árbol de cuello de botella mínimo}
% =====================================================================

La \emph{arista cuello de botella} de un árbol de cubrimiento es la más cara. Un árbol es \emph{de cuello de botella mínimo} (ACBM) si ningún otro árbol de cubrimiento tiene una arista cuello de botella más barata.

\begin{problema}{Examen feb.~2026, Ej.~2: MST y ACBM}
Costos positivos y distintos. (a) Probar por contradicción que todo MST es ACBM (sugerencia: propiedad de ciclo). (b) Mostrar con un contraejemplo que el recíproco es falso.

\smallskip
\textbf{Solución.} (a) Sean $T$ el MST, $e$ su arista más cara, y supongamos que existe un árbol $T'$ con todas sus aristas más baratas que $c_e$. Entonces $e\notin T'$. Al agregar $e$ a $T'$ se forma un ciclo en el que todas las demás aristas son de $T'$, o sea más baratas, y $e$ es la más cara. Por (4.20), $e$ no está en ningún MST, lo que contradice $e\in T$. (b) Aristas $c(a,b)=4$, $c(b,c)=1$, $c(b,d)=2$, $c(c,d)=3$. Todo árbol contiene a $(a,b)$, porque es la única arista de $a$, así que todos tienen cuello de botella $4$ y son ACBM. Pero solo $\{ab,bc,bd\}$, de costo $7$, es MST.
\end{problema}

\begin{center}
\begin{tikzpicture}[scale=1.05]
  \node[grv] (a) at (0,0) {$a$};
  \node[grv] (b) at (1.8,0) {$b$};
  \node[grv] (c) at (3.3,0.9) {$c$};
  \node[grv] (d) at (3.3,-0.9) {$d$};
  \draw[grp] (a)--node[grl,above]{$4$}(b);
  \draw[grp] (b)--node[grl,above left]{$1$}(c);
  \draw[grp] (b)--node[grl,below left]{$2$}(d);
  \draw[gra] (c)--node[grl,right]{$3$}(d);
  \node[grt] at (1.6,-1.25) {MST de costo $7$; $\{ab,bc,cd\}$ cuesta $8$ y también es ACBM};
\end{tikzpicture}
\end{center}

La pauta marca dos errores. Uno es agregar la arista cuello de botella de $T'$ a $T$ en lugar de $e$ a $T'$: puede que no se forme ciclo, o que $e$ no esté en él. El otro es concluir que ``$e$ está en un ciclo'' sin probar que es la \emph{más cara} de ese ciclo, que es lo que garantiza que las demás aristas sean de $T'$. También vale una prueba con la propiedad de corte. Se toma el corte que deja $T$ al sacar $e$: por (4.17), $e$ es la más barata de ese corte, y $T'$ tiene que cruzarlo con alguna arista de costo $\ge c_e$.\visto{Examen feb.~2026 Ej.~2: MST implica ACBM}

% =====================================================================
\section{Mapa de lo que se preguntó}
% =====================================================================

Relevamiento de los primeros y segundos parciales 2008--2025 y de los exámenes 2009--2026.

\begin{center}
\small
\begin{tabularx}{\textwidth}{@{}>{\raggedright\arraybackslash}p{2.8cm}>{\raggedright\arraybackslash}p{2.4cm}>{\raggedright\arraybackslash}X@{}}
\toprule
Fuente & Tema & Qué se pedía \\
\midrule
2P 2009 Ej.~3.1 & Caminos mínimos & Camino mínimo que visita una secuencia en orden; condiciones sobre los costos. \\
Ex.\ dic.~2009 Ej.~2 & Dijkstra vs.\ MST & ¿Algún origen hace que Dijkstra dé un MST? No, contraejemplo. \\
Ex.\ dic.~2009 Ej.~4 & Esquema greedy & Forma general del greedy y correspondencias con Prim. \\
Ex.\ feb.~2010 Ej.~3 & Dijkstra a mano & Explicar y ejecutar desde $1$ con costos y predecesores. \\
Ex.\ dic.~2010 Ej.~1--2 & Prim, Dijkstra & Prim a mano desde $A$; Dijkstra no da MST; el camino mínimo no está en el MST. \\
2P 2011 Ej.~2 & Cuello de botella & Carga máxima entre todo par de ciudades, con PD al estilo Floyd. \\
Ex.\ jul.~2011 Ej.~3 & Dijkstra & Esquema greedy; qué resuelve Dijkstra y pseudocódigo. \\
Ex.\ dic.~2011 Ej.~1--2 & MST & V/F sobre caminos mínimos en el MST; Kruskal y Prim paso a paso. \\
2P 2012 Ej.~2(a) & Kruskal & Qué resuelve, estrategia, correspondencia con el esquema. \\
Ex.\ jul.~2012 Ej.~3 & Kruskal & Qué resuelven Prim y Kruskal; Kruskal paso a paso. \\
Ex.\ dic.~2012 Ej.~1 & Propiedad MST & Enunciar, demostrar y decir para qué sirve. \\
Ex.\ feb.~2013 Ej.~2 & Caminos mínimos & Principio de optimalidad para el camino de costo mínimo. \\
2P 2013 Ej.~2 & Caminos y MST & ¿Kruskal o Prim dan caminos mínimos?; optimalidad; Floyd; ¿Dijkstra da MST? \\
Ex.\ jul.~2013 Ej.~1, 4 & MST & V/F ``Prim y Kruskal son greedy''; definir MST y ejecutar Kruskal. \\
Ex.\ dic.~2013 Ej.~1 & Propiedad MST & Enunciar y demostrar; usarla para probar que $(2,3)$ está en un MST. \\
2P 2014 Ej.~2 & MST, Dijkstra & Definiciones, propiedad MST, Prim (código, técnica, a mano), Dijkstra. \\
Ex.\ dic.~2014 Ej.~3 & Dijkstra a mano & Ejecutar desde el vértice $2$ justificando cada paso. \\
Ex.\ feb.~2015 Ej.~3 & Prim & Tendido eléctrico con una sola componente conexa en cada paso. \\
Ex.\ jul.~2015 Ej.~3, Probl.~2 & Propiedad MST & Demostrar y aplicarla en Prim y Kruskal; MST con $m=n$. \\
Ex.\ dic.~2015 Ej.~3 & BFS vs.\ Dijkstra & BFS sin costos; con costos no sirve, se usa Dijkstra. \\
Ex.\ feb.~2016 Ej.~3 & Empates & Grafo en el que Prim y Kruskal dan MST distintos. \\
Ex.\ jul.~2016 Ej.~4 & Técnicas & Qué resuelven Dijkstra y Prim y con qué técnica. \\
\bottomrule
\end{tabularx}
\end{center}

\begin{center}
\small
\begin{tabularx}{\textwidth}{@{}>{\raggedright\arraybackslash}p{2.8cm}>{\raggedright\arraybackslash}p{2.4cm}>{\raggedright\arraybackslash}X@{}}
\toprule
Fuente & Tema & Qué se pedía \\
\midrule
Ex.\ dic.~2016 Probl.~A & Dijkstra & Recorridos a paradas: predecesores, corte temprano, camino en $\Oh(r)$. \\
Ex.\ feb.~2017 P.~2 & Propiedad MST & Enunciado con empates incorrecto: contraejemplo y versión correcta. \\
1P 2017 Ej.~2 & Ciclo, modificar & Demostrar la propiedad de ciclo; ¿$T$ sigue siendo MST al agregar una arista? \\
Ex.\ dic.~2017 Probl.~A & Borrado inverso & Peor caso, orden $\Oh(m(m+n))$, ¿da MST?, ¿es greedy? \\
Ex.\ feb.~2018 P.~4 & Técnicas & Qué resuelven Dijkstra y Kruskal y con qué técnica. \\
Ex.\ jul.~2018 P.~3 & Corte, unicidad & Enunciar (4.17); unicidad con costos distintos; Prim vs.\ Kruskal. \\
Ex.\ dic.~2018 Ej.~1 & Contar caminos & Largo mínimo de las historias y cuántas hay: BFS con conteo. \\
1P 2019 Ej.~3 & Dijkstra & Comodín que anula una arista: dos Dijkstra, $\Oh(m\log n)$. \\
Ex.\ jul.~2019 Ej.~1 & Correctitud & Paso central de (4.14); falla con longitudes negativas. \\
Ex.\ jul.~2019 Ej.~2 & Contar caminos & El mismo ejercicio de diciembre de 2018. \\
1P 2020 Ej.~1 & Ciclo & ¿$(u,v)$ es la arista más cara de algún ciclo?, en $\Oh(m+n)$. \\
1P 2020 Ej.~3 & Modificar & Quitar una arista del MST y reemplazarla: correctitud y $\Oh(m)$. \\
Ex.\ jul.~2020 Ej.~1 & BFS & Robot más cercano en una grilla, en $\Oh(M^2)$. \\
Ex.\ dic.~2020 Ej.~1 & MST & Repavimentar caminos: modelar como MST y probar $\Oh(m\log n)$. \\
1P 2021 Ej.~2 & Dijkstra & Robot en grilla con diagonales, en $\Oh(n\log n)$. \\
1P 2022 Ej.~3 & DAG & Caminos mínimos en un DAG en $\Oh(m+n)$ con orden topológico. \\
Ex.\ jul.~2023 Ej.~2 & Costos en nodos & Procesamiento más transmisión (con Bellman–Ford; ver PD). \\
Ex.\ feb.~2024 Ej.~1 & Árbol máximo & Árbol de peso máximo vía MST con costos $M-d_{ij}$. \\
1P 2024 Ej.~2 & Dijkstra & Aspiradora en grilla, igual que 1P 2021. \\
1P 2025 Ej.~1(c--e) & Ciclo, corte, Dijkstra & Enunciar la propiedad de ciclo; arista en todo MST; Dijkstra a mano. \\
1P 2025 Ej.~2 & BFS & Viñedo: BFS con varias fuentes; Dijkstra no da el orden. \\
Ex.\ dic.~2025 Ej.~1(c) & Prim a mano & Ejecutar Prim desde $s$ indicando la arista de cada paso. \\
Ex.\ feb.~2026 Ej.~1(c), 2 & Dijkstra, ACBM & Propiedad del próximo vértice de $S$; MST implica ACBM, no al revés. \\
\bottomrule
\end{tabularx}
\end{center}

\textbf{Lectura rápida.} Lo que más se repite es la propiedad de corte (enunciarla, demostrarla o aplicarla a un ejemplo), con nueve apariciones, y las ejecuciones a mano de Dijkstra, Prim o Kruskal, también con nueve. Le siguen la pregunta ``¿Dijkstra da un MST? ¿un MST da caminos mínimos?'' y los problemas de modelado que se resuelven con Dijkstra o con un MST reescribiendo el algoritmo y su análisis de $\Oh(m\log n)$. Desde 2017 se agregan las modificaciones de un MST (agregar o quitar una arista, la más cara de un ciclo) y las variantes de cuello de botella.

% =====================================================================
%  Parte 1 — Dividir y conquistar
%  Fragmento para \input desde el documento maestro (usa estilo2p.tex)
% =====================================================================

% Macros locales de esta parte.
% \dycmas{...}: agrega otra referencia dentro de un mismo \visto, para
% que varias referencias ancladas al mismo párrafo no se superpongan.
\newcommand{\dycmas}[1]{\par\smallskip{\color{acento}\rule[0.15em]{0.9em}{1.4pt}}\hspace{0.35em}#1}

\parte{dyc}{Dividir y conquistar}{T(n)\le 2T(n/2)+cn \;\Rightarrow\; T(n)=\Oh(n\log n)}{Recurrencias y cómo resolverlas, los algoritmos del capítulo~5 del libro (mergesort, inversiones, par de puntos más cercano, multiplicación de enteros), quicksort, selección, búsqueda binaria y los problemas de examen que se resuelven partiendo a la mitad. Desde 2022 cae casi en todos los semestres un ejercicio o una parte de ejercicio: diseñar el algoritmo, plantear la recurrencia y justificar el orden.}

% ---------------------------------------------------------------------
\section{El esquema y qué se pide en el parcial}

Un algoritmo de dividir y conquistar (también lo vas a ver escrito como divide y vencerás) tiene tres partes: \emph{dividir} la instancia en subinstancias más chicas del mismo problema, \emph{resolverlas} recursivamente y \emph{combinar} sus soluciones. Además hay un \emph{caso base} que se resuelve directamente. El análisis del tiempo de ejecución sale de esas mismas cuatro piezas.\visto{Examen feb. 2018 P.~3: qué tiempos entran en la recurrencia}

\begin{resultado}[Los tiempos que forman la recurrencia]
Para plantear la relación de recurrencia hay que considerar el tiempo de \textbf{dividir} la instancia, el de \textbf{resolver las subinstancias} (un término $T(\cdot)$ por llamada recursiva), el de \textbf{combinar} las soluciones parciales y el de resolver el \textbf{caso base} (la condición inicial). Dividir y combinar se suman en un único término independiente, por ejemplo $cn$.
\end{resultado}

En los parciales y exámenes recientes el formato es muy estable. Se pide (a)~escribir un algoritmo de dividir y conquistar con cierto orden, (b)~escribir la relación de recurrencia ``explicando el origen de cada término'' y (c)~a veces, demostrar la correctitud o resolver la recurrencia por inducción. Muchas veces la letra da una sugerencia (``siga una estrategia similar a MergeSort'', ``similar a la que emplea QuickSort'') que indica exactamente qué división usar.\visto{Examen feb. 2025 Ej.~2: criterios de corrección explícitos}

\begin{receta}[estructura de una respuesta completa]
\begin{enumerate}
  \item \textbf{Algoritmo} con nombre, parámetros y precondición (por ejemplo, ``$n$ es potencia de~2'', ``$O$ está ordenado''). Líneas numeradas, para poder referirse a ellas después.
  \item Indicá la \textbf{invocación inicial} (por ejemplo, \textsc{Faltante}$(O,F,n)$ o \textsc{Inv}$(A,1,n)$).
  \item Verificá en voz alta lo que los correctores chequean: hay caso base, \emph{en cada llamada recursiva los parámetros se acercan al caso base}, en todos los casos se devuelve algo del tipo correcto.
  \item \textbf{Recurrencia con condición inicial}: $T(1)\le c_b$ (o $T(2)\le c$) y el caso general. Explicá cada término: cuántas llamadas recursivas hay y por qué tienen tamaño $n/2$, y de dónde sale el término independiente (qué pasos son $\Oh(n)$ o $\Oh(1)$).
  \item \textbf{Cerrá} citando el resultado del libro que resuelve la recurrencia: ``por (5.2) es $\Oh(n\log n)$'', ``por (5.5) es $\Oh(n)$'', o demostrándolo por inducción si se pide.
  \item Si se pide correctitud: \textbf{inducción en $n$} (o en $i$ con $n=2^i$). Paso base con el caso base del algoritmo, paso inductivo usando la hipótesis sobre las llamadas recursivas y justificando la combinación.
\end{enumerate}
\end{receta}

% ---------------------------------------------------------------------
\section{Recurrencias}

\subsection{Los resultados del libro: (5.1) a (5.6)}

El libro trabaja con desigualdades (``$\le$'') y con $n$ potencia de~2, y así lo hacen las soluciones oficiales. Usar ``$\le$'' alcanza para cotas $\Oh$; para $\Theta$ hace falta la cota inferior también.

\begin{resultado}[Capítulo 5, secciones 5.1 y 5.2]
\textbf{(5.1)} Para alguna constante $c$: $T(n)\le 2T(n/2)+cn$ si $n>2$, y $T(2)\le c$.\\
\textbf{(5.2)} Toda función $T$ que satisface (5.1) es $\Oh(n\log n)$ para $n>1$.

\smallskip
\textbf{(5.3)} Para alguna constante $c$: $T(n)\le qT(n/2)+cn$ si $n>2$, y $T(2)\le c$.\\
\textbf{(5.4)} Si $q>2$, toda función que satisface (5.3) es $\Oh(n^{\log_2 q})$.\\
\textbf{(5.5)} Si $q=1$, toda función que satisface (5.3) es $\Oh(n)$. (Con $q=2$ es (5.1).)

\smallskip
\textbf{(5.6)} Para alguna constante $c$: $T(n)\le 2T(n/2)+cn^2$ si $n>2$, y $T(2)\le c$; toda función que la satisface es $\Oh(n^2)$.
\end{resultado}

La pregunta directa ``escriba la recurrencia para $q$ subproblemas de tamaño $n/2$ y diga qué da para $q=1$ y $q=2$'' apareció tal cual, sin pedir demostración. La parte (a) del 2P 2025 pedía la recurrencia de Karatsuba y la cota ajustada, también sin demostrar.\visto{Examen jul. 2026 Ej.~1(e): $qT(n/2)+cn$, casos $q=1,2$\dycmas{2P 2025 Ej.~1(a): recurrencia de Karatsuba}}

\subsection{Cómo se resuelven: desenrollar, árbol, sustitución}

\paragraph{Desenrollar la recursión (árbol de recursión).} Se mira el nivel $j$ del árbol de llamadas: cuántos subproblemas hay, de qué tamaño, cuánto cuesta cada uno sin contar sus hijos, y se suma sobre los $\log_2 n$ niveles. Para (5.1): en el nivel $j$ hay $2^j$ subproblemas de tamaño $n/2^j$, cada uno con costo $c\,n/2^j$, así que cada nivel cuesta $cn$. Hay $\log_2 n$ niveles, total $cn\log_2 n$.

\begin{center}
\begin{tikzpicture}[
  every node/.style={font=\small},
  nodo/.style={draw=acento,rounded corners=2pt,inner sep=2pt,minimum width=1.0cm},
  level distance=0.85cm,
  level 1/.style={sibling distance=4.4cm},
  level 2/.style={sibling distance=2.2cm},
  edge from parent/.style={draw=gris}]
\node[nodo] {$cn$}
  child {node[nodo] {$cn/2$}
    child {node[nodo] {$cn/4$}}
    child {node[nodo] {$cn/4$}}}
  child {node[nodo] {$cn/2$}
    child {node[nodo] {$cn/4$}}
    child {node[nodo] {$cn/4$}}};
\node[gris,anchor=west] at (4.6,0) {nivel 0: $cn$};
\node[gris,anchor=west] at (4.6,-0.85) {nivel 1: $2\cdot cn/2=cn$};
\node[gris,anchor=west] at (4.6,-1.7) {nivel 2: $4\cdot cn/4=cn$};
\node[gris] at (0,-2.3) {$\vdots$};
\node[gris,anchor=west] at (4.6,-2.3) {$\log_2 n$ niveles $\Rightarrow cn\log_2 n$};
\end{tikzpicture}
\end{center}

Con $q$ subproblemas el nivel $j$ cuesta $q^j\cdot c\,n/2^j=(q/2)^j\,cn$, una serie geométrica:
\begin{itemize}
  \item $q=1$: $\sum_j (1/2)^j cn\le 2cn$, o sea $\Oh(n)$. Domina la raíz.
  \item $q=2$: todos los niveles iguales, $\Oh(n\log n)$.
  \item $q>2$: con $r=q/2>1$, $\sum_{j<\log_2 n} r^j cn\le \frac{cn}{r-1}r^{\log_2 n}$, y como $r^{\log_2 n}=n^{\log_2 r}=n^{\log_2 q-1}$, queda $\Oh(n^{\log_2 q})$. Dominan las hojas.
\end{itemize}
Para (5.6), el nivel $j$ cuesta $2^j c(n/2^j)^2=cn^2/2^j$, que suma a lo sumo $2cn^2$.

\paragraph{Sustitución (conjetura e inducción).} Se conjetura una cota y se prueba por inducción. Para (5.1), con $T(n)\le cn\log_2 n$: $T(n)\le 2c\frac n2\log_2\frac n2+cn=cn(\log_2 n-1)+cn=cn\log_2 n$. Es la técnica que se pide cuando la letra dice ``demuestre que $T(n)$ es $\Oh(\cdot)$''. Si la conjetura directa no cierra, se resta un término de orden menor (ver Karatsuba, sección~\ref{sec:dyc-karatsuba}).\visto{1P 2018 Ej.~2(b): $2T(n/2)+c$ es $\Oh(n)$\dycmas{Examen feb. 2022 Ej.~1(b): inducción $f(n)\le cn$}}

\begin{problema}{Combinación constante: $T(n)\le 2T(n/2)+c$ es $\Oh(n)$}
Probar que si $T(2)\le c$ y $T(n)\le 2T(n/2)+c$ para $n>2$ (potencia de~2), entonces $T(n)$ es $\Oh(n)$.

\textbf{Solución.} La conjetura $T(n)\le cn$ no cierra ($2cn/2+c=cn+c$). Se prueba la más fuerte $T(n)\le cn-c$. Base: $T(2)\le c=2c-c$. Paso: $T(n)\le 2\bigl(c\frac n2-c\bigr)+c=cn-c$. Esta recurrencia aparece en todos los algoritmos que hacen dos llamadas y combinan en tiempo constante (mín--máx, búsqueda en arreglo sin orden, conjunto independiente en un árbol, la leche ``estropeada'' en la versión mejorada).
\end{problema}

\subsection{Tabla de recurrencias frecuentes}

\begin{center}\small
\begin{tabularx}{\textwidth}{@{}>{$}l<{$}>{$}l<{$}X@{}}
\toprule
\text{Recurrencia} & \text{Solución} & Dónde aparece \\
\midrule
T(n/2)+c & \Theta(\log n) & búsqueda binaria, faltante en progresión \\
T(n/2)+cn & \Theta(n) & (5.5): selección, faltante con \textsc{separar} \\
2T(n/2)+c & \Theta(n) & mín--máx, búsqueda sin orden, árbol \\
2T(n/2)+cn & \Theta(n\log n) & (5.2): mergesort, inversiones, par más cercano \\
2T(n/2)+cn\log n & \Theta(n\log^2 n) & par más cercano ordenando la franja en cada llamada \\
3T(n/2)+cn & \Theta(n^{\log_2 3}) & (5.4): Karatsuba ($\approx n^{1{,}59}$) \\
4T(n/2)+cn & \Theta(n^2) & (5.4): multiplicación ingenua por mitades \\
4T(n/2)+c & \Theta(n^2) & mín--máx de una matriz $n\times n$ (lineal en $n^2$ datos) \\
7T(n/2)+cn^2 & \Theta(n^{\log_2 7}) & Strassen ($\approx n^{2{,}81}$) \\
2T(n/2)+cn^2 & \Theta(n^2) & (5.6) \\
3T(n/3)+cn^2 & \Theta(n^2) & domina la raíz \\
T(n-1)+T(n-2)+1 & \Theta(\varphi^n) & Fibonacci recursivo \\
T(n-1)+cn & \Theta(n^2) & peor caso de quicksort \\
\bottomrule
\end{tabularx}
\end{center}

Como control rápido sirve la forma general $T(n)=aT(n/b)+\Theta(n^d)$: si $d>\log_b a$ domina la raíz y $T=\Theta(n^d)$; si $d=\log_b a$, $T=\Theta(n^d\log n)$; si $d<\log_b a$ dominan las hojas y $T=\Theta(n^{\log_b a})$. En el libro no figura como ``teorema maestro'', así que en el parcial conviene justificar con (5.1)--(5.6), con el árbol o por inducción. En los cursos anteriores a 2017 se pedía resolver recurrencias concretas, casi siempre desenrollando.\visto{2P 2008 Ej.~1(b): $T(n/2)+n+\log n$\dycmas{Examen dic. 2009 Pr.~1(b): recurrencia de mergesort}}

\begin{problema}{Recurrencias resueltas en parciales viejos}
\textbf{(i)} $T(1)=1$, $T(n)=T(n/2)+n+\log_2 n$. Con $n=2^k$ y desenrollando: $T(n)=1+\sum_{i=0}^{k-1}\bigl(\frac{n}{2^i}+\log_2\frac{n}{2^i}\bigr)=1+(2n-2)+\frac{k(k+1)}2=2n-1+\frac{\log_2 n(\log_2 n+1)}{2}$, que es $\Theta(n)$. (Control: $T(2)=4$, $T(4)=10$.)

\textbf{(ii)} $T(1)=0$, $T(n)=T(\lfloor n/2\rfloor)+T(\lceil n/2\rceil)+n-1$ (comparaciones de mergesort en el peor caso). Para $n=2^k$: $T(n)=n\log_2 n-n+1$.

\textbf{(iii)} $T(N)=3T(N/3)+aN^2$. Con $N=3^m$, el nivel $j$ cuesta $aN^2/3^j$, así que $T(N)\le\frac32 aN^2+N\,T(1)=\Theta(N^2)$.

\textbf{(iv)} Mín--máx de una matriz $n\times n$ partiendo en 4 cuadrantes: $T(1)=0$, $T(n)=4T(n/2)+6$ (3 comparaciones para el máximo de 4 y 3 para el mínimo). Da $T(n)=2(n^2-1)$.
\end{problema}

Otra variante de los cursos viejos daba el código de una función recursiva que parte el rango $[i,j]$ en un \texttt{corte} y llama a una función \texttt{calcular} de costo conocido. Se pedía analizar las alternativas de un paso genérico (qué llamadas recursivas se ejecutan según lo que devuelvan) y plantear la recurrencia del peor caso. El peor caso es la rama que hace \emph{las dos} llamadas recursivas más el \texttt{calcular} más caro; el costo de \texttt{calcular} entra como término independiente.\visto{Examen feb. 2011 Ej.~1: $3T(N/3)+aN^2$\dycmas{Examen dic. 2013 Pr.~1: mín--máx de matriz}\dycmas{Examen jul. 2009 Ej.~1.2: recurrencia de función dada}\dycmas{Examen dic. 2011 Pr.~2(b): ídem, corte en mitades}}

\begin{trampa}
\begin{itemize}
  \item Olvidar la \textbf{condición inicial} de la recurrencia. En los criterios del examen de feb.~2025 figura explícitamente: ``debe incluir la condición inicial''.
  \item No explicar el origen de cada término. ``$2T(n/2)$ por las dos llamadas sobre mitades, $cn$ por el merge'' es lo mínimo que se espera.
  \item Confundir $2T(n/2)+c$ (lineal) con $2T(n/2)+cn$ ($n\log n$), o $T(n/2)+cn$ (lineal) con $T(n/2)+c$ (logarítmica).
  \item Hacer la inducción con la conjetura exacta cuando no cierra ($T(n)\le cn$ para $2T(n/2)+c$, $T(n)\le kn^{\log_2 3}$ para Karatsuba). Hay que restar un término de orden menor.
  \item No aclarar que se asume $n$ potencia de~2 (o de~3) cuando se usa.
\end{itemize}
\end{trampa}

% ---------------------------------------------------------------------
\section{Mergesort}

Mergesort es el ejemplo base del capítulo~5 y el que motiva (5.1). Divide el arreglo en dos mitades, ordena cada una recursivamente y las intercala (\emph{merge}).\visto{2P 2010 Ej.~3: implementar, comparaciones, ¿es óptimo?\dycmas{2P 2012 Ej.~2(b): implementar, estabilidad}\dycmas{2P 2014 Ej.~3: peor caso y caso medio}\dycmas{Examen jul. 2014 Ej.~2: merge sort dado \texttt{merge}}}

\begin{algorithmic}[1]
\Function{MergeSort}{$A,\,ini,\,fin$} \Comment{ordena $A[ini..fin]$}
  \If{$ini<fin$}
    \State $m\gets\lfloor (ini+fin)/2\rfloor$
    \State \Call{MergeSort}{$A,ini,m$};\; \Call{MergeSort}{$A,m+1,fin$};\; \Call{Merge}{$A,ini,m,fin$}
  \EndIf
\EndFunction
\Function{Merge}{$A,\,ini,\,m,\,fin$} \Comment{$A[ini..m]$ y $A[m+1..fin]$ ordenados}
  \State $i\gets ini$;\; $j\gets m+1$;\; $B\gets$ lista vacía
  \While{$i\le m$ \textbf{y} $j\le fin$}
    \If{$A[i]\le A[j]$} agregar $A[i]$ a $B$;\; $i\gets i+1$ \Comment{$\le$: estable}
    \Else\ agregar $A[j]$ a $B$;\; $j\gets j+1$
    \EndIf
  \EndWhile
  \State agregar a $B$ el resto de $A[i..m]$ o de $A[j..fin]$;\; copiar $B$ en $A[ini..fin]$
\EndFunction
\end{algorithmic}

\paragraph{Tiempo.} \textsc{Merge} hace a lo sumo $n-1$ comparaciones (cada comparación coloca un elemento, y el último se coloca sin comparar) y tiempo $\Theta(n)$. Dividir es $\Oh(1)$. Por lo tanto $T(n)\le 2T(n/2)+cn$, que es (5.1), y por (5.2) $T(n)=\Oh(n\log n)$. Como el merge siempre recorre todo, también es $\Omega(n\log n)$: mergesort es $\Theta(n\log n)$ en el peor, el mejor y el caso medio.

\paragraph{Comparaciones exactas en el peor caso.} $W(1)=0$, $W(n)=2W(n/2)+n-1$. Desenrollando, el nivel $j$ aporta $2^j(n/2^j-1)=n-2^j$, y sumando sobre $j=0,\dots,\log_2n-1$ queda $W(n)=n\log_2 n-(n-1)=n\log_2 n-n+1$. El peor caso del merge se da cuando los dos últimos elementos colocados vienen de mitades distintas (por ejemplo, las mitades se ``entrelazan''). Para $n$ cualquiera, $W(5)=W(2)+W(3)+4=1+3+4=8$.\visto{Examen jul. 2017 P.~3(d): mergesort con $n=5$ cuesta 8}

\paragraph{Correctitud.} Por inducción en el largo: los subarreglos de largo $\le 1$ están ordenados; por hipótesis inductiva las dos llamadas dejan ordenadas las mitades, y \textsc{Merge} produce la intercalación ordenada porque en cada paso coloca el menor de los dos mínimos restantes, que es el mínimo de lo que falta. \textbf{Estabilidad.} Un algoritmo de ordenación es \emph{estable} si elementos con igual clave quedan en el mismo orden relativo que tenían en la entrada. Mergesort es estable si, ante empate, \textsc{Merge} toma de la mitad izquierda (comparación $\le$). Con $<$ deja de serlo.

Un ejercicio combinó mergesort con insertion sort para los bloques chicos.\visto{2P 2016 Ej.~3: ordenación híbrida insertion + merge}

\begin{problema}{Ordenación híbrida: insertion sort en bloques de tamaño $k$ y luego merges}
Se parte el arreglo en $n/k$ bloques de tamaño $k$, se ordena cada uno con InsertionSort y se mezclan por etapas como en mergesort. El costo de las mezclas es $c_3n\log(n/k)$. (a)~Calcular el costo total. (b)~¿Para $k=\log n$ y para $k=\sqrt n$ se alcanza la cota inferior $\Theta(n\log n)$?

\textbf{Solución.} (a)~Cada InsertionSort de largo $k$ cuesta $c_2k^2$ en el peor caso; son $n/k$ bloques, total $c_2nk$. Costo: $c_2nk+c_3n\log n-c_3n\log k$. (b)~Con $k=\log n$: $(c_2+c_3)n\log n-c_3n\log\log n=\Theta(n\log n)$, sí. Con $k=\sqrt n$: $c_2n\sqrt n+\frac{c_3}{2}n\log n=\Theta(n^{3/2})$, no.
\end{problema}

Otras preguntas de los cursos viejos: ``escriba un algoritmo de sorting que use dividir y conquistar e indique su orden en mejor, peor y caso medio'' (mergesort: $\Theta(n\log n)$ en los tres; quicksort: $\Theta(n\log n)$, $\Theta(n^2)$, $\Theta(n\log n)$), o dar un \texttt{Ordenar(A,p,q)} genérico con \texttt{Combinar} de costo $n-1$ y pedir $T_W(n)$, el orden exacto y si es óptimo.\visto{Examen feb. 2015 Ej.~4: sorting por DyC, órdenes\dycmas{Examen dic. 2011 Ej.~4: $T_W$ y optimalidad de un mergesort}}

\begin{receta}[recurrencia de mergesort ``indicando la procedencia de cada término'']
$T(1)=0$ (o $\le c$), $T(n)=2T(n/2)+n$ (o $+cn$). ``El término $2T(n/2)$ surge de resolver los dos subproblemas: ordenar la primera mitad y ordenar la segunda mitad. El término $n$ es el tiempo de dividir y combinar: dividir es $\Oh(1)$ (calcular el medio) y combinar, el merge, es $\Oh(n)$.'' Por (5.2), $\Oh(n\log n)$.
\end{receta}

\begin{trampa}
\begin{itemize}
  \item Merge con un solo índice de fin, o copiando el auxiliar a $A[0..]$ en vez de a $A[ini..]$: el código de la solución oficial de jul.~2022 tiene ese desliz, no lo copies.
  \item Decir que mergesort es $\Theta(n^2)$ en el peor caso, o que su mejor caso es $\Theta(n)$ (eso es insertion sort).
  \item Usar $<$ en el merge y afirmar que es estable.
\end{itemize}
\end{trampa}

% ---------------------------------------------------------------------
\section{Cota inferior para ordenar: árboles de decisión}

Este tema venía junto con mergesort en los cursos hasta 2017 y reapareció en el 2P~2021 como ``encontrar el error'' en un argumento de cota inferior.\visto{2P 2011 Ej.~3: nodos internos y externos\dycmas{2P 2015 Ej.~3: AD, $m\ge n!$, $\log n!$}\dycmas{Examen jul. 2017 P.~3: hojas, altura, $n=5$}}

\begin{definicion}
El \emph{árbol de decisión} de un algoritmo de ordenación por comparaciones, para entradas de tamaño $n$, es un árbol binario donde cada nodo interno es una comparación $a_i:a_j$ que hace el algoritmo, sus dos hijos corresponden a los dos resultados posibles, y cada hoja (nodo externo) es una salida, es decir, la permutación que ordena la entrada. Un camino de la raíz a una hoja es una ejecución; su largo es la cantidad de comparaciones de esa ejecución.
\end{definicion}

\begin{itemize}
  \item La \textbf{altura} del árbol es la cantidad de comparaciones en el peor caso, $T_W(n)$.
  \item La \textbf{profundidad media de las hojas} (con entradas equiprobables) es el caso medio, $T_A(n)$.
  \item La mínima altura entre los árboles de todos los algoritmos que resuelven el problema es la complejidad del problema en el peor caso.
\end{itemize}

Los dos lemas del teórico y la cota que se deduce de ellos se pedían ``justificando detalladamente los pasos''.\visto{2P 2009 Ej.~3.2: cota inferior, peor y medio\dycmas{2P 2013 Ej.~3(a,b): ídem}\dycmas{Examen dic. 2010 Ej.~4: cota en caso medio}\dycmas{Examen feb. 2012 Ej.~4: lemas del AD}}

\begin{resultado}[Lemas y cota inferior]
\textbf{Lema 1.} Un árbol binario de altura $h$ tiene a lo sumo $2^h$ hojas. (Inducción en $h$.)

\textbf{Lema 2.} En un árbol binario con $m$ hojas, la suma de las profundidades de las hojas es al menos $m\log_2 m$.

\textbf{Cota en el peor caso.} Con $n$ elementos distintos, cada una de las $n!$ permutaciones necesita una hoja propia, así que $m\ge n!$. Por el lema~1, $2^h\ge m\ge n!$, de donde $T_W(n)=h\ge\lceil\log_2 n!\rceil$. Como $\log_2 n!\ge n\log_2 n-1{,}5n$ (dato que daban las letras), $T_W(n)=\Omega(n\log n)$.

\textbf{Cota en el caso medio.} Con las $n!$ entradas equiprobables, $T_A(n)=\frac1m\sum_i k_i\ge\frac{m\log_2 m}{m}=\log_2 m\ge\log_2 n!$, también $\Omega(n\log n)$.
\end{resultado}

\paragraph{Consecuencia: mergesort es óptimo en orden.}\visto{Examen jul. 2016 Ej.~2: $T(\textrm{mergesort})\in\Theta(\log n!)$\dycmas{Examen dic. 2009 Pr.~1(b.2): caso medio de mergesort}} $W(n)=n\log_2 n-n+1\le\log_2 n!+1{,}5n-n+1=\log_2 n!+0{,}5n+1$, y como $\log_2 n!=\Omega(n)$, $W(n)=\Theta(\log_2 n!)$. Optimalidad ``en orden'' no quiere decir exacta: para $n=5$, mergesort hace 8 comparaciones y $\lceil\log_2 120\rceil=7$, sin contradicción. El caso medio de mergesort está entre $\log_2 n!$ (cota inferior) y $W(n)$, así que también es $\Theta(n\log n)$.

\paragraph{Cómo dibujar el árbol de un algoritmo concreto para $n=3$.} Se ejecuta el algoritmo con la entrada simbólica $[a,b,c]$. Cada vez que compara, se abre un nodo con dos ramas y se sigue cada rama con el estado del arreglo actualizado (conviene anotar el arreglo en cada nodo, como pedían las letras). Insertion sort, para $n=3$, tiene altura 3 y 6 hojas. Selection sort hace siempre 3 comparaciones y su árbol tiene 8 hojas: dos corresponden a resultados contradictorios (por ejemplo $a<b$, $b<c$ y $c<a$) y son ejecuciones que nunca pueden darse.\visto{Examen feb. 2010 Ej.~4: AD de insertion sort\dycmas{Examen jul. 2010 Ej.~3: AD de selection sort}\dycmas{Examen dic. 2015 Ej.~4: hoja que nunca se da}}

Bubble sort con la entrada $[a,b,c]$ también hace 3 comparaciones y tiene hojas inalcanzables. Costos de referencia (comparaciones): insertion sort hace $n(n-1)/2$ en el peor caso (entrada en orden inverso) y $n-1$ en el mejor; selection sort hace siempre $n(n-1)/2$; bubble sort, $n(n-1)/2$ en la versión sin corte temprano. Ninguno es óptimo en el peor caso. La versión reciente del tema, en cambio, pide detectar cuándo la cota no se aplica.\visto{2P 2008 Ej.~2: bubble sort, AD, estabilidad\dycmas{Examen dic. 2013 Ej.~3: selection sort, ¿óptimo?}\dycmas{Examen dic. 2014 Ej.~2: peor caso de insertion}\dycmas{Examen jul. 2012 Ej.~4: insertion vs.\ selection}\dycmas{2P 2021 Ej.~3: error en la cota $\Omega(n\log n)$}}

\begin{problema}{Ordenar correos con a lo sumo 10 remitentes (2P 2021, Ej.~3)}
(a) Mostrar que un algoritmo que mantiene una lista ordenada de a lo sumo 10 listas (una por remitente) e inserta cada correo en la suya ordena en $\Oh(n)$. (b) Señalar el error en este argumento: ``el algoritmo debe distinguir entre $n!$ permutaciones, entonces $2^{M_n}\ge n!$ y $M_n=\Omega(n\log n)$''.

\textbf{Solución.} (a) Cada iteración busca entre a lo sumo 10 listas, $\Oh(1)$; insertar es $\Oh(1)$; concatenar las listas es $\Oh(1)$; el ciclo es $\Oh(n)$. (b) El error está en ``debe distinguir $n!$ permutaciones''. Con a lo sumo 10 claves posibles hay a lo sumo $10^n$ entradas distintas, y una misma permutación ordena muchas entradas distintas (por ejemplo, todas las que tienen claves repetidas en las mismas posiciones relativas). La cota $\log_2 n!$ vale cuando los elementos son distintos.
\end{problema}

\begin{trampa}
\begin{itemize}
  \item Decir que la cota inferior vale para \emph{todo} algoritmo de ordenación: vale para los que solo comparan pares de elementos. Counting sort o bucket sort no la cumplen.
  \item Confundir las hojas del árbol (salidas) con las permutaciones posibles: el árbol de un algoritmo concreto puede tener más de $n!$ hojas, algunas inalcanzables.
  \item Concluir de la cota que un algoritmo concreto hace exactamente $\lceil\log_2 n!\rceil$ comparaciones.
\end{itemize}
\end{trampa}

% ---------------------------------------------------------------------
\section{Conteo de inversiones}

Dada una secuencia $a_1,\dots,a_n$ de números distintos, una \emph{inversión} es un par de índices $i<j$ con $a_i>a_j$. La cantidad de inversiones mide qué tan lejos está la secuencia de estar ordenada: cero si está ordenada, $\binom n2$ si está en orden inverso. Sirve para comparar dos rankings: se escriben los índices del segundo ranking en el orden del primero y se cuentan las inversiones (cada inversión es un ``cruce'' de la figura de la letra).\visto{Examen jul. 2022 Ej.~1: inversiones entre rankings, $\Oh(n\log n)$}

Contar todos los pares es $\Theta(n^2)$. La idea del libro (sección~5.3) es modificar mergesort: además de ordenar cada mitad, cada llamada devuelve cuántas inversiones tiene. Las inversiones de la secuencia son las de la mitad izquierda $A$, más las de la mitad derecha $B$, más las \emph{cruzadas} ($a\in A$, $b\in B$, $a>b$). Las cruzadas se cuentan durante el merge: cuando se coloca un elemento de $B$, todos los que \emph{quedan} en $A$ son mayores que él, y cada uno forma una inversión con él.

\begin{algorithmic}[1]
\Function{Sort-and-Count}{$L$}
  \If{$|L|=1$} \Return $(0,L)$ \EndIf
  \State dividir $L$ en $A$ (primeros $\lceil n/2\rceil$) y $B$ (el resto)
  \State $(r_A,A)\gets$ \Call{Sort-and-Count}{$A$};\; $(r_B,B)\gets$ \Call{Sort-and-Count}{$B$}
  \State $(r,L)\gets$ \Call{Merge-and-Count}{$A,B$}
  \State \Return $(r_A+r_B+r,\;L)$
\EndFunction
\Function{Merge-and-Count}{$A,B$} \Comment{$A$ y $B$ ordenadas}
  \State $Cont\gets 0$;\; $C\gets$ lista vacía
  \While{$A$ y $B$ no vacías} \Comment{$a$, $b$: primeros de $A$ y $B$}
    \If{$a<b$} pasar $a$ al final de $C$
    \Else\ pasar $b$ al final de $C$;\; $Cont\gets Cont+|A|$ \Comment{lo que queda en $A$}
    \EndIf
  \EndWhile
  \State agregar a $C$ lo que quede de $A$ o de $B$;\; \Return $(Cont,C)$
\EndFunction
\end{algorithmic}

\paragraph{Correctitud.} Por inducción: las llamadas recursivas devuelven las inversiones internas de $A$ y de $B$ y las dejan ordenadas. Ordenar cada mitad no cambia las inversiones cruzadas (es un par por cada $a\in A$, $b\in B$ con $a>b$, sin importar las posiciones dentro de cada mitad). En el merge, cuando $b$ pasa a $C$, los elementos que siguen en $A$ son exactamente los $a\in A$ con $a>b$ (los menores ya salieron, porque $A$ está ordenada), así que cada inversión cruzada se cuenta una vez, en el momento en que sale su elemento de $B$.

\textbf{Tiempo.} \textsc{Merge-and-Count} es $\Oh(n)$: cada iteración hace trabajo constante y saca un elemento. Entonces $T(n)\le 2T(n/2)+cn$ y por (5.2) es $\Oh(n\log n)$.

\begin{ejemplo}
$I_2=(2,1,5,3,4)$. Mitades $(2,1)$ y $(5,3,4)$. En $(2,1)$ hay 1 inversión; en $(5,3,4)$ hay 2, porque al mezclar $(5)$ con $(3,4)$ salen el 3 y el 4 con un elemento restante en $A$ cada vez. Al mezclar $(1,2)$ con $(3,4,5)$ salen primero el 1 y el 2, así que no hay cruzadas. Total: 3.
\end{ejemplo}

\begin{receta}[inversiones o cualquier ``modificá mergesort'']
Escribí el algoritmo completo (las dos funciones). Señalá en el merge la línea donde se cuenta y \emph{por qué} se suma la cantidad de elementos restantes de la mitad izquierda. Para el tiempo: ``\textsc{Merge-and-Count} recorre cada mitad una vez con operaciones $\Oh(1)$ por elemento, $\Oh(n)$; la recursión cumple $T(1)\le c$, $T(n)\le 2T(n/2)+cn$; es la recurrencia (5.1) y por (5.2) es $\Oh(n\log n)$''. La solución oficial también aceptó desenrollar: $T(n)\le 2^kT(n/2^k)+kcn+(2^k-1)c$ con $k=\log_2 n$.
\end{receta}

\begin{trampa}
\begin{itemize}
  \item Sumar 1 en vez de la cantidad de elementos restantes de $A$ (eso cuenta pasos del merge, no inversiones).
  \item Olvidar devolver la suma de las tres partes, o no devolver el arreglo ordenado: la llamada de arriba necesita las mitades ordenadas.
  \item Copiar el código oficial de jul.~2022 tal cual: le falta el \texttt{return inv} en \texttt{mergeSort\_Inv}, el segundo ciclo de ``restos'' usa \texttt{fin1} donde va \texttt{fin2}, y la copia final debe ir a \texttt{A[ini+posAux]}.
\end{itemize}
\end{trampa}

% ---------------------------------------------------------------------
\section{Par de puntos más cercano}

Dado un conjunto $P=\{p_1,\dots,p_n\}$ de puntos del plano, se busca un par $p_i\neq p_j$ con distancia euclídea $d(p_i,p_j)$ mínima. La fuerza bruta es $\Theta(n^2)$. El algoritmo de la sección~5.4 es $\Oh(n\log n)$. En exámenes se preguntó describirlo, explicar los términos de su recurrencia e identificar su técnica.\visto{1P 2017 Ej.~3(a): definir el problema, términos de la recurrencia\dycmas{Examen feb. 2018 P.~4: qué técnica usa cada algoritmo}}

\paragraph{Algoritmo.}
\begin{enumerate}
  \item Al principio, una sola vez: construir $P_x$ y $P_y$, los puntos ordenados por $x$ y por $y$. Cuesta $\Oh(n\log n)$.
  \item \textbf{Dividir:} $Q$ = primera mitad de $P_x$ y $R$ = segunda mitad. La recta vertical $L$ pasa por el punto de $Q$ con mayor $x$. Construir $Q_x,Q_y,R_x,R_y$ recorriendo $P_x$ y $P_y$ en $\Oh(n)$.
  \item \textbf{Resolver} recursivamente en $Q$ y en $R$. Sea $\delta$ la menor de las dos distancias obtenidas.
  \item \textbf{Combinar:} sea $S$ el conjunto de puntos a distancia menor que $\delta$ de $L$ y $S_y$ esos puntos en orden de $y$ (se obtiene de $P_y$ en $\Oh(n)$). Para cada punto de $S_y$, calcular la distancia a los 15 siguientes de $S_y$. Devolver el mejor par entre el de $Q$, el de $R$ y el de la franja.
\end{enumerate}

\begin{resultado}[Sección 5.4 — por qué alcanza con 15 vecinos]
Si existen $q\in Q$ y $r\in R$ con $d(q,r)<\delta$, ambos están a distancia menor que $\delta$ de $L$, es decir, en $S$. Además, si $s,s'\in S$ cumplen $d(s,s')<\delta$, entonces están a lo sumo a 15 posiciones uno del otro en $S_y$.

\emph{Idea:} se divide la franja en cuadrados de lado $\delta/2$. Cada cuadrado tiene a lo sumo un punto (dos puntos en el mismo cuadrado estarían del mismo lado de $L$ y a distancia $\le\delta/\sqrt2<\delta$, absurdo). Si $s$ y $s'$ están separados por 16 o más posiciones en $S_y$, entre ellos hay al menos tres filas completas de cuadrados y $|y(s)-y(s')|\ge 3\delta/2>\delta$.
\end{resultado}

\paragraph{Tiempo.} Con $P_x,P_y$ preordenados, dividir y combinar son $\Oh(n)$: $T(n)\le 2T(n/2)+cn$, y con el ordenamiento inicial el total es $\Oh(n\log n)$. Si en cada llamada se ordena la franja por $y$, la recurrencia es $T(n)\le 2T(n/2)+cn\log n$ y da $\Oh(n\log^2 n)$. Esta es la recurrencia que daba la letra del 1P~2017, que pedía explicar sus términos: el 0 es el caso base, $2T(n/2)$ son las dos llamadas separadas por $L$ y $cn\log n$ es ordenar la franja.

\begin{trampa}
\begin{itemize}
  \item Decir que la franja tiene $\Oh(1)$ puntos: puede tener todos. Lo constante es la cantidad de vecinos a revisar por punto.
  \item Ordenar por $y$ en cada llamada y afirmar $\Oh(n\log n)$.
  \item Olvidar los pares que cruzan $L$ (sin la franja, el algoritmo es incorrecto).
\end{itemize}
\end{trampa}

% ---------------------------------------------------------------------
\section{Multiplicación de enteros (Karatsuba)}\label{sec:dyc-karatsuba}

Multiplicar dos enteros de $n$ bits con el algoritmo escolar es $\Theta(n^2)$. Escribiendo $x=x_1 2^{n/2}+x_0$ y $y=y_1 2^{n/2}+y_0$ (con $x_1,x_0,y_1,y_0$ de $n/2$ bits):
\[
xy=x_1y_1\,2^n+(x_1y_0+x_0y_1)\,2^{n/2}+x_0y_0 .
\]
Con cuatro productos recursivos queda $T(n)\le 4T(n/2)+cn$, que por (5.4) es $\Oh(n^{\log_2 4})=\Oh(n^2)$: no se gana nada. El truco es obtener el término del medio con un solo producto extra, usando $(x_1+x_0)(y_1+y_0)=x_1y_1+(x_1y_0+x_0y_1)+x_0y_0$.\visto{Examen dic. 2022 Ej.~2: diseñar, recurrencia y resolverla\dycmas{2P 2025 Ej.~1(a): recurrencia y cota, sin demostrar}}

\begin{algorithmic}[1]
\Function{MultiplicacionEnteros}{$x,y,n$}
  \If{$n>1$}
    \State sean $x=x_1 2^{n/2}+x_0$ e $y=y_1 2^{n/2}+y_0$
    \State $r_2\gets$ \Call{MultiplicacionEnteros}{$x_1,y_1,n/2$}
    \State $r_0\gets$ \Call{MultiplicacionEnteros}{$x_0,y_0,n/2$}
    \State $r_1\gets$ \Call{MultiplicacionEnteros}{$x_1+x_0,\,y_1+y_0,\,n/2$}$\;-\;r_2-r_0$
    \State \Return $r_2\cdot 2^n+r_1\cdot 2^{n/2}+r_0$
  \Else
    \State \Return $x\cdot y$ \Comment{producto de dos bits}
  \EndIf
\EndFunction
\end{algorithmic}

\paragraph{Recurrencia.}
\[
T(n)\le\begin{cases} c_b & \text{si } n=1,\\ 3T(n/2)+c_g n & \text{en otro caso.}\end{cases}
\]
El término $3T(n/2)$ son las tres llamadas con instancias de tamaño $n/2$. El término $c_gn$ sale de una cantidad fija de sumas, restas y desplazamientos (multiplicar por $2^n$ es agregar ceros), cada una $\Oh(n)$. $c_b$ es el caso base, multiplicar dos bits. Por (5.4) con $q=3$: $T(n)=\Oh(n^{\log_2 3})=\Oh(n^{1{,}59})$, y en particular $o(n^2)$.

\begin{problema}{Demostrar por inducción que $T(n)\le kn^{\log_2 3}-dn$ (examen dic.~2022)}
\textbf{Base:} $T(1)\le c_b$ y se necesita $c_b\le k-d$.

\textbf{Paso:} con la hipótesis para $n/2$,
\[
T(n)\le 3\Bigl(k\bigl(\tfrac n2\bigr)^{\log_2 3}-d\tfrac n2\Bigr)+c_gn
= kn^{\log_2 3}-dn+n\Bigl(c_g-\tfrac d2\Bigr),
\]
usando $3(n/2)^{\log_2 3}=n^{\log_2 3}$. Es $\le kn^{\log_2 3}-dn$ si $c_g-d/2\le 0$. Alcanza con elegir $d\ge 2c_g$ y $k\ge c_b+d$.

\textbf{Por qué restar $dn$:} con la conjetura $T(n)\le kn^{\log_2 3}$ queda $T(n)\le kn^{\log_2 3}+c_gn$, y el término $c_gn$ impide cerrar el paso inductivo.
\end{problema}

\paragraph{Lo mismo con matrices (Strassen).} Partiendo matrices $n\times n$ en cuatro bloques, el producto ingenuo hace 8 productos de bloques: $8T(n/2)+cn^2=\Theta(n^3)$. Strassen arma el resultado con 7 productos de bloques (y $\Theta(n^2)$ sumas): $T(n)=7T(n/2)+cn^2=\Theta(n^{\log_2 7})\approx n^{2{,}81}$. En el examen se pedía explicar la idea a nivel de operaciones matriciales, sin código.\visto{Examen jul. 2011 Ej.~2: producto de matrices y Strassen}

\begin{trampa}
\begin{itemize}
  \item Escribir $4T(n/2)$: con cuatro llamadas no hay mejora.
  \item Restar mal: $r_1$ es el producto de las sumas \emph{menos} $r_2$ y $r_0$. La solución oficial de dic.~2022 tiene la errata ``$-r_2-r_1$''.
  \item Escribir $3T(n/2)+c$: las sumas de números de $n$ bits no son de tiempo constante.
  \item $x_1+x_0$ puede tener $n/2+1$ bits. No cambia el orden (el libro lo comenta), pero si te lo preguntan, decilo.
\end{itemize}
\end{trampa}

% ---------------------------------------------------------------------
\section{Quicksort, selección y mediana}

\subsection{Quicksort}

Quicksort también es dividir y conquistar, pero el trabajo está en la división. \textsc{Partition} elige un pivote y reordena el subarreglo para que queden primero los menores, después el pivote en su lugar definitivo y después los mayores. Luego se ordena recursivamente cada lado. Combinar no cuesta nada.\visto{2P 2011 Ej.~3(d,e): Partition y AD de quicksort\dycmas{2P 2013 Ej.~3(c): quicksort, técnica, estabilidad}\dycmas{Examen feb. 2011 Ej.~4: quicksort, complejidad, estabilidad}\dycmas{Examen dic. 2012 Ej.~4: quicksort, peor caso}}

\begin{algorithmic}[1]
\Function{QuickSort}{$A,ini,fin$}
  \If{$ini<fin$}
    \State $p\gets$ \Call{Partition}{$A,ini,fin$}
    \State \Call{QuickSort}{$A,ini,p-1$};\; \Call{QuickSort}{$A,p+1,fin$}
  \EndIf
\EndFunction
\Function{Partition}{$A,ini,fin$} \Comment{pivote $A[ini]$}
  \State $piv\gets A[ini]$;\; $p\gets ini$
  \For{$i\gets ini+1$ \textbf{hasta} $fin$}
    \If{$A[i]<piv$} $p\gets p+1$;\; intercambiar $A[p]$ y $A[i]$ \EndIf
  \EndFor
  \State intercambiar $A[ini]$ y $A[p]$;\; \Return $p$
\EndFunction
\end{algorithmic}

\paragraph{Costo.} \textsc{Partition} sobre $m$ elementos hace $m-1$ comparaciones. En el \textbf{peor caso} el pivote es siempre el mínimo o el máximo (por ejemplo, arreglo ya ordenado con pivote al principio): $T(n)=T(n-1)+(n-1)$, total $n(n-1)/2=\Theta(n^2)$, y no es óptimo. En el \textbf{caso medio} (todas las permutaciones equiprobables, o pivote aleatorio) es $\Theta(n\log n)$, alrededor de $1{,}39\,n\log_2 n$ comparaciones. El libro lo analiza en la sección~13.5 con pivote aleatorio: el tiempo esperado es $\Oh(n\log n)$. Quicksort \textbf{no es estable}: los intercambios de \textsc{Partition} pueden invertir elementos iguales.

\begin{ejemplo}
$[5,2,7,4,9,1,8,3,6]$ con pivote 5 y la partición de arriba. Los menores que 5 (2, 4, 1, 3) se van llevando a las posiciones 1 a 4, y queda $[5,2,4,1,3,7,8,9,6]$. El intercambio final pone el pivote en la posición 4: $[3,2,4,1,\mathbf{5},7,8,9,6]$. Los subvectores a ordenar recursivamente son $[3,2,4,1]$ y $[7,8,9,6]$.\visto{Examen feb. 2016 Ej.~2: traza de la partición}
\end{ejemplo}

\subsection{Selección y mediana}

Para encontrar el $k$-ésimo menor elemento (la mediana es $k=\lceil n/2\rceil$) no hace falta ordenar. \textsc{Select}$(S,k)$ elige un separador $s$, arma $S^-=\{a\in S: a<s\}$ y $S^+=\{a\in S:a>s\}$ en $\Oh(n)$, y recursa en \emph{un solo} lado: si $|S^-|=k-1$ devuelve $s$; si $|S^-|\ge k$ busca el $k$-ésimo en $S^-$; si no, busca el $(k-1-|S^-|)$-ésimo en $S^+$. Si el separador deja de cada lado una fracción fija, por ejemplo cada lado con a lo sumo $\frac34 n$, entonces $T(n)\le T(\frac34 n)+cn$, una serie geométrica que da $\Oh(n)$. Con separador aleatorio el tiempo esperado es $\Oh(n)$ (sección~13.5).

El ejercicio típico de examen usa esta idea con la mediana como separador.\visto{Examen feb. 2025 Ej.~2: faltante con \textsc{separar}\dycmas{Simulacro 1P 2018 Ej.~3: faltante, $B$ sin ordenar}}

\begin{problema}{Elemento faltante con \textsc{separar} (examen feb.~2025, Ej.~2)}
$O$ es un arreglo ordenado sin repetidos con $n-1$ elementos ($n$ potencia de 2). $F$ es un conjunto no ordenado con todos los elementos de $O$ menos uno. Se dispone de \textsc{separar}$(A,s)$, que en $\Theta(|A|)$ devuelve $(A^<,A^>)$. Dar un algoritmo de dividir y conquistar $\Oh(n)$ que encuentre el faltante, y justificar el tiempo con una recurrencia.

\textbf{Solución.} Sea $m=n/2$, de modo que $O[m]$ es la mediana de $O$. Se calcula $(F^<,F^>)\gets$ \textsc{separar}$(F,O[m])$. Si $|F^<|=|F^>|$, falta $O[m]$. Si $|F^<|<|F^>|$, el faltante está en $O[1..m-1]$ y se recursa con $(O[1..m-1],F^<)$. Si no, se recursa con $(O[m+1..n-1],F^>)$. Ambos subarreglos tienen $n/2-1$ elementos, una instancia ``de tamaño $n/2$''. Caso base $n=2$: $F$ es vacío y se devuelve $O[1]$.

\textbf{Recurrencia:} $T(2)\le c$, $T(n)\le T(n/2)+cn$. Hay una sola llamada recursiva sobre tamaño $n/2$, y el término $cn$ sale de \textsc{separar} sobre $n-2$ elementos más operaciones $\Oh(1)$ (y la copia de medio arreglo, si se copia). Por (5.5) es $\Oh(n)$.

\textbf{Correctitud (idea):} $F^<\subseteq O[1..m-1]$ y $F^>\subseteq O[m+1..n-1]$, y cada uno tiene $m-1$ elementos, salvo el que contiene al faltante, que tiene $m-2$.
\end{problema}

Hay una variante de quicksort con dos conjuntos que solo se comparan entre sí.\visto{1P 2017 Ej.~3(b): tuercas y tornillos}

\begin{problema}{Tuercas y tornillos (1P 2017, Ej.~3(b))}
$n$ tuercas y sus $n$ tornillos, de diámetros distintos. Solo se puede comparar una tuerca con un tornillo. Emparejarlos con $\Oh(n\log n)$ comparaciones en promedio.

\textbf{Solución.} Es quicksort ``cruzado''. Se toma una tuerca $s_m$ como pivote y se la compara con todos los tornillos: quedan los más chicos $T_p$, los más grandes $T_g$ y su tornillo $t_m$. Con $t_m$ como pivote se separan las tuercas en $S_p$ y $S_g$. Se recursa en $(S_p,T_p)$ y en $(S_g,T_g)$. En promedio los tamaños son cercanos a $n/2$ y el costo es $\Oh(n\log n)$. El peor caso, cuando el pivote es siempre el menor o el mayor, es $\Theta(n^2)$.
\end{problema}

\begin{trampa}
\begin{itemize}
  \item \textbf{Ordenar $F$} para después recorrerlo: es $\Theta(n\log n)$ y no cumple $\Oh(n)$. La solución oficial lo marca expresamente.
  \item Pasar a \textsc{separar} el índice $m$ en vez del elemento $O[m]$, o tratar a $F$ como arreglo indexable: es un conjunto.
  \item Usar sumas (``suma de $O$ menos suma de $F$''): los elementos son de un dominio con orden pero sin suma definida.
  \item Hacer dos llamadas recursivas en selección: con una sola es $T(n/2)+cn=\Oh(n)$; con dos es $\Oh(n\log n)$.
\end{itemize}
\end{trampa}

% ---------------------------------------------------------------------
\section{Búsqueda binaria y variantes}

Búsqueda binaria en un arreglo ordenado compara $x$ con el elemento del medio y sigue en una sola mitad: $T(n)\le T(n/2)+c$, $T(1)\le c$, que da $\Oh(\log n)$ (por inducción, $T(n)\le c\log_2 n+c$). Es el caso $q=1$ con combinación constante. El patrón general sirve cada vez que se puede decidir en $\Oh(1)$ en qué mitad está la respuesta, aunque no se compare con un $x$ dado.\visto{Examen jul. 2016 Ej.~4: búsqueda binaria y su técnica\dycmas{1P 2021 Ej.~3: faltante en progresión, $\Oh(\log n)$}}

\begin{problema}{Faltante en una progresión aritmética (1P 2021, Ej.~3)}
$A$ era una progresión aritmética ordenada. Se le quitó un elemento que no es ni el primero ni el último, y quedó un arreglo de largo $n$, potencia de 2. Encontrar el faltante en $\Oh(\log n)$.

\textbf{Solución.} La diferencia es $d=(A[n]-A[1])/n$, porque hay $n$ saltos de $d$ entre el primero y el último del arreglo original. Para el rango $[i,j]$: si $A[j]-A[i]=(j-i)d$ no falta nada en el rango. Se compara la mitad izquierda $[i,m]$ con la derecha $[m+1,j]$. Si falta en una de ellas se recursa ahí; si no falta en ninguna, el faltante es $A[m]+d$, entre las dos mitades. Caso base: dos elementos, se devuelve $A[i]+d$. Recurrencia $T(n)\le T(n/2)+c$, $\Oh(\log n)$.
\end{problema}

Un ejercicio viejo pidió una ``búsqueda'' por dividir y conquistar sin suponer orden.\visto{Examen feb. 2014 Ej.~3: búsqueda DyC sin orden}

\begin{problema}{``Búsqueda'' por dividir y conquistar sin suponer orden (examen feb.~2014, Ej.~3)}
\textsc{Busqueda}$(A,ini,fin,x)$: si $ini=fin$ devuelve $A[ini]=x$; si no, devuelve \textsc{Busqueda}$(A,ini,m,x)$ \textbf{o} \textsc{Busqueda}$(A,m+1,fin,x)$. Descomposición: partir el rango a la mitad. Combinación: el ``o''. Caso base: rango de un elemento. Tiempo: $T(n)\le 2T(n/2)+c=\Oh(n)$. Sin orden no se puede descartar una mitad, así que no es más rápido que recorrer el arreglo. Sirve para practicar cómo identificar las partes de la técnica.
\end{problema}

\begin{trampa}
\begin{itemize}
  \item \textbf{Copiar el subarreglo} en cada llamada: cada copia cuesta $\Theta(n)$ y la recurrencia pasa a ser $T(n/2)+cn=\Oh(n)$. Para $\Oh(\log n)$ hay que pasar índices de inicio y fin. La solución del 1P~2021 lo pide explícitamente, y la sugerencia de la letra lo insinuaba.
  \item En el faltante en progresión, olvidar el caso en que el faltante está \emph{entre} las dos mitades.
\end{itemize}
\end{trampa}

% ---------------------------------------------------------------------
\section{Otros problemas de examen resueltos con dividir y conquistar}

\subsection{Mínimo y máximo con $3n/2-2$ comparaciones}

Recorrer el arreglo buscando mínimo y máximo por separado usa $2(n-1)$ comparaciones. Por dividir y conquistar: si $n=2$, una comparación decide cuál es el mínimo y cuál el máximo. Si $n>2$, se resuelve cada mitad y se combinan con dos comparaciones (mínimo de los mínimos, máximo de los máximos).\visto{Examen feb. 2019 Ej.~1: mín--máx, $C(n)=3n/2-2$, correctitud\dycmas{Examen feb. 2012 Ej.~3: \texttt{minmax} por DyC}}

\begin{algorithmic}[1]
\Function{MinMax}{$A$} \Comment{$n$ potencia de 2, $n>1$}
  \If{$n=2$}
    \If{\Call{EsMenor}{$A[1],A[2]$}} \Return $(A[1],A[2])$ \Else\ \Return $(A[2],A[1])$ \EndIf
  \EndIf
  \State $(min_1,max_1)\gets$ \Call{MinMax}{$A[1:n/2]$}
  \State $(min_2,max_2)\gets$ \Call{MinMax}{$A[n/2+1:n]$}
  \State $min\gets$ \textbf{si} \Call{EsMenor}{$min_1,min_2$} \textbf{entonces} $min_1$ \textbf{sino} $min_2$
  \State $max\gets$ \textbf{si} \Call{EsMenor}{$max_1,max_2$} \textbf{entonces} $max_2$ \textbf{sino} $max_1$
  \State \Return $(min,max)$
\EndFunction
\end{algorithmic}

\textbf{Conteo.} $C(2)=1$ y $C(n)=2C(n/2)+2$: dos llamadas sobre $n/2$ más las dos comparaciones de la combinación. Por inducción en $i$ con $n=2^i$: para $i=1$, $3\cdot2/2-2=1=C(2)$; en el paso, $C(n)=2(3n/4-2)+2=3n/2-2$. \textbf{Correctitud.} Inducción en $i$. Para $n=2$, la comparación identifica al mayor. Para $n>2$, por hipótesis inductiva las llamadas devuelven el mínimo y el máximo de cada mitad. Como las mitades cubren $A$, el mínimo global es el menor de $min_1,min_2$ y el máximo global es el mayor de $max_1,max_2$.

\subsection{Subarreglo de suma máxima}

Encontrar $i\le j$ con $A_i+\dots+A_j$ máxima. Todo intervalo está en la mitad izquierda, en la derecha, o cruza el medio. Los dos primeros casos se resuelven recursivamente. El que cruza se arma con la mejor suma que termina en $n/2$ (un ciclo hacia la izquierda acumulando) y la mejor que empieza en $n/2+1$ (un ciclo hacia la derecha), en $\Oh(n)$. Se devuelve el mejor de los tres: $T(n)\le 2T(n/2)+cn$, $\Oh(n\log n)$ por (5.2).\visto{1P 2019 Ej.~2: suma máxima, $\Oh(n\log n)$ y $\Oh(n)$}

\emph{Mejora a $\Oh(n)$:} cada llamada devuelve también la suma total $X$, el mejor prefijo $Y$ y el mejor sufijo $Z$ de su rango. Con los valores de las mitades $I$ y $D$: $X=X_I+X_D$, $Y=\max(Y_I,X_I+Y_D)$, $Z=\max(Z_D,X_D+Z_I)$, y la mejor suma que cruza es $Z_I+Y_D$. Todo en $\Oh(1)$, así que $T(n)\le 2T(n/2)+c=\Oh(n)$.

\subsection{Leche ``estropeada''}

Dadas $n$ temperaturas $t_1,\dots,t_n$, decidir si existen $i<j$ con $t_j-t_i>d$. Con dividir y conquistar: el par está en la mitad izquierda, en la derecha, o tiene $i$ a la izquierda y $j$ a la derecha. Este último caso existe si y solo si $\max(\text{derecha})-\min(\text{izquierda})>d$. Calculando mínimo y máximo con un recorrido, $T(n)\le 2T(n/2)+cn=\Oh(n\log n)$. Si cada llamada devuelve además el mínimo y el máximo de su rango, la combinación es $\Oh(1)$ y queda $2T(n/2)+c=\Oh(n)$. El ejercicio pedía también demostrar esto último por inducción, con $T(n)\le cn-c$.\visto{1P 2018 Ej.~2: temperaturas, $\Oh(n\log n)$ y luego $\Oh(n)$}

\subsection{Conjunto independiente de peso máximo en un árbol binario}

$p(T)$ es el peso máximo de un conjunto independiente de $T$ (ningún nodo es padre de otro) y $q(T)$ el máximo sin usar la raíz. Si $T$ es una hoja $v$: $p=w(v)$, $q=0$. Si no, con subárboles $T_I,T_D$: $q(T)=p(T_I)+p(T_D)$ y $p(T)=\max\{w(v)+q(T_I)+q(T_D),\;q(T)\}$. La unión de conjuntos independientes de $T_I$ y $T_D$ es independiente en $T$, porque ningún nodo de un subárbol es padre de uno del otro. El algoritmo devuelve el par $(p,q)$. Para un árbol perfecto con $n$ nodos: $f(1)\le c$, $f(n)\le 2f\bigl(\frac{n-1}2\bigr)+c$, y por inducción $f(n)\le cn$, porque $2c\frac{n-1}{2}+c=cn$.\visto{Examen feb. 2022 Ej.~1: independiente en árbol, $p(T)$ y $q(T)$} Este ejercicio es el puente con programación dinámica: los subproblemas no se repiten porque los subárboles son disjuntos, así que la recursión directa ya es lineal.

\subsection{Fibonacci recursivo: cuándo dividir y conquistar no sirve}

\textsc{FibDC}$(n)=$ \textsc{FibDC}$(n-1)+$\textsc{FibDC}$(n-2)$, con $F_0=F_1=1$ como en la solución oficial. La cantidad de sumas cumple $S(0)=S(1)=0$ y $S(n)=1+S(n-1)+S(n-2)$. Los valores para $n=0,\dots,10$ son $0,0,1,2,4,7,12,20,33,54,88$, es decir $S(n)=F_n-1$, que crece exponencialmente. La versión de programación dinámica con espacio $\Oh(1)$ hace $n-1$ sumas. La diferencia se explica porque los subproblemas \emph{se superponen} y el algoritmo recursivo resuelve muchas veces el mismo. Dividir y conquistar funciona bien cuando los subproblemas son disjuntos.\visto{Examen feb. 2017 P.~3: FibDC, $S(n)$ y tabla vs.\ FibPD\dycmas{Práctica 2P: el mismo ejercicio resuelto}}

\subsection{Cubrir una plaza $2^n\times 2^n$ con baldosas en L}

Una plaza de $N\times N$ con $N=2^n$ tiene una celda ocupada por un monumento, y hay que cubrir el resto con baldosas en forma de L (3 celdas). Se usan $(N^2-1)/3$ baldosas. Dividir y conquistar: se parte la plaza en 4 cuadrantes de $N/2\times N/2$. El monumento cae en uno; se coloca una L en el centro que ocupa una celda de cada uno de los otros tres, y esas celdas pasan a ser el ``monumento'' de su cuadrante. Se recursa en los 4. El caso base es $2\times2$ con una celda ocupada: una L. $T(N)=4T(N/2)+c=\Theta(N^2)$, lineal en el área.\visto{Examen dic. 2010 Pr.~2: baldosas en L por DyC}

\begin{receta}[cuando el problema no es uno del libro]
\begin{enumerate}
  \item Preguntate qué tiene que devolver cada llamada para que la combinación sea barata. Casi siempre hay que devolver \emph{más} que la respuesta (mín y máx, $p$ y $q$, prefijo y sufijo, el arreglo ordenado). Ese es el truco que pasa de $\Oh(n\log n)$ a $\Oh(n)$.
  \item Escribí la clasificación de casos de la combinación (``todo intervalo cae en una de estas tres categorías'') y justificá que no falta ninguno: es la base de la correctitud.
  \item Correctitud por inducción en $i$ con $n=2^i$, como en la solución de feb.~2019.
\end{enumerate}
\end{receta}

\begin{trampa}
\begin{itemize}
  \item Recalcular en la combinación algo que la recursión podía devolver, y perder el orden pedido.
  \item En mín--máx, contar mal el caso base: con $n=2$ es \emph{una} comparación, no dos.
  \item Olvidar el caso ``cruzado'' (intervalo o par que cruza el medio).
\end{itemize}
\end{trampa}

% ---------------------------------------------------------------------
\section{Mapa de lo que se preguntó}

Las fuentes van con la numeración de cada letra (Ex.~= examen, Ej.~= ejercicio, P.~= pregunta, Pr.~= problema). Hasta 2017 el curso usaba otro material: árboles de decisión, cota inferior de la ordenación y algoritmos de ordenación concretos aparecían en el segundo parcial. Entre 2017 y 2021 dividir y conquistar se tomaba en el primer parcial; esos ejercicios se incluyen porque son del mismo tipo que los recientes. Desde 2022 aparece en exámenes y, en 2025, en el segundo parcial.

{\small
\renewcommand{\arraystretch}{1.05}
\begin{tabularx}{\textwidth}{@{}>{\raggedright\arraybackslash}p{2.95cm}>{\raggedright\arraybackslash}p{2.75cm}>{\raggedright\arraybackslash}X@{}}
\toprule
Fuente & Tema & Qué se pedía \\
\midrule
2P 2008 Ej.~1(b) & Recurrencias & Resolver $T(n)=T(n/2)+n+\log_2 n$ \\
2P 2008 Ej.~2 & Árbol de decisión & Bubble sort: árbol $n=3$, estabilidad, optimalidad \\
\end{tabularx}\par\vspace{-\parskip}
\begin{tabularx}{\textwidth}{@{}>{\raggedright\arraybackslash}p{2.95cm}>{\raggedright\arraybackslash}p{2.75cm}>{\raggedright\arraybackslash}X@{}}
2P 2009 Ej.~3.2 & Cota inferior & Cota inferior en peor caso y caso medio \\
Ex. jul. 2009 Ej.~1.2 & Recurrencias & Función recursiva dada: recurrencia del peor caso \\
\end{tabularx}\par\vspace{-\parskip}
\begin{tabularx}{\textwidth}{@{}>{\raggedright\arraybackslash}p{2.95cm}>{\raggedright\arraybackslash}p{2.75cm}>{\raggedright\arraybackslash}X@{}}
Ex. dic. 2009 Pr.~1(b) & Mergesort & Resolver $T(\lfloor n/2\rfloor)+T(\lceil n/2\rceil)+n-1$; caso medio $\Theta(n\log n)$ \\
2P 2010 Ej.~3 & Mergesort & Implementar, comparaciones, optimalidad, árbol $n=3$ \\
Ex. feb. 2010 Ej.~4 & Árbol de decisión & Definición y árbol de insertion sort para $n=3$ \\
Ex. jul. 2010 Ej.~3 & Árbol de decisión & Selection sort y su árbol para $n=3$ \\
\end{tabularx}\par\vspace{-\parskip}
\begin{tabularx}{\textwidth}{@{}>{\raggedright\arraybackslash}p{2.95cm}>{\raggedright\arraybackslash}p{2.75cm}>{\raggedright\arraybackslash}X@{}}
Ex. dic. 2010 Ej.~4 & Cota inferior & Cota en el caso medio, $\Theta(n\log n)$ \\
Ex. dic. 2010 Pr.~2 & Otros DyC & Cubrir plaza $2^n\times2^n$ con baldosas en L \\
\end{tabularx}\par\vspace{-\parskip}
\begin{tabularx}{\textwidth}{@{}>{\raggedright\arraybackslash}p{2.95cm}>{\raggedright\arraybackslash}p{2.75cm}>{\raggedright\arraybackslash}X@{}}
2P 2011 Ej.~3 & AD, quicksort & Árbol de decisión; quicksort con \textsc{Partition} \\
Ex. feb. 2011 Ej.~1 & Recurrencias & Orden de $3T(N/3)+aN^2$ \\
Ex. feb. 2011 Ej.~4 & Quicksort & Implementar; complejidad y estabilidad \\
Ex. jul. 2011 Ej.~2 & Matrices & Costo del producto clásico; idea de Strassen \\
\end{tabularx}\par\vspace{-\parskip}
\begin{tabularx}{\textwidth}{@{}>{\raggedright\arraybackslash}p{2.95cm}>{\raggedright\arraybackslash}p{2.75cm}>{\raggedright\arraybackslash}X@{}}
Ex. jul. 2011 Ej.~4 & Árbol de decisión & Definición; insertion sort estable \\
Ex. dic. 2011 Ej.~4 & Mergesort & $T_W$, caso medio y optimalidad de un mergesort \\
Ex. dic. 2011 Pr.~2(b) & Recurrencias & Función recursiva dada: peor caso \\
\end{tabularx}\par\vspace{-\parskip}
\begin{tabularx}{\textwidth}{@{}>{\raggedright\arraybackslash}p{2.95cm}>{\raggedright\arraybackslash}p{2.75cm}>{\raggedright\arraybackslash}X@{}}
2P 2012 Ej.~2(b) & Mergesort & Implementar, árbol $n=3$, comparaciones, estabilidad \\
Ex. feb. 2012 Ej.~3 & Mín--máx & \texttt{minmax} de un arreglo por dividir y conquistar \\
Ex. feb. 2012 Ej.~4 & Árbol de decisión & Hojas $\le 2^h$, suma de profundidades, cota inferior \\
Ex. jul. 2012 Ej.~4 & Sort elemental & Insertion y selection sort: esfuerzo paso a paso \\
\end{tabularx}\par\vspace{-\parskip}
\begin{tabularx}{\textwidth}{@{}>{\raggedright\arraybackslash}p{2.95cm}>{\raggedright\arraybackslash}p{2.75cm}>{\raggedright\arraybackslash}X@{}}
Ex. dic. 2012 Ej.~4 & Quicksort & Implementar con \texttt{AUXILIAR}; peor caso \\
\bottomrule
\end{tabularx}

\medskip
\begin{tabularx}{\textwidth}{@{}>{\raggedright\arraybackslash}p{2.95cm}>{\raggedright\arraybackslash}p{2.75cm}>{\raggedright\arraybackslash}X@{}}
\toprule
Fuente & Tema & Qué se pedía \\
\midrule
2P 2013 Ej.~3 & Cota, quicksort & Cotas inferiores; quicksort y estabilidad \\
Ex. feb. 2013 Ej.~3 & Árbol de decisión & Selection sort; partes del árbol con $b<a$ \\
Ex. dic. 2013 Ej.~3 & Ordenación elemental & Selection sort: peor, mejor y medio; optimalidad \\
Ex. dic. 2013 Pr.~1 & Mín--máx & Mín y máx de matriz $n\times n$, conteo \\
\end{tabularx}\par\vspace{-\parskip}
\begin{tabularx}{\textwidth}{@{}>{\raggedright\arraybackslash}p{2.95cm}>{\raggedright\arraybackslash}p{2.75cm}>{\raggedright\arraybackslash}X@{}}
2P 2014 Ej.~3 & Mergesort & Implementar dado \texttt{Merge}; peor caso y caso medio \\
Ex. feb. 2014 Ej.~3 & Búsqueda & Búsqueda por dividir y conquistar en arreglo sin orden \\
\end{tabularx}\par\vspace{-\parskip}
\begin{tabularx}{\textwidth}{@{}>{\raggedright\arraybackslash}p{2.95cm}>{\raggedright\arraybackslash}p{2.75cm}>{\raggedright\arraybackslash}X@{}}
Ex. jul. 2014 Ej.~2 & Mergesort & Implementar sin vectores auxiliares, técnica usada \\
Ex. dic. 2014 Ej.~2 & Ordenación elemental & Insertion sort, peor caso, comparaciones \\
Ex. feb. 2015 Ej.~4 & Mergesort & Algoritmo de sorting por DyC; mejor, peor y medio \\
2P 2015 Ej.~3 & Árbol de decisión & Construcción, $m\ge n!$, $\log n!$, altura mínima con $n=5$ \\
\end{tabularx}\par\vspace{-\parskip}
\begin{tabularx}{\textwidth}{@{}>{\raggedright\arraybackslash}p{2.95cm}>{\raggedright\arraybackslash}p{2.75cm}>{\raggedright\arraybackslash}X@{}}
Ex. jul. 2015 Ej.~2 & Árbol de decisión & Selection sort, árbol para $n=3$ \\
Ex. dic. 2015 Ej.~4 & Árbol de decisión & Árbol de selection sort; hojas que nunca se dan \\
\end{tabularx}\par\vspace{-\parskip}
\begin{tabularx}{\textwidth}{@{}>{\raggedright\arraybackslash}p{2.95cm}>{\raggedright\arraybackslash}p{2.75cm}>{\raggedright\arraybackslash}X@{}}
Ex. feb. 2016 Ej.~2 & Quicksort & Traza de la partición de $[5,2,7,4,9,1,8,3,6]$ \\
2P 2016 Ej.~3 & Mergesort & Híbrido insertion + merge; $k=\log n$ vs.\ $\sqrt n$ \\
Ex. jul. 2016 Ej.~2 & Mergesort, cota & Orden de mergesort; $T(\text{mergesort})\in\Theta(\log n!)$ \\
Ex. jul. 2016 Ej.~4 & Técnicas & Qué resuelve la búsqueda binaria y con qué técnica \\
\end{tabularx}\par\vspace{-\parskip}
\begin{tabularx}{\textwidth}{@{}>{\raggedright\arraybackslash}p{2.95cm}>{\raggedright\arraybackslash}p{2.75cm}>{\raggedright\arraybackslash}X@{}}
Ex. feb. 2017 P.~3 & Fibonacci & $S(n)$ de FibDC; FibPD en espacio $\Oh(1)$; tabla comparativa \\
Ex. jul. 2017 P.~3 & Árbol de decisión & Hojas y altura; mergesort $n=5$ cuesta 8 vs.\ cota 7 \\
1P 2017 Ej.~3 & Par cercano, quicksort & Términos de $2T(n/2)+cn\log n$; tuercas y tornillos \\
\bottomrule
\end{tabularx}

\medskip
\begin{tabularx}{\textwidth}{@{}>{\raggedright\arraybackslash}p{2.95cm}>{\raggedright\arraybackslash}p{2.75cm}>{\raggedright\arraybackslash}X@{}}
\toprule
Fuente & Tema & Qué se pedía \\
\midrule
Ex. feb. 2018 P.~3 & Recurrencias & Tiempos de la recurrencia; la de mergesort \\
Ex. feb. 2018 P.~4 & Técnicas & Par más cercano: qué resuelve y técnica \\
Simulacro 1P 2018 Ej.~3 & Selección & Faltante entre $A$ ordenado y $B$ sin ordenar, $\Oh(n)$ \\
1P 2018 Ej.~2 & Otros DyC & Leche: $\Oh(n\log n)$; $2T(n/2)+c$ es $\Oh(n)$ \\
\end{tabularx}\par\vspace{-\parskip}
\begin{tabularx}{\textwidth}{@{}>{\raggedright\arraybackslash}p{2.95cm}>{\raggedright\arraybackslash}p{2.75cm}>{\raggedright\arraybackslash}X@{}}
\end{tabularx}\par\vspace{-\parskip}
\begin{tabularx}{\textwidth}{@{}>{\raggedright\arraybackslash}p{2.95cm}>{\raggedright\arraybackslash}p{2.75cm}>{\raggedright\arraybackslash}X@{}}
Ex. feb. 2019 Ej.~1 & Mín--máx & Algoritmo con $C(n)=2C(n/2)+2$; probar $3n/2-2$; correctitud \\
1P 2019 Ej.~2 & Otros DyC & Subarreglo de suma máxima en $\Oh(n\log n)$ \\
1P 2021 Ej.~3 & Búsqueda binaria & Faltante en progresión aritmética en $\Oh(\log n)$ \\
2P 2021 Ej.~3 & Cota inferior & Correos en $\Oh(n)$; error en cota $\Omega(n\log n)$ \\
\end{tabularx}\par\vspace{-\parskip}
\begin{tabularx}{\textwidth}{@{}>{\raggedright\arraybackslash}p{2.95cm}>{\raggedright\arraybackslash}p{2.75cm}>{\raggedright\arraybackslash}X@{}}
Ex. feb. 2022 Ej.~1 & Otros DyC & Independiente en árbol; $\Oh(n)$ por inducción \\
Ex. jul. 2022 Ej.~1 & Inversiones & Algoritmo tipo mergesort y prueba de $\Oh(n\log n)$ \\
Ex. dic. 2022 Ej.~2 & Karatsuba & Diseñar $\Oh(n^{\log 3})$; resolver por inducción \\
Ex. feb. 2025 Ej.~2 & Selección & Faltante con \textsc{separar}; $T(n/2)+\Oh(n)$, cita a (5.5) \\
\end{tabularx}\par\vspace{-\parskip}
\begin{tabularx}{\textwidth}{@{}>{\raggedright\arraybackslash}p{2.95cm}>{\raggedright\arraybackslash}p{2.75cm}>{\raggedright\arraybackslash}X@{}}
2P 2025 Ej.~1(a) & Karatsuba & Recurrencia de Karatsuba y cota ajustada \\
Ex. jul. 2026 Ej.~1(e) & Recurrencias & $qT(n/2)+cn$; resultado para $q=1$ y $q=2$ \\
\bottomrule
\end{tabularx}
}

% =====================================================================
%  Parte 2 — Programación dinámica
%  Fragmento: se incluye con \input desde el documento maestro.
%  Macros locales con prefijo pd.
% =====================================================================

\newcommand{\pdalgo}[1]{\par\medskip\noindent{\sans\small\bfseries\color{acento}#1}\par\nopagebreak\smallskip}
\newcommand{\pdargmin}{\operatorname*{arg\,min}}
\newcommand{\pdargmax}{\operatorname*{arg\,max}}
\newcommand{\pdbar}[1]{\overline{#1}}

\parte{pd}{Programación dinámica}{\OPT(j)=\max\{v_j+\OPT(p(j)),\ \OPT(j-1)\}}{Se arma una familia chica de subproblemas ordenados, se escribe cómo sale el óptimo de cada uno a partir de los más chicos y se llena una tabla. Es el tema que más se pregunta: hubo un ejercicio de programación dinámica en todos los segundos parciales desde 2018.}

% =====================================================================
\section{El método y cómo se corrige}
% =====================================================================

La programación dinámica (PD) sirve cuando el problema se puede escribir en
función de una cantidad polinomial de subproblemas que cumplen tres cosas:
tienen un orden natural de ``más chico'' a ``más grande''; el óptimo de cada
uno se obtiene con una cuenta sencilla a partir de óptimos de subproblemas
más chicos; y el problema original es uno de ellos (o sale de ellos de forma
directa). A diferencia de dividir y conquistar, los subproblemas se
\emph{solapan}: el mismo subproblema aparece en muchas ramas de la recursión,
así que se calcula una sola vez y se guarda en una tabla.

La propiedad que lo hace funcionar es el \emph{principio de optimalidad}: una
solución óptima de una instancia está formada por soluciones óptimas de las
subinstancias en que se descompone. Se demuestra siempre igual, ``cortando y
pegando'': si la parte que corresponde a la subinstancia no fuera óptima, la
reemplazo por una mejor y obtengo una solución mejor de la instancia
original, lo cual es absurdo.\visto{Examen feb. 2013 Ej. 2: enunciar principio de optimalidad}

En los exámenes viejos (hasta 2017) se pedía enunciarlo y demostrarlo para
un problema dado, por ejemplo la mochila o los caminos mínimos de Floyd. En los
actuales no se pide con ese nombre, pero la justificación de la recurrencia
es exactamente ese argumento, y es la parte que más puntos pierde.\visto{Examen dic. 2015 Ej. 2: mochila y principio de optimalidad}

\subsection{Qué espera la corrección}

Las soluciones oficiales de 2024 y 2025 traen la lista explícita de lo que se
evalúa. Es siempre la misma, sin importar el problema:\visto{2P 2024 Ej. 1 y 2P 2025 Ej. 3: criterios de corrección}

\begin{enumerate}[label=\textbf{\arabic*.}]
  \item \textbf{Definir la función.} Escribir $\OPT(\ldots)$ con su
  significado en palabras (``la ganancia máxima que se puede obtener desde el
  día $i$ hasta el día $n$'') y el \emph{rango} de cada parámetro. Si hay un
  parámetro de estado (la cuchilla instalada, los neumáticos usados), decir
  con precisión qué representa. Decir dónde queda la respuesta: $\OPT(n)$,
  $\OPT(1)$, $\min_X \OPT(n,X)$, $\max_j \OPT(j)$.
  \item \textbf{Casos base.} Los subproblemas que se resuelven sin consultar
  otros. Hacen falta tantos como ``saltos'' tenga la recurrencia: si
  $\OPT(i)$ usa $\OPT(i+7)$, hay que definir $\OPT(i)=0$ para
  $n<i\le n+7$.
  \item \textbf{Recurrencia justificada.} Se distingue por casos según la
  última (o la primera) decisión y se explica de dónde sale cada término. Lo
  ideal es mostrar las dos desigualdades: que cada término es el valor de
  alguna solución factible ($\ge$ o $\le$ según se maximice o minimice) y que
  el óptimo cae en alguno de los casos.
  \item \textbf{Algoritmo iterativo.} Una tabla con una celda por
  subproblema; se inicializan los casos base; se recorre en un orden tal que,
  al calcular una celda, las que usa ya estén calculadas; se devuelve la celda
  (o la combinación de celdas) que contiene la respuesta.
  \item \textbf{Reconstrucción.} Recorrer la tabla en sentido opuesto al de
  llenado; en cada paso decidir \emph{con la misma ecuación} con la que se
  calculó la celda, mantener el estado elegido y saltar las celdas que la
  decisión vuelve incompatibles.
  \item \textbf{Orden del tiempo de ejecución.} Cantidad de subproblemas por
  costo de calcular cada uno, más los precálculos (ordenar, calcular
  $p(j)$, errores $e_{i,j}$, etc.).
\end{enumerate}

\begin{resultado}[Plantilla de la justificación de una recurrencia]
Sea $S$ una solución óptima de la instancia $I$ (con valor $\OPT(I)$).
Mirando la última decisión de $S$ hay $k$ casos posibles. En el caso $c$, lo
que queda de $S$ al sacar esa decisión es una solución factible de la
subinstancia $I_c$, y tiene que ser óptima para ella: si hubiera una mejor,
le agrego la decisión y obtengo una solución de $I$ mejor que $S$. Entonces
$\OPT(I)=\text{costo}(c)+\OPT(I_c)$ para el caso en que cae $S$. Como
recíprocamente cada $\text{costo}(c)+\OPT(I_c)$ es el valor de una solución
factible de $I$, el óptimo es el mejor de los $k$ valores.
\end{resultado}

La versión rigurosa de este argumento, con las dos desigualdades por
separado, está en la solución oficial del ladrón de casas
(sección~\ref{pd:sec:wis}); vale la pena copiarle el formato.

\subsection{Memoización e iterativo}

Escribir la recurrencia como una función recursiva directa suele ser
exponencial, porque el mismo subproblema se resuelve muchas veces. Hay dos
formas de evitarlo, equivalentes en orden:
\begin{itemize}
  \item \textbf{Memoización (de arriba hacia abajo).} La función recursiva
  guarda cada resultado en un arreglo global $M$ y, antes de calcular,
  consulta si ya lo tiene.
  \item \textbf{Iterativa (de abajo hacia arriba).} Se llena la tabla en un
  orden compatible con las dependencias. Es lo que piden \emph{siempre} los
  parciales (``dé un algoritmo iterativo'').
\end{itemize}

El ejemplo canónico es Fibonacci. La versión dividir y conquistar
$F(n)=F(n-1)+F(n-2)$ hace $S(n)=1+S(n-1)+S(n-2)$ sumas, con $S(0)=S(1)=0$,
lo que da $S(n)=F(n)-1$, exponencial. La versión iterativa hace $n-1$
sumas y, si solo guarda los dos últimos valores, usa espacio
$\Oh(1)$.\visto{Examen feb. 2017 Preg. 3: Fibonacci con DyC y con PD}

\pdalgo{Fibonacci con PD en espacio constante}
\begin{algorithmic}[1]
\Function{FibPD}{$n$}
  \State $a \gets 0$;\ $b \gets 1$ \Comment{$a=F(i-1)$, $b=F(i)$}
  \For{$i=2$ \textbf{hasta} $n$}
    \State $c \gets a+b$;\ $a \gets b$;\ $b \gets c$
  \EndFor
  \State \Return $b$
\EndFunction
\end{algorithmic}

La moraleja que se pide explicar: la diferencia de crecimiento se debe a la
\emph{superposición de subproblemas}; la PD resuelve cada uno una sola vez.
El mismo truco de guardar solo la última fila (o las últimas $k$ celdas)
reduce la memoria de cualquier tabla cuando solo se pide el valor óptimo y
no la solución.\visto{2P 2017 Ej. 1: reducir la memoria de una tabla}

\begin{receta}[estructura de cualquier respuesta de PD]
\begin{enumerate}[leftmargin=1.2em]
  \item ``Para $\ldots$ (rango), definimos $\OPT(\ldots)$ como $\ldots$ (en
  palabras). La respuesta es $\ldots$''.
  \item ``Casos base: $\ldots$ porque $\ldots$''.
  \item ``Para $\ldots$, consideramos una solución óptima $S$ de $\ldots$.
  O bien $\ldots$ (caso 1), y entonces $\ldots$ es óptima para $\ldots$, con
  valor $\ldots$; o bien $\ldots$ (caso 2) $\ldots$. Por lo tanto
  $\OPT(\ldots)=\max\{\ldots\}$''.
  \item Algoritmo: casos base, ciclos en el orden correcto, una línea que
  aplica la recurrencia, \textbf{retornar}.
  \item Reconstrucción: \textbf{mientras} (no \textbf{para}) que recorre al
  revés y repite las comparaciones de la recurrencia.
  \item Orden: ``hay $\Theta(\ldots)$ subproblemas y cada uno se calcula en
  $\Oh(\ldots)$''.
\end{enumerate}
\end{receta}

\begin{trampa}
\begin{itemize}[leftmargin=1.2em]
  \item Una ``recurrencia'' que decide localmente lo mejor en cada paso (y
  por lo tanto podría construir la solución en la misma pasada) es un
  greedy disfrazado; casi siempre es incorrecta y se corrige como tal.
  \item Falta el parámetro que hace falta para que la decisión dependa de lo
  anterior (qué cuchilla está instalada, cuántas vueltas tienen los
  neumáticos).
  \item El significado de $\OPT$ cambia entre el lado izquierdo y el derecho
  de la ecuación.
  \item Casos base insuficientes: la recurrencia consulta $\OPT(i+2)$ o
  $\OPT(i-d)$ fuera del rango definido.
  \item La reconstrucción compara entre sí los valores de la tabla en cada
  posición (``si $\OPT(i,A)<\OPT(i,B)$ elijo $A$'') en lugar de seguir la
  decisión que produjo el valor.
  \item Algoritmo recursivo sin memoización cuando se pide iterativo.
  \item Orden ``$\Oh(nW)$, polinomial'': es \emph{pseudo}polinomial.
\end{itemize}
\end{trampa}

\subsection{Catálogo de patrones}

Casi todos los ejercicios de examen caen en alguno de estos moldes. Reconocer
el molde es la mitad del ejercicio: te dice qué parámetros usar, qué casos
base escribir y en qué orden llenar la tabla.

{\small
\begin{tabularx}{\textwidth}{@{}>{\raggedright\arraybackslash}p{2.9cm}>{\raggedright\arraybackslash}p{3.4cm}>{\raggedright\arraybackslash}X@{}}
\toprule
\textbf{Patrón} & \textbf{Subproblema} & \textbf{Ejemplos de examen}\\
\midrule
Prefijo con decisión binaria & $\OPT(i)$ sobre $1..i$; tomo $i$ o no & ladrón, baldosas, edificios, intervalos ponderados\\
Sufijo con saltos & $\OPT(i)$ sobre $i..n$; salto a $i+k$ & tarifa plana, pedidos, ravioles, dinosaurio\\
Prefijo con estado & $\OPT(i,X)$, $X$ = última elección & neumáticos A/B, cuchillas, F1\\
Segmentación & $\OPT(j)$; elijo dónde empieza el último tramo & fragmentos, prefijos, palíndromos\\
Último elemento & $\OPT(j)$ = mejor que termina en $j$ & subsecuencia creciente\\
Capacidad extra & $\OPT(i,w)$ con $w\le W$ & mochila, cambio, varillas\\
Dos secuencias & $\OPT(i,j)$ sobre prefijos de $x$ e $y$ & edición, alineación, ejecución paralela\\
Intervalo & $\OPT(i,j)$ sobre $x_i..x_j$ & matrices, paréntesis, palíndromo\\
Grilla & $\OPT(i,j)$ por casilla & sargento, minero, recolector\\
DAG o árbol & un valor por vértice, en orden topológico & camino más largo, conteo, conjunto independiente\\
Cantidad de aristas & $\OPT(i,v)$, a lo sumo $i$ aristas & Bellman--Ford, routers\\
\bottomrule
\end{tabularx}}

Las secciones que siguen recorren los moldes en ese orden. En cada una está la
teoría del libro, el pseudocódigo modelo y los ejercicios que se tomaron.
Las referencias (6.x) siguen la numeración del capítulo 6 de Kleinberg y
Tardos; si tu edición corre algún número, lo que importa es el enunciado.

% =====================================================================
\section{Intervalos ponderados y prefijos con decisión binaria}\label{pd:sec:wis}
% =====================================================================

\subsection{Weighted interval scheduling}

Hay $n$ pedidos; el pedido $j$ ocupa el intervalo $[s_j,f_j)$ y tiene valor
$v_j>0$. Se busca un subconjunto de pedidos disjuntos dos a dos de valor
máximo. Greedy por fin más temprano ya no sirve (solo sirve con valores
iguales). Se ordenan los pedidos por tiempo de fin,
$f_1\le f_2\le\dots\le f_n$, y se define $p(j)$ como el mayor índice $i<j$
tal que $i$ y $j$ son compatibles ($p(j)=0$ si no hay ninguno). Para
$0\le j\le n$, $\OPT(j)$ es el valor óptimo usando solo los pedidos
$1,\dots,j$; $\OPT(0)=0$.

Para el pedido $j$ hay dos casos: si está en la solución óptima $\mathcal{O}_j$,
ninguno de $p(j)+1,\dots,j-1$ puede estar (se solapan con $j$) y lo que queda
es una solución óptima para $1..p(j)$; si no está, $\mathcal{O}_j$ es una
solución óptima para $1..j-1$.

\begin{resultado}[(6.1)]
$\OPT(j)=\max\bigl(v_j+\OPT(p(j)),\ \OPT(j-1)\bigr)$.
\end{resultado}

\begin{resultado}[(6.2)]
El pedido $j$ pertenece a alguna solución óptima del conjunto
$\{1,\dots,j\}$ si y solo si $v_j+\OPT(p(j))\ge\OPT(j-1)$.
\end{resultado}

Implementada tal cual en forma recursiva (\textproc{Compute-Opt}), la
recurrencia es correcta (6.3, por inducción en $j$) pero exponencial: si
$p(j)=j-2$ para todo $j$, la cantidad de llamadas crece como Fibonacci. Con
memoización o en forma iterativa, cada valor se calcula una vez.

\pdalgo{Intervalos ponderados: memoización, iterativo y reconstrucción}
\begin{algorithmic}[1]
\Function{M-Compute-Opt}{$j$}
  \If{$j=0$} \Return $0$
  \ElsIf{$M[j]$ está definido} \Return $M[j]$
  \Else
    \State $M[j]\gets\max\bigl(v_j+\Call{M-Compute-Opt}{p(j)},\ \Call{M-Compute-Opt}{j-1}\bigr)$
    \State \Return $M[j]$
  \EndIf
\EndFunction
\Statex
\Function{Iterative-Compute-Opt}{\,}
  \State $M[0]\gets 0$
  \For{$j=1$ \textbf{hasta} $n$}
    \State $M[j]\gets\max\bigl(v_j+M[p(j)],\ M[j-1]\bigr)$
  \EndFor
\EndFunction
\Statex
\Function{Find-Solution}{$j$}
  \If{$j=0$} \Return $\emptyset$
  \ElsIf{$v_j+M[p(j)]\ge M[j-1]$}
    \State \Return $\{j\}\cup\Call{Find-Solution}{p(j)}$
  \Else\ \Return \Call{Find-Solution}{$j-1$}
  \EndIf
\EndFunction
\end{algorithmic}

\begin{resultado}[(6.4) y (6.5)]
Con los pedidos ordenados por fin y los $p(j)$ calculados,
\textproc{M-Compute-Opt}$(n)$ corre en $\Oh(n)$: cada llamada que no
devuelve un valor ya guardado llena una celda nueva y genera dos llamadas, y
hay $n$ celdas. Dado el arreglo $M$, \textproc{Find-Solution} devuelve una
solución óptima en $\Oh(n)$. Ordenar y calcular los $p(j)$ (con búsqueda
binaria) cuesta $\Oh(n\log n)$.
\end{resultado}

\subsection{El ladrón: la solución modelo}

Este es el ejercicio que conviene tener memorizado, porque la solución
oficial es la referencia de cómo se justifica una recurrencia. Las casas
$c_1,\dots,c_n$ están en fila, la casa $c_i$ tiene valor $v_i>0$ y no se
pueden robar dos contiguas. Es intervalos ponderados con $p(i)=i-2$: cada
casa se ``solapa'' solo con sus vecinas.\visto{2P 2019 Ej. 2: ladrón de casas contiguas}

\begin{problema}{Ladrón de casas no contiguas (2P 2019, Ej. 2)}
\textbf{Subproblema.} Para $0\le i\le n$, $\OPT(i)$ es la ganancia máxima
robando casas de $\{c_1,\dots,c_i\}$. La respuesta es $\OPT(n)$.

\textbf{Recurrencia.} $\OPT(0)=0$, $\OPT(1)=v_1$ y, para $i>1$,
\[\OPT(i)=\max\{\OPT(i-1),\ \OPT(i-2)+v_i\}.\]

\textbf{Justificación ($\ge$).} Sean $S_1\subseteq\{c_1..c_{i-1}\}$ y
$S_2\subseteq\{c_1..c_{i-2}\}$ compatibles de valores $\OPT(i-1)$ y
$\OPT(i-2)$. Tanto $S_1$ como $S_2\cup\{c_i\}$ son compatibles dentro de
$\{c_1..c_i\}$, así que $\OPT(i)$ es al menos el máximo de ambos valores.

\textbf{Justificación ($\le$).} Sea $S$ compatible con
$v(S)=\OPT(i)$. Si $c_i\notin S$, entonces $S\subseteq\{c_1..c_{i-1}\}$ y
$v(S)\le\OPT(i-1)$. Si $c_i\in S$, entonces $c_{i-1}\notin S$ y
$S\setminus\{c_i\}$ es compatible dentro de $\{c_1..c_{i-2}\}$, así que
$v(S)-v_i\le\OPT(i-2)$. En ambos casos $v(S)$ es a lo sumo el máximo.

\textbf{Algoritmo.} $M[0]\gets0$, $M[1]\gets v_1$; para $i=2..n$ aplicar la
recurrencia; devolver $M[n]$. Tiempo $\Oh(n)$.

\textbf{Reconstrucción.} $i\gets n$; \emph{mientras} $i\ge1$: si
$M[i]=M[i-1]$, $i\gets i-1$; si no, agregar $c_i$ e $i\gets i-2$.
\end{problema}

El mismo enunciado, con baldosas en lugar de casas (no se pueden retirar dos
baldosas seguidas), se tomó en diciembre de 2017. La solución guarda un
arreglo booleano de \emph{traceback} con la decisión tomada en cada $i$, que
es una alternativa válida a recalcular la comparación al
reconstruir.\visto{Examen dic. 2017 Prob. B: baldosas no contiguas}

Si la incompatibilidad alcanza a más vecinos, cambia solo el salto. En el
parcial de 2015 la restricción era una distancia mínima $d$ entre edificios,
así que el salto es $p(i)=i-d$; el ejercicio pedía también refutar un greedy
y corregir una implementación dada.\visto{2P 2015 Ej. 1: edificios a distancia mínima $d$}

\begin{problema}{Edificios separados por al menos $d$ cuadras (2P 2015, Ej. 1)}
La cuadra $i$ aloja $h_i$ habitantes (0 si no hay solicitud) y dos edificios
aceptados deben estar a distancia $\ge d$. Con $T(i)$ el máximo para las
cuadras $1..i$:
\[
T(0)=0,\qquad T(i)=\max\bigl\{h_i+T(\max\{0,i-d\}),\ T(i-1)\bigr\},\quad i\ge1.
\]
Es intervalos ponderados con $p(i)=\max\{0,i-d\}$. Se pedía además mostrar
que el greedy ``aceptar en orden decreciente de habitantes si es compatible''
falla: con $d=2$ y $h=(2,3,2)$ elige $\{2\}$ (valor 3) y el óptimo es
$\{1,3\}$ (valor 4).
\end{problema}

\subsection{Sufijos con saltos}

Muchas veces es más natural mirar la \emph{primera} decisión: $\OPT(i)$ es el
óptimo para los días $i,i+1,\dots,n$, se llena con $i$ decreciente y la
respuesta es $\OPT(1)$. Si la decisión del día $i$ ``consume'' $k$ días, la
recurrencia salta a $\OPT(i+k)$ y hacen falta casos base hasta
$n+k$.\visto{2P 2021 Ej. 1: tarifa plana de 7 días}

\begin{problema}{Tarifa por consumo o tarifa plana (2P 2021, Ej. 1)}
El día $i$ se paga $t_iK_c$ por consumo, o se contrata un bloque de 7 días
consecutivos por $K_p$ (aunque queden menos de 7 días). Para
$1\le i\le n+7$, $\OPT(i)$ es el costo óptimo de los días $i..n$:
\begin{gather*}
\OPT(i)=0\qquad (n<i\le n+7),\\
\OPT(i)=\min\{t_iK_c+\OPT(i+1),\ K_p+\OPT(i+7)\}\qquad (1\le i\le n).
\end{gather*}
El día $i$ solo puede pagarse de esas dos formas; con la plana quedan pagos
hasta $\min\{i+6,n\}$. Se llena con $i=n$ bajando hasta 1. Reconstrucción:
$i\gets1$; mientras $i\le n$, si $\OPT(i)=K_p+\OPT(i+7)$ agregar $i$ a $P$ e
$i\gets i+7$, si no $i\gets i+1$.
\end{problema}

\begin{problema}{Ravioles y sorrentinos (examen dic. 2024, Ej. 2)}
Cada día se fabrica un solo tipo de pasta; si el día $i$ se hacen
sorrentinos, los días $i+1$ e $i+2$ solo pueden ser ravioles. Las ventas
estimadas son $r_i$ y $s_i$ unidades, con ganancias unitarias $g_r$ y $g_s$.
Con $\OPT(i)$ la ganancia máxima de los días $i..n$ y $r_{n+1}=r_{n+2}=0$:
\begin{gather*}
\OPT(n+d)=0\qquad (1\le d\le3),\\
\OPT(i)=\max\bigl\{s_ig_s+(r_{i+1}+r_{i+2})g_r+\OPT(i+3),\ r_ig_r+\OPT(i+1)\bigr\}.
\end{gather*}
Si se elige sorrentinos el día $i$, los dos días siguientes quedan fijados y
lo que sigue desde $i+3$ no tiene restricciones. Reconstrucción con
\textbf{mientras}: si $\OPT(i)=r_ig_r+\OPT(i+1)$ poner $R$ y avanzar 1; si no,
poner $S,R,R$ y avanzar 3.
\end{problema}

El de pedidos de soporte de 2025 es el mismo molde, pero con tres casos que
dependen de los datos y no de la decisión, y la solución oficial lista
errores típicos de la reconstrucción que conviene leer.\visto{2P 2025 Ej. 3: pedidos que duran uno o dos días}

\begin{problema}{Pedidos de soporte de uno o dos días (2P 2025, Ej. 3)}
El pedido del día $i$ requiere $r_i$ horas y da ganancia $g_i$; el día $i$
hay $d_i$ horas. Si $r_i\le d_i$ se hace en el día; si
$d_i<r_i\le d_i+d_{i+1}$ ocupa también el día $i+1$ (y se pierde el pedido
$i+1$); si $r_i>d_i+d_{i+1}$ se rechaza. El día $n$ solo se acepta si
$r_n\le d_n$. Con $\OPT(i)$ la ganancia máxima de los días $i..n$:
\[
\OPT(n+1)=0,\qquad \OPT(n)=\begin{cases}g_n & r_n\le d_n\\ 0 & r_n>d_n\end{cases}
\]
\[
\OPT(i)=\begin{cases}
\OPT(i+1) & r_i>d_i+d_{i+1}\\
g_i+\OPT(i+1) & r_i\le d_i\\
\max\{\OPT(i+1),\ g_i+\OPT(i+2)\} & d_i<r_i\le d_i+d_{i+1}
\end{cases}\qquad 1\le i<n.
\]
Se llena desde $n-1$ hasta 1. En el ejemplo de la letra la respuesta es
$G=11$.

\textbf{Reconstrucción.} Con \textbf{mientras}: si $\OPT(i)=\OPT(i+1)$, se
rechaza y se avanza 1; si no, se acepta, y se avanza 1 si $r_i\le d_i$ y 2 en
otro caso. Hay que decidir mirando también $r$ y $d$, no solo los valores de
la tabla.
\end{problema}

\begin{trampa}[reconstrucción en sufijos con saltos]
\begin{itemize}[leftmargin=1.2em]
  \item Usar \textbf{para} $i=1..n$: no permite saltar el día que quedó
  ocupado. Con $n=2$, $g=(20,10)$, $d=(1,1)$, $r=(2,1)$ la tabla da
  $\OPT(2)=10$, $\OPT(1)=20$ y un \textbf{para} acepta los dos pedidos.
  \item Comparar solo valores: puede valer $\OPT(i)=g_i+\OPT(i+2)$ aunque
  $r_i>d_i+d_{i+1}$ (pedido imposible). Con $n=3$, $g=(10,10,5)$,
  $d=(1,1,1)$, $r=(3,2,2)$ se tiene $\OPT(3)=0$, $\OPT(2)=10$, $\OPT(1)=10$,
  y el pedido 1 se aceptaría por error.
  \item Olvidar el caso base extra ($n+1$, o $n+1..n+7$) cuando la recurrencia
  salta más de un lugar.
\end{itemize}
\end{trampa}

\subsection{Recurrencias booleanas}

Si la pregunta es ``¿se puede?'', $\OPT$ es booleana y el $\max$ se vuelve
$\lor$. La reconstrucción es más simple que en optimización: basta ir a
cualquier sucesor que valga \texttt{true}.\visto{Examen feb. 2020 Ej. 3: dinosaurio que salta obstáculos}

\begin{problema}{Dinosaurio que avanza 1 o salta 4 (examen feb. 2020, Ej. 3)}
Hay $n$ baldosas, libres o con obstáculo ($B[i]$). Desde $i$ se va a $i+1$ o a
$i+4$; se gana al pasar más allá de $n$. Con $M(i)$ verdadero si se puede
ganar partiendo de $i$:
\begin{gather*}
M(i)=\texttt{true}\qquad (n<i\le n+4),\\
M(i)=\bigl(B[i]=\text{libre}\bigr)\wedge\bigl(M(i+1)\vee M(i+4)\bigr)\qquad (1\le i\le n).
\end{gather*}
Se calcula con $i$ decreciente. Si $M(1)$ es falso no hay solución; si no,
desde $i=1$, mientras $i\le n$: si $M(i+4)$ saltar, si no avanzar.
\end{problema}

\begin{receta}[prefijos y sufijos]
\begin{itemize}[leftmargin=1.2em]
  \item Si la letra sugiere ``desde el día $i$ hasta el $n$'', usá sufijos y
  llená con $i$ decreciente; si sugiere ``los primeros $i$'', prefijos.
  \item Escribí los casos base para todos los índices fuera de rango que
  aparecen en la recurrencia (o agregá ceros ficticios, $r_{n+1}=0$, y decilo).
  \item En cada término explicá en una línea qué decisión representa y por
  qué lo que queda es independiente de esa decisión.
  \item Si se pide comparar con un greedy, el contraejemplo de tres elementos
  casi siempre alcanza.
\end{itemize}
\end{receta}

% =====================================================================
\section{Prefijos con un parámetro de estado}
% =====================================================================

Cuando el costo de una decisión depende de la decisión anterior (cambiar de
cuchilla cuesta $T$, cambiar neumáticos cuesta $P$), $\OPT(i)$ no alcanza: al
escribir la recurrencia no sabés qué se eligió en $i-1$. Se agrega un
parámetro que guarda exactamente esa información, y la tabla pasa a tener
una fila por cada valor del estado. Este molde salió en el segundo parcial
de 2024 con la mitad del puntaje.\visto{2P 2024 Ej. 1: cuchillas A y B}

\begin{problema}{Cuchillas A y B (2P 2024, Ej. 1)}
Hay $n$ cortes en orden sobre materiales $m_1..m_n$; cortar $m$ con cuchilla
$X\in\{A,B\}$ cuesta $C_X(m)$ y cambiar de cuchilla cuesta $T$. Al empezar
está instalada la $A$. Escribimos $\pdbar{X}$ para la cuchilla opuesta.

\textbf{Subproblema.} Para $1\le i\le n$ y $X\in\{A,B\}$, $\OPT(i,X)$ es el
costo mínimo de hacer las primeras $i$ tareas \emph{terminando con la
cuchilla $X$} (la tarea $i$ se hace con $X$).

\textbf{Recurrencia.}
\[
\OPT(1,A)=C_A(m_1),\qquad \OPT(1,B)=C_B(m_1)+T,
\]
\[
\OPT(i,X)=C_X(m_i)+\min\bigl\{\OPT(i-1,X),\ \OPT(i-1,\pdbar{X})+T\bigr\},\quad 1<i\le n.
\]
Si la tarea $i$ se hace con $X$, se paga el desgaste $C_X(m_i)$; las
primeras $i-1$ tareas deben resolverse de forma óptima terminando con $X$
(sin cambio) o con $\pdbar{X}$ (y se paga el cambio). La respuesta es
$\min\{\OPT(n,A),\OPT(n,B)\}$. Tiempo $\Oh(n)$.

\textbf{Reconstrucción.} $X\gets\pdargmin_X\OPT(n,X)$; para $i=n$ bajando
hasta 2: $\mathit{SOL}[i]\gets X$; si $\OPT(i-1,\pdbar{X})+T<\OPT(i-1,X)$ entonces
$X\gets\pdbar{X}$. Al final $\mathit{SOL}[1]\gets X$.
\end{problema}

La solución oficial de este ejercicio es la más útil para estudiar, porque
analiza tres recurrencias incorrectas que aparecieron en las pruebas. Las
tres fallan por lo mismo: la información de qué cuchilla se usó en la tarea
anterior no está en el subproblema, o está con un significado distinto del
que se usa en el lado derecho.

\begin{trampa}[las recurrencias incorrectas del 2P 2024]
\begin{itemize}[leftmargin=1.2em]
  \item $\OPT(i)=\OPT(i-1)+\min\{C_A(m_i),T+C_B(m_i)\}$ ``si la tarea $i-1$
  se hizo con $A$'': no hay forma de saber eso desde $\OPT(i-1)$; es greedy.
  \item ``El último cambio fue en la tarea $k$'': para sumar $T$ hay que
  saber que la tarea $k-1$ se hizo con la otra cuchilla, y eso no está en
  $\OPT(k-1)$.
  \item $\OPT(i,X)$ con $X$ = ``cuchilla instalada al empezar la tarea $i$''
  pero usando $\OPT(i-1,X)$ como si $X$ fuera la cuchilla con la que se hizo
  $i-1$: el parámetro cambia de significado entre los dos lados.
  \item Reconstruir eligiendo en cada $i$ la cuchilla con menor $\OPT(i,X)$:
  con $T=5$, $C_A=(10,10)$, $C_B=(1,100)$ la tabla es
  $\OPT(\cdot,A)=(10,20)$, $\OPT(\cdot,B)=(6,106)$ y ese método devuelve
  $[B,A]$ de costo 21, cuando el óptimo es $[A,A]$ de costo 20.
\end{itemize}
\end{trampa}

El de neumáticos de diciembre de 2019 es idéntico salvo por el caso base: el
tipo de neumático de la largada se elige libremente, así que
$\OPT(1,X)=t(1,X)$ para ambos $X$, y la parada cuesta
$P$.\visto{Examen dic. 2019 Ej. 2: neumáticos A o B por vuelta}

\pdalgo{Neumáticos A/B: costo mínimo y solución}
\begin{algorithmic}[1]
\State $\OPT[1,X]\gets t(1,X)$ para $X\in\{A,B\}$
\For{$i=2$ \textbf{hasta} $n$}
  \ForAll{$X\in\{A,B\}$}
    \State $\OPT[i,X]\gets t(i,X)+\min\{\OPT[i-1,X],\ \OPT[i-1,\pdbar{X}]+P\}$
  \EndFor
\EndFor
\State $X\gets\pdargmin_{X}\OPT[n,X]$
\For{$i=n$ \textbf{bajando hasta} $2$}
  \State $\mathit{SOL}[i]\gets X$
  \If{$\OPT[i-1,\pdbar{X}]+P<\OPT[i-1,X]$} $X\gets\pdbar{X}$ \EndIf
\EndFor
\State $\mathit{SOL}[1]\gets X$;\ \Return $\min_X\OPT[n,X]$ y $SOL$
\end{algorithmic}

Un estado puede ser también un número. En la estrategia de Fórmula 1 de 2022
el tiempo de cada vuelta depende de cuántas vueltas tienen los neumáticos,
así que ese contador entra en el subproblema. La solución oficial lo vincula
con el ejercicio 6.9 del libro (procesamiento con reinicios).\visto{2P 2022 Ej. 1: paradas de Fórmula 1}

\begin{problema}{Paradas en boxes de Fórmula 1 (2P 2022, Ej. 1)}
$s_{i,j}$ es el tiempo de la vuelta $i$ con neumáticos de $j$ vueltas de uso
($1\le j\le i\le n$); parar al final de una vuelta cuesta $P$. Con
$\OPT(i,j)$ el tiempo mínimo desde la vuelta $i$ hasta la $n$ si los
neumáticos llevan $j$ vueltas al correr la vuelta $i$:
\begin{gather*}
\OPT(n,j)=s_{n,j},\\
\OPT(i,j)=s_{i,j}+\min\bigl\{P+\OPT(i+1,1),\ \OPT(i+1,j+1)\bigr\}\qquad(1\le j\le i<n).
\end{gather*}
Respuesta $\OPT(1,1)$; se llena con $i$ decreciente y, para cada $i$, todos
los $j\le i$: $\Oh(n^2)$. Reconstrucción: $j\gets1$; para $i=1..n-1$, si
$\OPT(i,j)=s_{i,j}+P+\OPT(i+1,1)$ se para en $i$ y $j\gets1$, si no
$j\gets j+1$.

\textbf{Alternativa.} $\OPT'(i,j)$ = tiempo óptimo de las vueltas $1..i$
terminando con neumáticos de $j$ vueltas: $\OPT'(i,j)=\OPT'(i-1,j-1)+s_{i,j}$
si $j>1$, y $\OPT'(i,1)=P+s_{i,1}+\min_k\OPT'(i-1,k)$; respuesta
$\min_k\OPT'(n,k)$.
\end{problema}

En el programa anterior a 2018 apareció la versión ``dos agentes'': Ana y
Luis marcan los lugares $l_1,\dots,l_n$ en orden y en cada paso uno de los
dos va al siguiente lugar. El estado es dónde quedó el que no se movió.\visto{2P 2010 Ej. 2: dos corredores marcan lugares}

\begin{problema}{Dos integrantes marcan lugares (2P 2010, Ej. 2)}
$f(l,k)$, con $k<l$, es la distancia que falta recorrer si ya se marcaron
los lugares hasta $l_l$, uno de los integrantes está en $l_l$ y el otro en
$l_k$. Caso base: $f(n,k)=0$. En general, el lugar $l_{l+1}$ lo marca el que
está en $l_l$ o el que está en $l_k$:
\[
f(l,k)=\min\bigl\{d(l_l,l_{l+1})+f(l+1,k),\ d(l_k,l_{l+1})+f(l+1,l)\bigr\}.
\]
Hay $\Oh(n^2)$ subproblemas de costo $\Oh(1)$.
\end{problema}

\begin{receta}[cuándo agregar un estado]
Preguntate: ``para decidir lo que hago en el paso $i$, ¿qué necesito saber
de lo que pasó antes?''. Eso, y nada más, es el parámetro extra. Escribí
explícitamente si $X$ es la elección hecha \emph{en} el paso $i$ o la que
estaba vigente \emph{al empezar} el paso $i$, y usá el mismo significado en
ambos lados de la ecuación. En la reconstrucción llevá el estado en una
variable y actualizalo con la misma comparación de la recurrencia.
\end{receta}

% =====================================================================
\section{Segmentación: elegir dónde empieza el último tramo}
% =====================================================================

\subsection{Mínimos cuadrados segmentados}

Hay puntos $p_1,\dots,p_n$ ordenados por abscisa. Se quiere partirlos en
tramos consecutivos y ajustar una recta a cada tramo; el costo es la suma de
los errores cuadráticos $e_{i,j}$ de cada tramo $p_i..p_j$ más una constante
$C$ por tramo. La decisión ya no es binaria: el último tramo empieza en
algún $i$ entre 1 y $n$, y hay que probar todos (decisión \emph{de múltiples
opciones}).

\begin{resultado}[(6.6)]
Si el último tramo de la partición óptima es $p_i,\dots,p_n$, entonces
$\OPT(n)=e_{i,n}+C+\OPT(i-1)$.
\end{resultado}

\begin{resultado}[(6.7)]
Para los puntos $p_1,\dots,p_j$,
$\OPT(j)=\min_{1\le i\le j}\bigl(e_{i,j}+C+\OPT(i-1)\bigr)$, con
$\OPT(0)=0$; el tramo $p_i..p_j$ se usa en una solución óptima del
subproblema si y solo si el mínimo se alcanza en $i$.
\end{resultado}

\pdalgo{Segmented-Least-Squares y reconstrucción}
\begin{algorithmic}[1]
\State calcular $e_{i,j}$ para todo $1\le i\le j\le n$
\State $M[0]\gets 0$
\For{$j=1$ \textbf{hasta} $n$}
  \State $M[j]\gets\min_{1\le i\le j}\bigl(e_{i,j}+C+M[i-1]\bigr)$
\EndFor
\Statex
\Function{Find-Segments}{$j$}
  \If{$j>0$}
    \State hallar $i$ que minimiza $e_{i,j}+C+M[i-1]$
    \State emitir el tramo $p_i..p_j$;\ \Call{Find-Segments}{$i-1$}
  \EndIf
\EndFunction
\end{algorithmic}

Con los errores precalculados, el ciclo principal es $\Oh(n^2)$ ($n$
subproblemas, $\Oh(n)$ cada uno). Calcular todos los $e_{i,j}$ de forma
directa cuesta $\Oh(n^3)$ y puede bajarse a $\Oh(n^2)$ con sumas
acumuladas. En los parciales el costo del tramo suele darse como función
calculable en $\Oh(1)$.\visto{2P 2020 Ej. 2: cortar una secuencia en fragmentos}

\begin{problema}{Fragmentos con costo máximo menos mínimo (2P 2020, Ej. 2)}
Partir $x_1..x_n$ en fragmentos consecutivos; cada fragmento cuesta
$d_{i,j}=\max-\min$ de sus elementos y cada \emph{corte} cuesta $c$ (no cada
fragmento). Con $\OPT(j)$ el costo mínimo del prefijo $x_1..x_j$:
\[
\OPT(j)=\min\Bigl\{d_{1,j},\ c+\min_{2\le i\le j}\bigl(\OPT(i-1)+d_{i,j}\bigr)\Bigr\},\qquad 1\le j\le n,
\]
con la convención $\min\emptyset=\infty$ (para $j=1$ queda $d_{1,1}=0$). El
primer término es ``un solo fragmento''; el segundo, ``el último fragmento
empieza en $i>1$'', que agrega un corte. Orden $\Oh(n^2)$.

\textbf{Posiciones de inicio.} $j\gets n$; mientras $j>0$: si
$\OPT(j)=d_{1,j}$ entonces $i\gets1$, si no buscar $i\ge2$ con
$\OPT(j)=c+\OPT(i-1)+d_{i,j}$; marcar $p[i]\gets\texttt{true}$;
$j\gets i-1$.
\end{problema}

Ojo con la diferencia respecto del libro: allá la constante se paga por
tramo, acá por corte. Una forma limpia de unificar es definir
$\OPT(0)=-c$ y usar $\OPT(j)=\min_{1\le i\le j}(c+\OPT(i-1)+d_{i,j})$. En
febrero de 2024 se tomó otra variante, en la que el ``tramo'' es un prefijo
de una cadena fija.\visto{Examen feb. 2024 Ej. 2: concatenación de prefijos}

\begin{problema}{Concatenación de prefijos (examen feb. 2024, Ej. 2)}
Se quiere construir $Y=y_1..y_n$ como concatenación de prefijos no nulos de
$S=s_1..s_n$ que coincida con $X$ en la mayor cantidad de posiciones. Si el
último prefijo cubre $x_{i+1}..x_j$, aporta
$c(i,j)=|\{k: i<k\le j,\ x_k=s_{k-i}\}|$ coincidencias:
\[
\OPT(0)=0,\qquad \OPT(j)=\max_{0\le i<j}\bigl(\OPT(i)+c(i,j)\bigr),\quad 1\le j\le n.
\]
Precalcular $c$ ingenuamente cuesta $\Oh(n^3)$; con
$c(i,j)=c(i,j-1)+[x_j=s_{j-i}]$ baja a $\Oh(n^2)$, igual que el ciclo
principal.
\end{problema}

La segmentación en palíndromos del examen de febrero de 2026 es el mismo
molde con la condición de que cada tramo sea palíndromo. La solución oficial
la presenta explícitamente como ``parecida a mínimos cuadrados
segmentados''.\visto{Examen feb. 2026 Ej. 3: segmentación en palíndromos}

\begin{problema}{Segmentación mínima en palíndromos (examen feb. 2026, Ej. 3)}
Con $\mathit{esP}(i,j)$ verdadero si $a_i..a_j$ es palíndromo y $\OPT(j)$ el tamaño
mínimo de una segmentación en palíndromos de $a_1..a_j$:
\[
\OPT(0)=0,\qquad \OPT(j)=1+\min\{\OPT(i-1):\ 1\le i\le j,\ \mathit{esP}(i,j)\}.
\]
El rango de $i$ tiene que incluir $i=1$ (toda la cadena es un palíndromo) e
$i=j$ (tramos de largo 1, que siempre existen, así que el conjunto nunca es
vacío). Reconstrucción: $j\gets n$; mientras $j>0$, elegir $i$ con
$\mathit{esP}(i,j)$ y $\OPT(j)=1+\OPT(i-1)$, agregar $i$ al inicio de $S$,
$j\gets i-1$. Orden $\Oh(n^2)$ con $\mathit{esP}$ precalculada (que es a su vez una
PD sobre intervalos, sección~\ref{pd:sec:intervalos}).

El greedy ``el palíndromo inicial más largo posible'' falla: para
\texttt{dqdqdmmmmdq} da $\{\texttt{dqdqd},\texttt{mmmm},\texttt{d},\texttt{q}\}$
y el óptimo es $\{\texttt{dqd},\texttt{qdmmmmdq}\}$. Tomar el palíndromo
final más largo también falla.
\end{problema}

\begin{trampa}[segmentación]
\begin{itemize}[leftmargin=1.2em]
  \item Excluir $i=1$ o $i=j$ del rango del mínimo, dejando afuera ``un solo
  tramo'' o ``tramos de largo 1''.
  \item Sumar la constante también al primer tramo cuando se cobra por corte.
  \item Escribir $\OPT(j)=\OPT(j-1)+\ldots$: eso decide solo si $j$ empieza o
  no un tramo, sin saber dónde empezó el anterior.
\end{itemize}
\end{trampa}

% =====================================================================
\section{Subsecuencia creciente más larga}
% =====================================================================

Otra decisión de múltiples opciones: para que el subproblema tenga
estructura, se fija el \emph{último} elemento. Para $1\le j\le n$,
$\OPT(j)$ es el largo máximo de una subsecuencia creciente de $x_1..x_j$ que
\emph{termina en $x_j$}. Si la subsecuencia tiene largo mayor a 1, su
penúltimo elemento es algún $x_i$ con $i<j$ y $x_i<x_j$, y lo anterior es una
subsecuencia creciente óptima que termina en $x_i$.\visto{Examen dic. 2023 Ej. 2: subsecuencia creciente más larga}

\begin{problema}{Subsecuencia creciente más larga (dic. 2023; 2P 2016; feb. 2018)}
Sea $S_j=\{i: 1\le i<j,\ x_i<x_j\}$. Entonces
\[
\OPT(j)=\begin{cases}1 & S_j=\emptyset\\ 1+\max\{\OPT(i): i\in S_j\} & S_j\ne\emptyset\end{cases}
\]
y la respuesta es $\max_{1\le j\le n}\OPT(j)$, no $\OPT(n)$. Alternativa de
2018: $G_0=0$, $G_i=\max\{G_{i-1},L_i\}$ con $L_i=\OPT(i)$, y la respuesta es
$G_n$. Orden $\Oh(n^2)$.

\textbf{Reconstrucción.} $j\gets\pdargmax_j\OPT(j)$; $P\gets(x_j)$; para
$i=j-1$ bajando hasta 1: si $x_i<x_j$ y $\OPT(i)+1=\OPT(j)$, agregar $x_i$ al
principio de $P$ y $j\gets i$. También puede guardarse un arreglo $T[j]$ con el
índice anterior elegido.
\end{problema}

En 2016 se pedía, además, explicar qué subproblemas hacen falta para resolver
$P_i$ (los $P_j$ con $j<i$ y $A_j<A_i$) y por qué la versión recursiva
directa es exponencial: en el peor caso, la secuencia creciente, la llamada
para $j$ se repite $2^{n-j}$ veces.\visto{2P 2016 Ej. 2: largos $L_i$ de subsecuencias crecientes}

En febrero de 2018 se agregó una reducción: apilar la mayor cantidad de
bandejas rectangulares, donde cada una debe ser más chica en ancho y largo
que la de abajo. Se ordenan por ancho y se busca la subsecuencia creciente más
larga de los largos.\visto{Examen feb. 2018 Prob. B: LIS y pila de bandejas}

\begin{trampa}[subsecuencia creciente]
Devolver $\OPT(n)$ (la mejor que termina en el último elemento) en lugar del
máximo sobre todos los $j$; definir $\OPT(j)$ como ``la más larga en
$x_1..x_j$'' y después usar que termina en $x_j$; olvidar que con elementos
repetidos la condición es estricta ($x_i<x_j$).
\end{trampa}

% =====================================================================
\section{Mochila y suma de subconjuntos: una capacidad extra}
% =====================================================================

\subsection{Suma de subconjuntos}

Hay $n$ ítems con pesos enteros $w_i$ y una cota $W$; se busca un
subconjunto $S$ con $\sum_{i\in S}w_i\le W$ lo más grande posible. Con
$\OPT(i)$ sola no alcanza: si el ítem $n$ entra en la solución, lo que
queda es un problema sobre $1..n-1$ \emph{con capacidad $W-w_n$}, que no es
ninguno de los $\OPT(i)$. Hay que agregar la capacidad restante como
parámetro: $\OPT(i,w)$ es el óptimo con los ítems $1..i$ y cota $w$, para
$0\le i\le n$ y $0\le w\le W$.

\begin{resultado}[(6.8)]
Si $w<w_i$ entonces $\OPT(i,w)=\OPT(i-1,w)$. Si no,
$\OPT(i,w)=\max\bigl(\OPT(i-1,w),\ w_i+\OPT(i-1,w-w_i)\bigr)$.
Casos base: $\OPT(0,w)=0$ para todo $w$.
\end{resultado}

\pdalgo{Subset-Sum$(n,W)$ y reconstrucción}
\begin{algorithmic}[1]
\State $M[0,w]\gets 0$ para $0\le w\le W$
\For{$i=1$ \textbf{hasta} $n$}
  \For{$w=0$ \textbf{hasta} $W$}
    \If{$w<w_i$} $M[i,w]\gets M[i-1,w]$
    \Else\ $M[i,w]\gets\max\bigl(M[i-1,w],\ w_i+M[i-1,w-w_i]\bigr)$
    \EndIf
  \EndFor
\EndFor
\State $S\gets\emptyset$;\ $w\gets W$
\For{$i=n$ \textbf{bajando hasta} $1$}
  \If{$w\ge w_i$ \textbf{y} $M[i,w]=w_i+M[i-1,w-w_i]$}
    \State $S\gets S\cup\{i\}$;\ $w\gets w-w_i$
  \EndIf
\EndFor
\State \Return $M[n,W]$ y $S$
\end{algorithmic}

\begin{resultado}[(6.9), (6.10) y (6.11)]
\textproc{Subset-Sum}$(n,W)$ calcula correctamente el óptimo en tiempo
$\Oh(nW)$. Dada la tabla, el conjunto óptimo se reconstruye en $\Oh(n)$. La
mochila (cada ítem tiene además valor $v_i$, y se maximiza
$\sum_{i\in S}v_i$ sujeto a $\sum_{i\in S}w_i\le W$) se resuelve igual en
$\Oh(nW)$ con
$\OPT(i,w)=\max\bigl(\OPT(i-1,w),\ v_i+\OPT(i-1,w-w_i)\bigr)$ si $w_i\le w$.
\end{resultado}

El orden $\Oh(nW)$ es \emph{pseudopolinomial}: es polinomial en el
\emph{valor} de $W$, pero $W$ se escribe con $\log W$ bits, así que puede ser
exponencial en el tamaño de la entrada. No contradice que suma de
subconjuntos sea $\NP$-completo (capítulo 8), y en un parcial hay que decirlo
así y no ``polinomial''.

\subsection{Mochila 0/1}

La mochila es el problema de PD más preguntado del programa viejo (hasta
2017), siempre con la misma estructura: plantear la recurrencia, justificar
por qué greedy no sirve, implementar y, a veces, comparar con la versión
fraccionaria.\visto{Examen jul. 2016 Prob. C: mochila de Martín}

\begin{problema}{Mochila 0/1 (jul. 2016 Prob. C; dic. 2010 Prob. 1; 2P 2008 Ej. 1)}
Productos con precio $d_k$ y peso $w_k$, un solo ejemplar de cada uno,
capacidad $W$. Con $g(k,w)$ la ganancia óptima usando los productos
$1..k$ con peso máximo $w$:
\[
g(0,w)=0,\qquad
g(k,w)=\begin{cases}g(k-1,w) & w<w_k\\ \max\{g(k-1,w),\ g(k-1,w-w_k)+d_k\} & w\ge w_k.\end{cases}
\]
Se llena por filas $k=1..n$ y, dentro de cada fila, $w=0..W$; respuesta
$g(n,W)$.

\textbf{Greedy no sirve.} Con $W=10$ y ítems $(w,d)=(6,7),(5,5),(5,5)$,
elegir por mayor $d/w$ toma el primero y obtiene 7; el óptimo es 10.

\textbf{Versión fraccionaria.} Si se pueden llevar fracciones, greedy por
mayor $d_k/w_k$ es óptimo, en $\Oh(n\log n)$.
\end{problema}

La misma estructura aparece disfrazada: espacios de publicidad con precio y
rating y un presupuesto (2012, donde también se pedía mostrar que el greedy
por cociente falla) y la definición formal del problema más el principio de
optimalidad (diciembre de 2015).\visto{Examen feb. 2012 Ej. 2: publicidad con presupuesto}

\subsection{Ítems ilimitados: cambio de monedas}

Si cada ítem se puede usar cualquier cantidad de veces, al usar el ítem $i$
\emph{no} se pasa a $i-1$: se queda en la misma fila con menos capacidad.
Ese es el único cambio, y es el error más común.\visto{Examen dic. 2020 Ej. 3: cambio con menos billetes}

\begin{problema}{Cambio con la menor cantidad de billetes (dic. 2020 Ej. 3; jul. 2020 Ej. 3)}
Valores $1=w_1<w_2<\dots<w_n$, cantidad ilimitada, monto $C$. Con
$\OPT(i,R)$ la menor cantidad de billetes con valores en $\{w_1..w_i\}$ que
suman $R$:
\[
\OPT(i,0)=0,\qquad \OPT(1,R)=R,
\]
\[
\OPT(i,R)=\begin{cases}\OPT(i-1,R) & R<w_i\\ \min\{\OPT(i-1,R),\ 1+\OPT(i,R-w_i)\} & R\ge w_i.\end{cases}
\]
O no se usa ningún billete $w_i$, o se usa al menos uno y el resto
$R-w_i$ se paga de forma óptima con los mismos valores. Orden $\Oh(nC)$.

\textbf{Greedy falla.} Con $B=\{1,3,4\}$ y $C=6$, tomar siempre el mayor
billete da $4+1+1$ (3 billetes) y el óptimo es $3+3$. En el problema de las
varillas (longitudes $1,4,6$) las respuestas para $4,8,14,20$ son $1,2,3,4$
varillas; por ejemplo $14=6+4+4$ y $20=6+6+4+4$.

\textbf{Reconstrucción.} $i\gets n$, $R\gets C$; mientras $R>0$: si $i=1$,
usar $R$ billetes de 1; si no, si $R<w_i$ o $\OPT(i-1,R)<1+\OPT(i,R-w_i)$,
$i\gets i-1$; si no, usar un billete $w_i$ y $R\gets R-w_i$.
\end{problema}

También se puede usar una sola dimensión,
$\OPT(R)=1+\min\{\OPT(R-w_i): w_i\le R\}$ con $\OPT(0)=0$, que es
correcta porque el orden de los billetes no importa. La ventaja de la versión
con dos índices es que es la que sugiere la letra y la reconstrucción es más
directa.\visto{Examen jul. 2020 Ej. 3: varillas con mínimas soldaduras}

El greedy de monedas (tomar la de mayor valor que entra, tantas veces como
se pueda) es óptimo para sistemas como el uruguayo, pero no en general; en
febrero de 2018 se pedía identificar la técnica y mostrar que puede no
encontrar la solución óptima. La PD de arriba es la que resuelve el caso
general.\visto{Examen feb. 2018 Preg. 2: greedy de monedas}

\subsection{Variantes con cantidades acotadas o dos restricciones}

Si de cada ítem hay un stock limitado, la decisión pasa a ser \emph{cuántas}
unidades del último ítem usar, y se toma el mínimo (o el máximo) sobre esa
cantidad. Si además la capacidad debe llenarse exactamente, las celdas
imposibles valen $+\infty$ y la respuesta puede no existir.\visto{2P 2012 Ej. 1: carga exacta con stock}

\begin{problema}{Carga exacta con stock acotado (2P 2012, Ej. 1)}
Productos $0..n$ de volumen $T[p]$ con stock $Q[p]$; llenar \emph{exactamente}
la capacidad $C$ con la menor cantidad de unidades. Con $f(p,c)$ el mínimo
usando los productos $0..p$ para volumen exacto $c$ ($+\infty$ si es
imposible), una recurrencia posible es
\begin{gather*}
f(-1,0)=0,\qquad f(-1,c)=+\infty\quad (c>0),\\
f(p,c)=\min_{0\le q\le\min\{Q[p],\lfloor c/T[p]\rfloor\}}\bigl(q+f(p-1,c-q\,T[p])\bigr).
\end{gather*}
La respuesta es $f(n,C)$, y si vale $+\infty$ no hay solución. Orden
$\Oh(n\,C\,\max_pQ[p])$.
\end{problema}

Con dos recursos (por ejemplo dinero y peso) la tabla tiene una dimensión
por cada uno. En los exámenes viejos la letra solía sugerir cómo definir los
casos base, distinguiendo si el requerimiento mínimo ya se alcanzó o no.\visto{Examen jul. 2009 Ej. 2: menús con dos restricciones}

\begin{problema}{Dos restricciones (menús 2009; ringtones 2010; contenedor 2014)}
Si hay dos recursos (precio $\le M$ \emph{y} proteínas $\ge W$), cada uno es
un índice. Para minimizar calorías con menús $1..k$, presupuesto $m$ y
requerimiento proteico $w$:
\begin{gather*}
f(0,m,w)=\begin{cases}0 & w\le 0\\ +\infty & w>0\end{cases}\\
f(k,m,w)=\min\bigl\{f(k-1,m,w),\ c_k+f(k-1,m-p_k,\max\{0,w-q_k\})\bigr\},
\end{gather*}
donde la segunda opción solo existe si $p_k\le m$. Truncar el requerimiento
en 0 mantiene la tabla finita: $\Oh(nMW)$. El caso base distingue si el
requerimiento mínimo ya se cumplió, que es lo que pedía la sugerencia.
\end{problema}

Hay otras variantes del programa viejo que se resuelven igual. Invertir
$d$ pesos en $n$ empresas con interés $h(x,i)$ es una mochila en la que se
elige \emph{cuánto} asignar a la última empresa:
$f(i,x)=\max_{0\le y\le x}\bigl(h(y,i)+f(i-1,x-y)\bigr)$, en
$\Oh(nd^2)$.\visto{Examen feb. 2009 Ej. 3: invertir en $n$ empresas}

Cuando el objetivo tiene dos niveles (``gastar la mayor cantidad de dinero
y, entre esas, la menor cantidad de golosinas''), el valor de cada celda es
un par que se compara de forma lexicográfica. Cuando además del máximo se
exige un mínimo (programadores con experiencia al menos $E$), basta
calcular el óptimo y verificar al final que cumple la cota.\visto{Examen feb. 2016 Ej. 4: programadores con presupuesto}

\begin{trampa}[mochila]
\begin{itemize}[leftmargin=1.2em]
  \item En la versión ilimitada, usar $\OPT(i-1,R-w_i)$: eso es 0/1.
  \item Olvidar el caso $w<w_i$, que deja índices negativos.
  \item Decir que $\Oh(nW)$ es polinomial.
  \item Proponer greedy por cociente para la versión 0/1.
  \item Omitir cómo se reconstruye el conjunto, o reconstruirlo comparando
  $M[i,w]$ con $M[i,w-1]$ en lugar de con $M[i-1,w]$.
\end{itemize}
\end{trampa}

% =====================================================================
\section{Dos secuencias: alineación, edición y subsecuencias comunes}
% =====================================================================

\subsection{Alineación de secuencias}

Dadas $X=x_1..x_m$ e $Y=y_1..y_n$, una \emph{alineación} es un conjunto $M$
de pares $(i,j)$ que forma un emparejamiento sin cruces: no hay
$(i,j),(i',j')\in M$ con $i<i'$ y $j>j'$. Cada posición sin pareja (un hueco)
cuesta $\delta$ y cada par $(i,j)\in M$ cuesta $\alpha_{x_iy_j}$ (cero si las
letras son iguales). El subproblema es sobre los dos prefijos: $\OPT(i,j)$ es
el costo mínimo de alinear $x_1..x_i$ con $y_1..y_j$.

\begin{resultado}[(6.14)]
Sea $M$ una alineación de $X$ e $Y$. Si $(m,n)\notin M$, entonces la posición
$m$ de $X$ o la posición $n$ de $Y$ no está emparejada en $M$.
\end{resultado}

\begin{resultado}[(6.15) y (6.16)]
En una alineación óptima se cumple al menos una de: $(m,n)\in M$; la
posición $m$ de $X$ no está emparejada; la posición $n$ de $Y$ no está
emparejada. Por lo tanto, para $i,j\ge1$,
\[
\OPT(i,j)=\min\bigl\{\alpha_{x_iy_j}+\OPT(i-1,j-1),\ \delta+\OPT(i-1,j),\ \delta+\OPT(i,j-1)\bigr\},
\]
con $\OPT(i,0)=i\delta$ y $\OPT(0,j)=j\delta$.
\end{resultado}

La tabla se llena por filas (o por columnas) en $\Oh(mn)$ tiempo y espacio.
El valor también es el largo del camino mínimo de $(0,0)$ a $(m,n)$ en la
grilla $G_{XY}$, con aristas horizontales y verticales de costo $\delta$ y
diagonales de costo $\alpha$ (6.17). Para el \emph{valor} alcanza con
guardar dos columnas (espacio $\Oh(m)$); para recuperar la alineación en
espacio lineal está el algoritmo de Hirschberg, que combina esa idea con
dividir y conquistar en tiempo $\Oh(mn)$. Saber ubicar el problema (qué
resuelve y con qué técnica) se preguntó como teoría.\visto{Examen feb. 2018 Preg. 4: qué técnica usa alineamiento}

\subsection{Distancia de edición}

Es la alineación con $\delta=1$ y $\alpha_{pq}=1$ si $p\ne q$ (0 si son
iguales): un par desigual es una sustitución, una posición de $x$ sin pareja
es un borrado y una de $y$ sin pareja es una inserción.\visto{Examen feb. 2019 Ej. 2: distancia de edición}

\begin{problema}{Distancia de edición (examen feb. 2019, Ej. 2)}
\textbf{Base.} $\OPT(0,j)=j$ (hay que insertar $j$ caracteres, y no se puede
con menos) y $\OPT(i,0)=i$ (borrar $i$).

\textbf{Lema que se pedía.} En un emparejamiento óptimo $S$, la posición $n$
de $x$ o la $m$ de $y$ está emparejada: si ninguna lo estuviera, agregar
$(n,m)$ cambia un borrado y una inserción por una sustitución y baja el
costo. Si ambas lo están, es entre sí, porque si no habría un cruce.

\textbf{Recurrencia.} Con $\delta_{i,j}=[x_i\ne y_j]$,
\[
\OPT(i,j)=\min\bigl\{\OPT(i-1,j-1)+\delta_{i,j},\ \OPT(i-1,j)+1,\ \OPT(i,j-1)+1\bigr\},\quad i,j>0.
\]
Cada término corresponde a: $(i,j)\in S$; $x_i$ se borra; $y_j$ se inserta.
En los tres casos lo que queda de $S$ debe ser óptimo para los prefijos
correspondientes (corte y pegado). Orden $\Oh(nm)$.
\end{problema}

En diciembre de 2018 apareció el mismo molde con otra historia: dos programas
que se ejecutan a la vez y comparten un ciclo cuando la próxima instrucción
coincide. La letra preguntaba primero por los casos base, que es la forma
habitual de arrancar estos ejercicios.\visto{Examen dic. 2018 Ej. 3: ejecución paralela de dos programas}

\begin{problema}{Ejecución en paralelo de dos programas (examen dic. 2018, Ej. 3)}
En cada ciclo, si las próximas instrucciones de $X$ e $Y$ son iguales se
ejecutan juntas; si no, se ejecuta una sola. Con $\OPT(i,j)$ los ciclos
mínimos para las primeras $i$ de $X$ y $j$ de $Y$:
\begin{gather*}
\OPT(0,j)=j,\qquad \OPT(i,0)=i,\\
\OPT(i,j)=\begin{cases}1+\OPT(i-1,j-1) & x_i=y_j\\ 1+\min\{\OPT(i-1,j),\OPT(i,j-1)\} & x_i\ne y_j.\end{cases}
\end{gather*}
Cuando $x_i=y_j$ conviene ejecutarlas juntas: cualquier ejecución óptima en
la que no lo están se transforma en otra igual de buena postergando la
primera al último ciclo. Observación: el resultado es $n+m-\mathrm{LCS}(X,Y)$.
\end{problema}

\subsection{Subsecuencia y substring común más largos}

La subsecuencia común más larga (LCS) es el mismo molde con máximo:
$\OPT(i,j)=1+\OPT(i-1,j-1)$ si $x_i=y_j$, y
$\max\{\OPT(i-1,j),\OPT(i,j-1)\}$ si no, con bordes en 0. El parcial de 2017
daba el código de la variante para \emph{substrings} (caracteres
consecutivos) y pedía descifrarlo.\visto{2P 2017 Ej. 1: substring común más largo}

\begin{problema}{Substring común más largo (2P 2017, Ej. 1)}
El código usa dos tablas. $\mathit{Mp}[i,j]$ es el largo del sufijo común más largo
de $A_1..A_i$ y $B_1..B_j$ (el substring común que \emph{termina} en $i$ y
en $j$); $\mathit{Mt}[i,j]$ es el largo del substring común más largo de esos
prefijos:
\begin{gather*}
\mathit{Mp}[i,j]=\begin{cases}1+\mathit{Mp}[i-1,j-1] & A_i=B_j\\ 0 & A_i\ne B_j\end{cases}\\
\mathit{Mt}[i,j]=\max\{\mathit{Mp}[i,j],\mathit{Mt}[i-1,j],\mathit{Mt}[i,j-1]\},
\end{gather*}
con fila y columna 0 en cero. Para $A=\texttt{la}$, $B=\texttt{ala}$:

\begin{center}
\begin{tabular}{c|ccc}
$\mathit{Mp}$ & a & l & a\\\hline
l & 0 & 1 & 0\\
a & 1 & 0 & 2
\end{tabular}\qquad
\begin{tabular}{c|ccc}
$\mathit{Mt}$ & a & l & a\\\hline
l & 0 & 1 & 1\\
a & 1 & 1 & 2
\end{tabular}
\end{center}

El resultado es 2 (\texttt{la}). La versión recursiva sin tabla recalcula
los mismos $\mathit{Mp}(i,j)$ muchas veces. El espacio $\Oh(nm)$ baja a $\Oh(m)$
guardando solo la fila anterior.
\end{problema}

En 2013 se pidió la alineación de dos secuencias de cubos de colores con un
cubo ``comodín'' para los huecos. Es exactamente la sección 6.6 con la
función de costos de la letra.\visto{Examen jul. 2013 Prob. 1: alinear cubos de colores}

\begin{receta}[dos secuencias]
\begin{itemize}[leftmargin=1.2em]
  \item Definí $\OPT(i,j)$ sobre \emph{prefijos} $x_1..x_i$, $y_1..y_j$ con
  $0\le i\le n$, $0\le j\le m$, y escribí la fila y la columna 0.
  \item Justificá con ``¿qué pasa con las últimas posiciones $i$ y $j$?'': se
  emparejan entre sí, $x_i$ queda solo o $y_j$ queda solo.
  \item Llená con dos \textbf{para} anidados crecientes y devolvé $\OPT(n,m)$.
  \item Para reconstruir, partí de $(n,m)$ y en cada celda andá a la vecina
  (diagonal, arriba o izquierda) cuyo término dio el mínimo.
\end{itemize}
\end{receta}

% =====================================================================
\section{\texorpdfstring{Intervalos $i..j$}{Intervalos i..j}}\label{pd:sec:intervalos}
% =====================================================================

Cuando una decisión parte el problema en un pedazo izquierdo y uno derecho
que \emph{no} son prefijos, los subproblemas tienen que ser intervalos
arbitrarios $x_i..x_j$. Hay $\Theta(n^2)$ de ellos, y la tabla se llena por
largo creciente $\ell=j-i$ (o con $i$ decreciente y $j$ creciente), para que
los intervalos más cortos estén calculados.

\subsection{Estructura secundaria del ARN}

Una secuencia $B=b_1..b_n$ sobre $\{A,C,G,U\}$; una estructura secundaria es
un conjunto $S$ de pares $(i,j)$ tal que: no hay giros cerrados ($i<j-4$);
cada par es $\{A,U\}$ o $\{C,G\}$; ninguna base está en dos pares; y no hay
cruces (no existen $(i,j),(k,l)\in S$ con $i<k<j<l$). Se busca $|S|$
máximo.

Con $\OPT(j)$ sobre prefijos no alcanza. Si $b_j$ se empareja con $b_t$, por
la condición de no cruce el problema se parte en $b_1..b_{t-1}$ y
$b_{t+1}..b_{j-1}$, y el segundo no es un prefijo. Por eso se usa
$\OPT(i,j)$ sobre $b_i..b_j$.

\begin{resultado}[Recurrencia del ARN (sección 6.5)]
$\OPT(i,j)=0$ si $i\ge j-4$. Si no,
\[
\OPT(i,j)=\max\Bigl(\OPT(i,j-1),\ \max_{t}\bigl(1+\OPT(i,t-1)+\OPT(t+1,j-1)\bigr)\Bigr),
\]
donde $t$ recorre los $i\le t<j-4$ tales que $b_t$ y $b_j$ forman un par
permitido. Hay $\Oh(n^2)$ subproblemas de costo $\Oh(n)$: $\Oh(n^3)$.
\end{resultado}

\pdalgo{Llenado por largo de intervalo (ARN)}
\begin{algorithmic}[1]
\State $M[i,j]\gets0$ para todo $i\ge j-4$
\For{$k=5$ \textbf{hasta} $n-1$} \Comment{$k=j-i$}
  \For{$i=1$ \textbf{hasta} $n-k$}
    \State $j\gets i+k$
    \State $M[i,j]\gets\max\bigl(M[i,j-1],\ \max_t(1+M[i,t-1]+M[t+1,j-1])\bigr)$
  \EndFor
\EndFor
\State \Return $M[1,n]$
\end{algorithmic}

\subsection{Multiplicación de una cadena de matrices}

No está en el capítulo 6 del libro, pero es el ejemplo más directo del
molde: la última multiplicación parte la cadena en dos
subcadenas, y cada una se resuelve de forma óptima.\visto{Examen jul. 2021 Ej. 2: producto de cadena de matrices}

\begin{problema}{Producto $A_1\times\dots\times A_n$ (examen jul. 2021, Ej. 2)}
$A_i$ es de $d_{i-1}\times d_i$. Con $\OPT(i,j)$ el mínimo de
multiplicaciones escalares para $A_i\times\dots\times A_j$:
\[
\OPT(i,i)=0,\qquad \OPT(i,j)=\min_{i\le k<j}\bigl(\OPT(i,k)+\OPT(k+1,j)+d_{i-1}d_kd_j\bigr),\quad i<j.
\]
Calcular $A_i..A_j$ es calcular $U=A_i..A_k$, $V=A_{k+1}..A_j$ (de forma
óptima, por definición de $\OPT$) y $U\times V$, que cuesta
$d_{i-1}d_kd_j$. Se llena con $\ell=1..n-1$ e $i=1..n-\ell$, $j=i+\ell$;
respuesta $\OPT(1,n)$, orden $\Oh(n^3)$.
\end{problema}

Con la misma estructura se tomó en 2011 ``intercalar paréntesis en una suma
para minimizar la suma de las sumas intermedias'': el término de combinación
es $\sum_{t=i}^{j}x_t$ (que se obtiene en $\Oh(1)$ con sumas
acumuladas) en lugar de $d_{i-1}d_kd_j$.\visto{Examen dic. 2011 Prob. 2: paréntesis con sumas intermedias}

\subsection{Palíndromos}

\begin{problema}{Subsecuencia palíndroma más larga (examen jul. 2022, Ej. 2)}
Con $\OPT(i,j)$ el largo máximo de una subsecuencia de $x_i..x_j$ que es
palíndromo:
\begin{gather*}
\OPT(i,i)=1,\qquad \OPT(i,j)=0\quad (j<i),\\
\OPT(i,j)=\begin{cases}2+\OPT(i+1,j-1) & x_i=x_j\\ \max\{\OPT(i+1,j),\OPT(i,j-1)\} & x_i\ne x_j\end{cases}\quad (i<j).
\end{gather*}
Si $x_i=x_j$, hay un palíndromo óptimo que usa ambos extremos (si usa solo
uno, se lo cambia por el otro; si no usa ninguno, se lo puede alargar). Si
$x_i\ne x_j$, no puede usar los dos. Orden de llenado: $i$ desde $n-1$
bajando hasta 1 y $j$ desde $i+1$ subiendo hasta $n$. $\Oh(n^2)$.
\end{problema}

Para \emph{substrings} palíndromos (lo que necesita la segmentación de 2026)
la tabla booleana es $\mathit{esP}(i,i)=\texttt{true}$,
$\mathit{esP}(i,i+1)=(a_i=a_{i+1})$ y $\mathit{esP}(i,j)=(a_i=a_j)\wedge \mathit{esP}(i+1,j-1)$, con el
mismo orden de llenado; se calcula en $\Oh(n^2)$.\visto{Examen jul. 2022 Ej. 2: subsecuencia palíndroma más larga}

\begin{trampa}[intervalos]
\begin{itemize}[leftmargin=1.2em]
  \item Llenar con $i$ creciente: $\OPT(i,j)$ necesita $\OPT(i+1,\cdot)$, que
  todavía no está.
  \item Olvidar los intervalos vacíos ($j<i$) cuando la recurrencia los usa.
  \item Plantear $\OPT(j)$ sobre prefijos cuando la decisión deja un
  pedazo ``del medio''.
\end{itemize}
\end{trampa}

% =====================================================================
\section{Grillas y tableros}
% =====================================================================

Una grilla en la que solo se avanza en ciertas direcciones (hacia arriba,
hacia abajo, a la derecha) es un DAG implícito: cada casilla es un
subproblema y depende de sus dos o tres vecinas ``anteriores''. La tabla es
la grilla misma y se llena fila por fila en el sentido del movimiento. Los
bordes se manejan con centinelas $\pm\infty$ en columnas ficticias, que es lo
que sugieren las letras.\visto{2P 2023 Ej. 1: sargento en un tablero}

\begin{problema}{El sargento (2P 2023, Ej. 1)}
Desde $(i,j)$ el sargento va a $(i-1,j-1)$, $(i-1,j)$ o $(i-1,j+1)$; salir
de $(i,j)$ cuesta $C(i,j)$. Se va de $(n,1)$ a alguna casilla de la fila 1.
Con $\OPT(i,j)$ el costo mínimo de llegar a $(i,j)$:
\[
\OPT(n,1)=0,\quad \OPT(n,j)=\infty\ (j>1),\quad \OPT(i,0)=\OPT(i,n+1)=\infty,
\]
\[
\OPT(i,j)=\min_{\ell\in\{j-1,j,j+1\}}\bigl(\OPT(i+1,\ell)+C(i+1,\ell)\bigr),\quad 1\le i<n,\ 1\le j\le n.
\]
Se llena con $i$ desde $n-1$ bajando hasta 1; la respuesta es
$\min_k\OPT(1,k)$. Para reconstruir, se parte de la columna que da el mínimo
en la fila 1 y en cada fila se baja a la vecina cuyo término coincide con el
valor actual. $\Oh(n^2)$.
\end{problema}

El minero de 2022 es el espejo: la máquina perfora hacia abajo desde una
columna inicial a elección, maximizando el valor extraído, con
$v_{i,0}=v_{i,n+1}=-\infty$ como centinelas.\visto{Examen feb. 2022 Ej. 3: perforación de una veta}

\begin{problema}{Perforación minera (examen feb. 2022, Ej. 3)}
Con $\OPT(i,j)$ la ganancia máxima perforando desde $(i,j)$ (incluida) hasta
la profundidad $k$:
\begin{gather*}
\OPT(k,j)=v_{k,j},\\
\OPT(i,j)=v_{i,j}+\max\{\OPT(i+1,j-1),\OPT(i+1,j),\OPT(i+1,j+1)\},\quad 1\le i<k.
\end{gather*}
Respuesta $\max_j\OPT(1,j)$; el plan es la columna del $\pdargmax$ y, en cada
fila, la dirección $I$, $V$ o $D$ cuyo término da el valor de la celda.
$\Oh(kn)$.
\end{problema}

\begin{problema}{Contar caminos y recolectar en una grilla (jul. 2017; jul. 2018)}
\textbf{Contar.} Caminos de $(i,j)$ a $(m,n)$ moviéndose a la derecha o
hacia arriba: $\mathit{NC}(i,j)=1$ si $i=m$ o $j=n$, y
$\mathit{NC}(i,j)=\mathit{NC}(i+1,j)+\mathit{NC}(i,j+1)$ en otro caso. Los caminos se parten en
disjuntos según el primer paso, por eso se \emph{suma}. Para $m=6$, $n=3$
da $\binom{9}{3}=84$.

\textbf{Recolectar.} Con $E(i,j)$ objetos en la esquina $(i,j)$ y calles
flechadas hacia índices crecientes: $C(i,j)=0$ si $i=0$ o $j=0$, y
$C(i,j)=\max\{C(i-1,j),C(i,j-1)\}+E(i,j)$. Si hay esquinas intransitables
($E=-1$), se define la accesibilidad $A(i,j)$ (la esquina es transitable y
alguna de las dos anteriores es accesible) y se toma $C(i,j)=-\infty$ en las
no accesibles.
\end{problema}

Estos dos son del programa viejo, pero el de contar caminos marca un
patrón que volvió en 2019: cuando la pregunta es ``¿cuántos?'', los casos
se \emph{suman} en lugar de compararse, porque particionan el conjunto
de soluciones.\visto{Examen jul. 2017 Prob. A: caminos en una grilla}

El del recolector agrega una idea reutilizable: cuando algunas celdas están
prohibidas, conviene una recurrencia auxiliar booleana de alcanzabilidad (o
el valor $-\infty$) para no confundir ``no se puede llegar'' con ``se llega
sin recolectar nada''.\visto{Examen jul. 2018 Prob. A: recolector en calles flechadas}

\begin{trampa}[grillas]
Olvidar los bordes (consultar la columna 0 o $n+1$ sin haberlas definido);
devolver $\OPT(1,1)$ cuando el destino es ``cualquier casilla de la fila
superior''; recorrer las filas en el sentido opuesto al movimiento.
\end{trampa}

% =====================================================================
\section{DAGs, conteo y árboles}
% =====================================================================

\subsection{PD sobre un DAG}

En un grafo dirigido acíclico, cualquier recurrencia en la que $f(v)$ se
obtiene combinando los valores $f(w)$ de los sucesores $w$ de $v$ (con un
máximo, un mínimo o una suma) se calcula procesando
los vértices en orden topológico inverso: cuando se llega a $v$, todos sus
sucesores ya están resueltos. Cada arista se mira una vez, así que el tiempo
es $\Oh(n+m)$ con listas de adyacencia. El orden topológico se obtiene con
el algoritmo de los grados de entrada del capítulo 3, que se pidió
reescribir en diciembre de 2018 y en julio de 2026.\visto{Examen feb. 2023 Ej. 3: camino más largo en un DAG}

\begin{problema}{Camino más largo en un DAG (examen feb. 2023, Ej. 3)}
Los vértices $1..n$ ya están en orden topológico. Con $L(v)$ la longitud
(en aristas) del camino más largo que empieza en $v$:
\[
L(v)=\begin{cases}0 & v\text{ no tiene aristas salientes}\\ 1+\max_{(v,w)\in E}L(w) & \text{si no.}\end{cases}
\]
Se calcula con $v$ desde $n$ bajando hasta 1 y la respuesta es $\max_vL(v)$.
Para devolver el camino, guardar $\mathit{sig}[v]$ (el $w$ que dio el máximo) y el
vértice de inicio del mejor, y seguir los punteros. $\Oh(|V|+|E|)$.
\end{problema}

Cuando el orden topológico no viene dado, se calcula dentro del mismo
algoritmo. En julio de 2026 los valores se propagaban hacia adelante mientras
se sacaban vértices de la cola, y había que detectar además si el grafo tenía
ciclos.\visto{Examen jul. 2026 Ej. 2: tiempos de inicio con dependencias}

\begin{problema}{Tiempos de inicio con dependencias (examen jul. 2026, Ej. 2)}
La subtarea $i$ dura $d_i$ y solo puede empezar cuando terminaron todas las
que la preceden. El inicio mínimo es $s_i=0$ si no tiene previas y
$s_i=\max\{s_j+d_j:\ i\in p_j\}$ si no (camino crítico). Se calcula en el
orden en que el algoritmo de grados de entrada saca los vértices de la cola:
al sacar $h$, para cada $i\in p_h$ se hace $s_i\gets\max\{s_i,s_h+d_h\}$, se
decrementa $a_i$ y se encola $i$ si llegó a 0. Si al final queda algún
$a_i>0$ hay un ciclo y el proyecto no es factible. $\Oh(n+m)$.

\textbf{Error típico.} Recorrer $i=1..n$ en orden de índice y propagar
$s_j\gets\max\{s_j,s_i+d_i\}$: puede usar un $s_i$ todavía no definitivo.
Con $p_1=\{\}$, $p_2=\{1\}$, $p_3=\{2\}$ y $d=(5,10,15)$ da $s_1=10$ en vez
de 25.
\end{problema}

\subsection{Contar caminos}

Para contar caminos de $v$ a $t$ en un DAG: $N(t)=1$ y
$N(v)=\sum_{(v,w)\in E}N(w)$. Para contar caminos \emph{de largo mínimo}
desde un origen se restringe la suma a las aristas que respetan las
distancias: en un grafo sin pesos, las que van de un nivel BFS al
siguiente.\visto{Examen jul. 2019 Ej. 2: contar historias más cortas}

\begin{problema}{Historias más cortas de un libro (examen jul. 2019, Ej. 2)}
Los capítulos forman un DAG con a lo sumo 3 aristas por vértice; se pide el
largo mínimo $L$ de una historia (de 1 a un capítulo final) y la cantidad
$N$ de historias de ese largo, en $\Oh(n)$. Con $L(w)$ el nivel BFS de $w$
desde 1, los caminos mínimos a $w$ se parten según el último vértice antes
de $w$, que está en el nivel anterior:
\[
N(1)=1,\qquad N(w)=\sum_{t:\ (t,w)\in E,\ L(w)=L(t)+1}N(t).
\]
Se hace BFS por niveles acumulando $N$ y se corta al completar el primer
nivel con capítulos finales: $L$ es ese nivel y $N$ es la suma de $N(u)$
sobre los finales de ese nivel. El análisis es el de BFS con $|E|=\Oh(n)$.
\end{problema}

\subsection{Árboles}

En un árbol con raíz, los subproblemas son los subárboles y se resuelven en
postorden (hijos antes que el padre). El caso clásico es el conjunto
independiente de peso máximo (capítulo 10 del libro): para cada vértice $u$
se calculan dos valores, el óptimo del subárbol de $u$ que incluye a $u$ y
el que no lo incluye.\visto{Examen feb. 2022 Ej. 1: independiente de peso máximo en árbol}

\begin{problema}{Conjunto independiente de peso máximo en árbol binario (examen feb. 2022, Ej. 1)}
$q(T)$ es el peso máximo de un independiente de $T$ que no incluye la raíz
$v$; $p(T)$ el máximo sin restricción. Si $T$ es una hoja, $p(T)=w(v)$ y
$q(T)=0$. Si no, con subárboles $T_I$ y $T_D$:
\[
q(T)=p(T_I)+p(T_D),\qquad p(T)=\max\bigl\{w(v)+q(T_I)+q(T_D),\ q(T)\bigr\}.
\]
La unión de independientes de $T_I$ y $T_D$ es independiente; si se incluye
$v$ no se pueden incluir las raíces de los subárboles. La función recursiva
devuelve el par $(p,q)$; para un árbol perfecto,
$f(n)\le2f\bigl(\tfrac{n-1}{2}\bigr)+c$, que da $f(n)\le cn$ por inducción.
En general es $\Oh(n)$ porque cada nodo se procesa una vez.
\end{problema}

Sin pesos, en un bosque alcanza con el greedy que toma hojas (diciembre
de 2019, demostrado por intercambio). Con pesos ese greedy falla: en el
camino $a-b-c$ con pesos $1,10,1$ toma $\{a,c\}$ (peso 2) y el óptimo es
$\{b\}$ (peso 10). Por eso con pesos hace falta la PD.

\begin{receta}[DAGs y árboles]
\begin{itemize}[leftmargin=1.2em]
  \item Decí explícitamente en qué orden se procesan los vértices (orden
  topológico inverso, postorden) y por qué así los valores necesarios ya
  están calculados.
  \item Para $\Oh(n+m)$: cada vértice se procesa una vez y cada arista se mira
  una vez desde su origen, con listas de adyacencia.
  \item Si se cuenta, sumá sobre una partición de las soluciones y decí cuál
  es la partición.
\end{itemize}
\end{receta}

% =====================================================================
\section{Caminos mínimos con costos negativos: Bellman--Ford}
% =====================================================================

\subsection{Por qué no alcanza Dijkstra}

Con costos negativos, Dijkstra puede fijar mal una distancia: con aristas
$s\to a$ de costo 2, $s\to b$ de costo 3 y $b\to a$ de costo $-2$, fija
$d(a)=2$ cuando el camino $s,b,a$ cuesta 1. Sumarle una constante a todas las
aristas para hacerlas positivas tampoco sirve, porque penaliza más a los
caminos con más aristas. Y si hay un ciclo de costo negativo desde el que se
alcanza $t$, el costo mínimo no está definido (se puede dar vueltas
indefinidamente).

\begin{resultado}[(6.22)]
Si $G$ no tiene ciclos negativos, existe un camino mínimo de $s$ a $t$ que es
simple y, por lo tanto, tiene a lo sumo $n-1$ aristas.
\end{resultado}

Esto sugiere el parámetro: la \emph{cantidad de aristas permitidas}. En el
libro los caminos van \emph{hacia} $t$: $\OPT(i,v)$ es el costo mínimo de un
camino de $v$ a $t$ con a lo sumo $i$ aristas, con $\OPT(0,t)=0$ y
$\OPT(0,v)=\infty$ para $v\ne t$. La respuesta es $\OPT(n-1,s)$.

\begin{resultado}[(6.23)]
Si $i>0$, entonces
$\OPT(i,v)=\min\Bigl(\OPT(i-1,v),\ \min_{w\in V}\bigl(\OPT(i-1,w)+c_{vw}\bigr)\Bigr)$.
\end{resultado}

O el camino óptimo usa a lo sumo $i-1$ aristas (primer término), o usa
exactamente $i$ y la primera es alguna $(v,w)$; lo que sigue es un camino de
$w$ a $t$ con a lo sumo $i-1$ aristas, que tiene que ser óptimo.

\pdalgo{Bellman--Ford (Shortest-Path$(G,s,t)$)}
\begin{algorithmic}[1]
\State $M[0,t]\gets0$;\ $M[0,v]\gets\infty$ para todo $v\ne t$
\For{$i=1$ \textbf{hasta} $n-1$}
  \ForAll{$v\in V$}
    \State $M[i,v]\gets\min\bigl(M[i-1,v],\ \min_{(v,w)\in E}(M[i-1,w]+c_{vw})\bigr)$
  \EndFor
\EndFor
\State \Return $M[n-1,s]$
\end{algorithmic}

\begin{resultado}[(6.24) y (6.25)]
Si $G$ no tiene ciclos negativos, el algoritmo calcula correctamente el
costo mínimo de un camino $s$-$t$. La tabla tiene $n^2$ celdas y cada una
cuesta $\Oh(n)$, lo que da $\Oh(n^3)$; mirando solo las aristas que salen de
$v$, cada iteración cuesta $\Oh(\sum_vn_v)=\Oh(m)$ y el total es $\Oh(mn)$.
\end{resultado}

En la práctica se usa un solo arreglo $M[v]$ actualizado en el lugar, más un
puntero $\mathit{first}[v]$ al primer vértice del camino actual, lo que baja el
espacio a $\Oh(n)$ y permite cortar en cuanto una iteración no cambia nada.
Si el grafo de punteros $\{(v,\mathit{first}[v])\}$ llega a tener un ciclo, ese ciclo
tiene costo negativo.

\subsection{Las versiones de examen}

En el parcial de 2018 se pidió primero la recurrencia para ``exactamente
$i$ aristas'' y después obtener el óptimo con a lo sumo $K$ aristas como un
mínimo; la letra fijaba los nombres y había que respetarlos.\visto{2P 2018 Ej. 2: caminos con a lo sumo $K$ aristas}

\begin{problema}{Caminos mínimos con a lo sumo $K$ aristas (2P 2018, Ej. 2)}
$\mathit{CMC}(i,v)$ es el costo mínimo de un camino de $v$ a $t$ con
\emph{exactamente} $i$ aristas ($\infty$ si no hay):
\[
\mathit{CMC}(0,v)=\begin{cases}0&v=t\\ \infty&v\ne t\end{cases}\qquad
\mathit{CMC}(i,v)=\min_{(v,w)\in E}\bigl(\mathit{CMC}(i-1,w)+c_{vw}\bigr),\quad 1\le i\le K,
\]
con $\min\emptyset=\infty$, y
$\OPT(K,v)=\min_{0\le i\le K}\mathit{CMC}(i,v)$. La alternativa directa es
$\OPT(k,t)=0$ y $\OPT(k,v)=\min_{(v,w)}(\OPT(k-1,w)+c_{vw})$ para $v\ne t$.
Con listas de adyacencia el ciclo principal es $\Oh(K(n+m))$ (con matriz,
$\Oh(Kn^2)$), y el mínimo final $\Oh(nK)$; como $K<n$, es polinomial.
\end{problema}

En julio de 2023 se pidió lo mismo pero con caminos que salen del origen, y
con el costo de cada arista armado a partir de dos datos de la letra. Los
costos son positivos, así que Dijkstra también serviría; lo que se evaluaba
era la recurrencia en $i$.\visto{Examen jul. 2023 Ej. 2: routers con a lo sumo $i$ saltos}

\begin{problema}{Routers con a lo sumo $i$ saltos (examen jul. 2023, Ej. 2)}
Cada enlace $\{u,v\}$ da dos aristas $(u,v)$ y $(v,u)$; la arista $(u,v)$
cuesta $\mathit{TP}(u)+\mathit{TT}(u,v)$ (procesar en $u$ y transmitir). Aquí los caminos
salen \emph{desde} $s$: $\OPT(0,s)=0$, $\OPT(0,v)=\infty$ para $v\ne s$ y
\begin{multline*}
\OPT(i,v)=\min\Bigl\{\OPT(i-1,v),\\ \min_{(u,v)\in E}\bigl(\OPT(i-1,u)+\mathit{TP}(u)+\mathit{TT}(u,v)\bigr)\Bigr\},\quad 1\le i\le n-1.
\end{multline*}
Respuesta $\OPT(n-1,t)$, en $\Oh(nm)$. Para el camino: desde $v=t$ e $i=n-1$
bajando, si $\OPT(i,v)\ne\OPT(i-1,v)$ se busca el $u$ que dio el mínimo, se
lo agrega al principio y $v\gets u$.
\end{problema}

\subsection{Detección de ciclos negativos}

Para saber si $G$ tiene \emph{algún} ciclo negativo (no solo uno que llegue a
$t$) se agrega un vértice nuevo $t$ y una arista $(v,t)$ de costo 0 desde
cada $v$: todo ciclo de $G$ queda conectado a $t$, y $G$ tiene un ciclo
negativo si y solo si el grafo aumentado tiene uno desde el que se llega a
$t$. Después se corre Bellman--Ford una iteración más de lo necesario.

\begin{resultado}[(6.30) y (6.31)]
Si $\OPT(n,v)=\OPT(n-1,v)$ para todo $v$, no hay ciclos negativos que
lleguen a $t$. Si $\OPT(n,v)<\OPT(n-1,v)$ para algún $v$, un camino de $v$ a
$t$ de costo $\OPT(n,v)$ tiene un ciclo, y ese ciclo tiene costo negativo.
El ciclo se encuentra siguiendo ese camino, y todo se hace en $\Oh(mn)$.
\end{resultado}

La aplicación típica es el arbitraje: con tasas de cambio $r_{ij}$ entre
monedas, hay un ciclo de cambios con ganancia ($\prod r_{ij}>1$) si y solo si
hay un ciclo negativo con costos $c_{ij}=-\log r_{ij}$. En los exámenes
recientes no se tomó la detección de ciclos, pero es teoría del capítulo
que se puede preguntar como verdadero o falso.

\subsection{Floyd--Warshall (programa anterior)}

Hasta 2017 el curso daba Floyd para caminos mínimos entre todo par. No está
en el capítulo 6 del libro, pero los exámenes viejos lo preguntan seguido y
la idea es la misma: el parámetro extra limita los vértices intermedios
permitidos.\visto{Examen jul. 2015 Ej. 4: optimalidad y recurrencia de Floyd}

\begin{resultado}[Floyd]
Con $f_k(i,j)$ el costo mínimo de un camino de $i$ a $j$ cuyos vértices
intermedios están en $\{1..k\}$: $f_0(i,j)=W(i,j)$ y
$f_k(i,j)=\min\{f_{k-1}(i,j),\ f_{k-1}(i,k)+f_{k-1}(k,j)\}$. Se puede usar
una sola matriz; el tiempo es $\Theta(n^3)$. Funciona con costos negativos si
no hay ciclos negativos, y un valor negativo en la diagonal indica uno.
\end{resultado}

El principio de optimalidad que se pedía enunciar es que todo subcamino de
un camino mínimo es mínimo (si no, se lo reemplaza). Por eso el camino
mínimo de $v_1$ a $v_k$ que visita $v_2..v_{k-1}$ en orden es la
concatenación de los caminos mínimos entre consecutivos.\visto{2P 2009 Ej. 3: camino que visita una secuencia}

La variante de carga máxima de 2011 (el ``ancho'' de un camino es su arista
de menor capacidad y se quiere el camino más ancho) cambia las operaciones:
\[
f_k(i,j)=\max\bigl\{f_{k-1}(i,j),\ \min\{f_{k-1}(i,k),f_{k-1}(k,j)\}\bigr\}.
\]
Para reconstruir se guarda en $\mathit{Camino}[i][j]$ el último $k$ que mejoró la
celda y se arma recursivamente $i\leadsto k\leadsto j$.\visto{2P 2011 Ej. 2: carga máxima entre ciudades}

\begin{trampa}[caminos mínimos]
\begin{itemize}[leftmargin=1.2em]
  \item Usar Dijkstra (o sumar una constante) con costos negativos.
  \item Olvidar el término $\OPT(i-1,v)$ en la versión ``a lo sumo $i$
  aristas'' (sin él, es ``exactamente $i$'').
  \item Confundir el sentido: en el libro $\OPT(i,v)$ es de $v$ \emph{hacia}
  $t$ y se minimiza sobre aristas que \emph{salen} de $v$; si se define desde
  $s$, se minimiza sobre aristas que \emph{entran} a $v$.
  \item Analizar como $\Oh(n^3)$ cuando se pide $\Oh(mn)$: hay que decir que
  cada arista se mira una vez por iteración.
  \item Afirmar que Bellman--Ford ``funciona con ciclos negativos'': los
  detecta, pero con ellos no hay camino mínimo.
\end{itemize}
\end{trampa}

% =====================================================================
\section{Mapa de lo que se preguntó}
% =====================================================================

Todos los ejercicios de programación dinámica de los parciales y exámenes
disponibles, por año. Hasta 2017 el curso seguía otro programa (con
backtracking y Floyd, y el pedido típico era ``dé una fórmula recursiva
$f$''); desde 2018 sigue el libro de Kleinberg y Tardos y el pedido es
recurrencia, algoritmo iterativo, reconstrucción y orden.

\newcommand{\pdmapa}[1]{%
\par\smallskip{\small
\begin{tabularx}{\textwidth}{@{}>{\raggedright\arraybackslash}p{2.9cm}>{\raggedright\arraybackslash}p{2.5cm}>{\raggedright\arraybackslash}X@{}}
\toprule
\textbf{Fuente} & \textbf{Patrón} & \textbf{Qué se pedía}\\
\midrule
#1
\bottomrule
\end{tabularx}}\par\smallskip}

\pdmapa{
2P 2008 Ej.~1 & Mochila 0/1 & Fórmula recursiva para animales en un barco de capacidad $G$.\\
Ex. feb. 2009 Ej.~3 & Capacidad (asignación) & Interés óptimo invirtiendo $d$ pesos en $n$ empresas; versión iterativa y cómo reconstruir.\\
Ex. jul. 2009 Ej.~2 & Dos restricciones & Menús de mínimo valor calórico con precio $\le M$ y proteínas $\ge W$.\\
2P 2009 Ej.~3 & Caminos mínimos & Camino mínimo de $v_1$ a $v_k$ que visita la secuencia en orden.\\
Ex. dic. 2009 Prob.~1 & Capacidad con stock & Golosinas: gastar lo más posible con la menor cantidad.\\
Ex. jul. 2010 Prob.~2 & Dos restricciones & Ringtones de tamaño mínimo con presupuesto y ranking mínimo.\\
2P 2010 Ej.~2 & Prefijo con estado & Dos integrantes marcan lugares en orden; $f(l,k)$.\\
Ex. dic. 2010 Prob.~1 & Mochila 0/1 & Comerciante romano: contraejemplo greedy y recurrencia.\\
2P 2011 Ej.~2 & Floyd (máx-mín) & Carga máxima entre todo par de ciudades; implementación y camino.\\
Ex. dic. 2011 Prob.~2 & Intervalos & Paréntesis que minimizan la suma de sumas intermedias.\\
Ex. feb. 2012 Ej.~2 & Mochila 0/1 & Publicidad con presupuesto: greedy falla; recurrencia del rating.\\
Ex. jul. 2012 Ej.~2 & Floyd & Qué resuelve, qué técnica usa y aplicación paso a paso.\\
2P 2012 Ej.~1 & Capacidad exacta con stock & Carga exacta $C$ con mínimas unidades; $f(p,c)$, iterativo y orden.\\
Ex. feb. 2013 Ej.~2 & Teoría & Enunciar el principio de optimalidad y aplicarlo a caminos mínimos.\\
Ex. jul. 2013 Ej.~1 & Memoización & Verdadero o falso: Fibonacci se puede calcular en $\Oh(n)$.\\
Ex. jul. 2013 Prob.~1 & Dos secuencias & Alineación óptima de cubos de colores con comodines.\\
}

\pdmapa{
Ex. jul. 2014 Ej.~3 & Dos restricciones & Contenedor: máximo de compras con capacidad $C$ y valor $\le M$.\\
Ex. jul. 2015 Ej.~4 & Floyd & Por qué vale el principio de optimalidad; recurrencia de Floyd.\\
2P 2015 Ej.~1 & Prefijo, $p(i)=i-d$ & Edificios a distancia $\ge d$; greedy falla; recurrencia e implementación.\\
Ex. dic. 2015 Ej.~2 & Mochila 0/1 & Definir la variante, demostrar optimalidad, plantear la recurrencia.\\
Ex. feb. 2016 Ej.~4 & Dos restricciones & Programadores: experiencia máxima con presupuesto y mínimo $E$.\\
Ex. jul. 2016 Ej.~4 & Teoría & Qué resuelven Floyd, suma de subconjuntos y otros, y con qué técnica.\\
Ex. jul. 2016 Prob.~C & Mochila 0/1 & Técnica, implementación en C* y versión fraccionaria.\\
2P 2016 Ej.~2 & Último elemento & Largos $L_i$ de subsecuencias crecientes; versión recursiva y PD.\\
Ex. feb. 2017 Preg.~3 & Memoización & Fibonacci con DyC y con PD en espacio $\Oh(1)$; contar sumas.\\
Ex. jul. 2017 Prob.~A & Grilla (conteo) & Cantidad de caminos monótonos $\mathit{NC}(i,j)$; caso $m=6$, $n=3$.\\
2P 2017 Ej.~1 & Dos secuencias & Descifrar el código del substring común más largo; memoria.\\
Ex. dic. 2017 Prob.~B & Prefijo binario & Baldosas no contiguas de máximo valor; devolver las elegidas.\\
Ex. feb. 2018 Preg.~2 & Capacidad ilimitada & Greedy de monedas: qué técnica es y cuándo falla.\\
Ex. feb. 2018 Preg.~4 & Teoría & Qué resuelve alineamiento de secuencias y con qué técnica.\\
Ex. feb. 2018 Prob.~B & Último elemento & LIS con $G_i$ y $L_i$; pila de bandejas.\\
Ex. jul. 2018 Prob.~A & Grilla & Recolector con calles flechadas; esquinas intransitables.\\
}

\pdmapa{
2P 2018 Ej.~2 & Cantidad de aristas & Costo mínimo con $\le K$ aristas vía $\mathit{CMC}(i,v)$; algoritmo y orden.\\
Ex. dic. 2018 Ej.~3 & Dos secuencias & Ciclos mínimos para ejecutar dos programas en paralelo.\\
Ex. feb. 2019 Ej.~2 & Dos secuencias & Distancia de edición; lema de las últimas posiciones.\\
Ex. jul. 2019 Ej.~2 & DAG (conteo) & Largo mínimo de las historias y cuántas hay, en $\Oh(n)$.\\
2P 2019 Ej.~2 & Prefijo binario & Ladrón de casas no contiguas; valor y conjunto óptimo.\\
Ex. dic. 2019 Ej.~2 & Prefijo con estado & Neumáticos A/B por vuelta con costo de parada; $\OPT(i,X)$.\\
Ex. feb. 2020 Ej.~3 & Sufijo booleano & Dinosaurio que avanza 1 o salta 4: ¿hay secuencia ganadora?\\
Ex. jul. 2020 Ej.~3 & Capacidad ilimitada & Varillas con la mínima cantidad de soldaduras; $C(i,\ell)$.\\
2P 2020 Ej.~2 & Segmentación & Fragmentos con costo máximo menos mínimo y costo por corte.\\
Ex. dic. 2020 Ej.~3 & Capacidad ilimitada & Cambio con menos billetes; instancia donde greedy falla.\\
Ex. jul. 2021 Ej.~2 & Intervalos & Producto de una cadena de matrices.\\
2P 2021 Ej.~1 & Sufijo con salto & Tarifa por consumo o plana de 7 días; planificación óptima.\\
}

\pdmapa{
Ex. feb. 2022 Ej.~1 & Árbol & Independiente de peso máximo en árbol binario ($p$ y $q$); orden $\Oh(n)$.\\
Ex. feb. 2022 Ej.~3 & Grilla & Perforación minera de ganancia máxima; plan de perforación.\\
Ex. jul. 2022 Ej.~2 & Intervalos & Subsecuencia palíndroma más larga.\\
2P 2022 Ej.~1 & Prefijo con estado numérico & Paradas de F1 con desgaste de neumáticos; $\OPT(i,j)$.\\
Ex. feb. 2023 Ej.~3 & DAG & Camino más largo con orden topológico dado, en $\Oh(|V|+|E|)$.\\
Ex. jul. 2023 Ej.~2 & Cantidad de aristas & Routers: tiempo mínimo con a lo sumo $i$ saltos; camino.\\
2P 2023 Ej.~1 & Grilla & Sargento: costo mínimo hasta la fila superior; movimientos.\\
Ex. dic. 2023 Ej.~2 & Último elemento & Subsecuencia creciente más larga y reconstrucción.\\
Ex. feb. 2024 Ej.~2 & Segmentación & Concatenación de prefijos con máximas coincidencias.\\
2P 2024 Ej.~1 & Prefijo con estado & Cuchillas A y B con costo de cambio; solución óptima.\\
Ex. dic. 2024 Ej.~2 & Sufijo con salto & Ravioles y sorrentinos; planificación.\\
2P 2025 Ej.~3 & Sufijo con salto & Pedidos de soporte de uno o dos días.\\
Ex. feb. 2026 Ej.~3 & Segmentación & Segmentación mínima en palíndromos; posiciones de inicio.\\
Ex. jul. 2026 Ej.~2 & DAG & Tiempos de inicio con dependencias (camino crítico) en $\Oh(n+m)$.\\
}

% =====================================================================
%  Parte 3 — Flujo en redes (Kleinberg–Tardos, cap. 7)
%  Fragmento: se incluye con \input desde el documento maestro.
% =====================================================================

% ---------- macros y estilos locales (prefijo fl) ----------
\newcommand{\flout}{f^{\mathrm{out}}}
\newcommand{\flin}{f^{\mathrm{in}}}
\newcommand{\flGf}{G_f}
\newcommand{\flcuello}{\operatorname{cuello}}
\tikzset{
  flv/.style={circle,draw=tinta,fill=white,inner sep=1pt,minimum size=5.6mm,font=\small},
  fla/.style={-{Stealth[length=2.1mm]},draw=gris,semithick},
  flp/.style={-{Stealth[length=2.3mm]},draw=flujo,line width=1.3pt},
  flr/.style={-{Stealth[length=2.1mm]},draw=gris,semithick,dashed},
  flrp/.style={-{Stealth[length=2.3mm]},draw=flujo,line width=1.3pt,dashed},
  fll/.style={font=\scriptsize,fill=white,inner sep=0.8pt,text=tinta},
  flt/.style={font=\footnotesize\sffamily,text=gris},
}
% nodos del ejemplo s, a, b, t (Examen jul. 2023)
\newcommand{\flnodosej}{%
  \node[flv] (s) at (0,0) {$s$};
  \node[flv] (a) at (1.55,0.95) {$a$};
  \node[flv] (b) at (1.55,-0.95) {$b$};
  \node[flv] (t) at (3.1,0) {$t$};}

\parte{flujo}{Flujo en redes}
  {\max_{f}\, v(f) \;=\; \min_{(A,B)}\, c(A,B)}
  {Mandamos la mayor cantidad posible de algo de $s$ a $t$ sin pasarnos de las capacidades. Lo que frena al flujo es siempre un corte, y casi todos los ejercicios de parcial consisten en armar la red adecuada y leer la respuesta en el flujo máximo.}

% =====================================================================
\section{Redes, flujos y cortes}
% =====================================================================

Este capítulo del libro (KT cap.~7) entra en los parciales de dos maneras. Por un lado están las preguntas de teórico: dar definiciones, demostrar $v(f)\le c(A,B)$ o la proposición (7.9), explicar la cota $\Oh(mC)$ o correr Ford–Fulkerson a mano. Por otro, los ejercicios de modelado, que son la mayoría de los puntos: se arma una red cuyo flujo máximo resuelve un problema de asignación y se prueba la equivalencia en los dos sentidos.\visto{2P 2022 Ej.~2: definiciones y $v(f)\le c(A,B)$}

\begin{definicion}[red de flujo]
Una \emph{red de flujo} es un grafo dirigido $G=(V,E)$ con un nodo \emph{fuente} $s$ sin aristas entrantes, un nodo \emph{sumidero} $t$ sin aristas salientes y una \emph{capacidad} entera $c_e\ge 0$ en cada arista. Se supone además que todo nodo tiene alguna arista incidente, así que $m\ge n/2$.
\end{definicion}

\begin{definicion}[flujo $s$-$t$ y valor]
Un \emph{flujo $s$-$t$} es una función $f:E\to\mathbb{R}_{\ge 0}$ que cumple
\begin{enumerate}[label=(\roman*)]
  \item \emph{capacidad}: $0\le f(e)\le c_e$ para toda $e\in E$;
  \item \emph{conservación}: para todo $v\ne s,t$, $\displaystyle\sum_{e\text{ entra a }v} f(e)=\sum_{e\text{ sale de }v} f(e)$, es decir, $\flin(v)=\flout(v)$.
\end{enumerate}
El \emph{valor} del flujo es lo que sale de la fuente: $v(f)=\flout(s)$.
\end{definicion}

En el examen de diciembre de 2025 pidieron exactamente esta definición, y en el de julio de 2020 preguntaron qué cotas tiene que cumplir $f(e)$. Las dos condiciones tienen nombre, y conviene escribirlo: en las correcciones se pide nombrar ``capacidad'' y ``conservación'' cada vez que se usan.\visto{Examen dic.~2025 Ej.~1(d): definir flujo}

\begin{definicion}[corte y capacidad]
Un \emph{corte $s$-$t$} es una partición $(A,B)$ de $V$ con $s\in A$ y $t\in B$. Su \emph{capacidad} es la suma de las capacidades de las aristas que \emph{salen} de $A$:
\[ c(A,B)=\sum_{e\text{ sale de }A} c_e . \]
Las aristas que entran a $A$ no cuentan.
\end{definicion}

Para un conjunto $S$ de nodos escribimos $\flout(S)$ para la suma del flujo en las aristas que salen de $S$ y $\flin(S)$ para las que entran. La definición de corte y de capacidad se pidió tal cual en el segundo parcial de 2025.\visto{2P 2025 Ej.~1(b): corte y capacidad}

\begin{resultado}[(7.6) y (7.7)]
Sean $f$ un flujo $s$-$t$ y $(A,B)$ un corte $s$-$t$ cualesquiera. Entonces
\[ v(f)=\flout(A)-\flin(A) \qquad\text{y también}\qquad v(f)=\flin(B)-\flout(B). \]
\end{resultado}

\begin{proof}[Idea]
Por definición $v(f)=\flout(s)-\flin(s)$, porque $\flin(s)=0$. Para $v\in A$, $v\neq s$, la conservación da $\flout(v)-\flin(v)=0$. Sumando sobre todos los nodos de $A$:
$v(f)=\sum_{v\in A}\bigl(\flout(v)-\flin(v)\bigr)$. Una arista con los dos extremos en $A$ aparece una vez con $+$ y otra con $-$, así que se cancela; una que sale de $A$ aparece solo con $+$ y una que entra solo con $-$. Lo que queda es $\flout(A)-\flin(A)$.
\end{proof}

\begin{resultado}[(7.8)]
Sean $f$ un flujo $s$-$t$ y $(A,B)$ un corte $s$-$t$ cualesquiera. Entonces $v(f)\le c(A,B)$.
\end{resultado}

Esta es la demostración que pidieron en el 2P de 2022 y en el examen de julio de 2020, en ambos casos con (7.6) como dato. La corrección espera que se justifique cada desigualdad con la propiedad que se usa.\visto{Examen jul.~2020 Ej.~2: cota y $v(f)\le c(A,B)$}

\begin{proof}
\begin{align*}
v(f) &= \flout(A)-\flin(A) && \text{por (7.6)}\\
     &\le \flout(A) && \text{porque } \flin(A)\ge 0 \text{ (los flujos son no negativos)}\\
     &= \sum_{e\text{ sale de }A} f(e) && \text{definición de } \flout(A)\\
     &\le \sum_{e\text{ sale de }A} c_e = c(A,B) && \text{condición de capacidad.}
\end{align*}
La primera desigualdad es igualdad si y solo si ninguna arista que entra a $A$ lleva flujo. La segunda lo es si y solo si todas las que salen de $A$ están saturadas.
\end{proof}

\begin{trampa}
\begin{itemize}
  \item Contar en $c(A,B)$ las aristas que entran a $A$. Solo cuentan las salientes, y esa asimetría es justamente lo que hace funcionar la demostración de (7.8).
  \item Escribir la conservación también para $s$ y $t$. Vale solo para $v\ne s,t$.
  \item Demostrar (7.8) sin decir en qué paso se usa la no negatividad y en cuál la capacidad.
\end{itemize}
\end{trampa}

% =====================================================================
\section{Grafo residual y algoritmo de Ford–Fulkerson}
% =====================================================================

\subsection{Grafo residual}

Dado un flujo $f$ en $G$, el \emph{grafo residual} $\flGf$ tiene los mismos nodos que $G$ y, por cada arista $e=(u,v)$ de $G$:
\begin{itemize}
  \item si $f(e)<c_e$, una \emph{arista hacia adelante} $(u,v)$ con capacidad residual $c_e-f(e)$: lo que todavía podemos empujar;
  \item si $f(e)>0$, una \emph{arista hacia atrás} $(v,u)$ con capacidad residual $f(e)$: lo que podemos ``deshacer'' devolviendo flujo.
\end{itemize}
Cada arista de $G$ genera a lo sumo dos aristas en $\flGf$, así que $\flGf$ tiene a lo sumo $2m$ aristas.

La definición del grafo residual se pidió como parte (a) del examen de diciembre de 2021, y con esa misma redacción (adelante/atrás, sobrante/deshacer) aparece en la solución oficial. Las aristas con capacidad residual $0$ no se dibujan.\visto{Examen dic.~2021 Ej.~3(a): definir grafo residual}

\subsection{Aumentar a lo largo de un camino}

Sea $P$ un camino simple de $s$ a $t$ en $\flGf$ (un \emph{camino de aumento}). Su \emph{cuello de botella}, $\flcuello(P,f)$, es la menor capacidad residual de las aristas de $P$.

\begin{algorithmic}[1]
\Function{Aumentar}{$f,P$}
  \State $b\gets \flcuello(P,f)$
  \ForAll{arista $(u,v)\in P$}
    \If{$(u,v)$ es arista hacia adelante}
      \State $f(u,v)\gets f(u,v)+b$ \Comment{empujamos más}
    \Else
      \State $f(v,u)\gets f(v,u)-b$ \Comment{devolvemos flujo de $(v,u)\in E$}
    \EndIf
  \EndFor
  \State \Return $f$
\EndFunction
\end{algorithmic}

\begin{resultado}[(7.1) y (7.3)]
Si $P$ es un camino simple $s$-$t$ en $\flGf$, entonces $f'=\textsc{Aumentar}(f,P)$ es un flujo en $G$ y
$v(f')=v(f)+\flcuello(P,f)>v(f)$.
\end{resultado}

Que $f'$ sea flujo se ve en dos partes. La capacidad se cumple porque $b$ no supera la capacidad residual de ninguna arista del camino. La conservación vale porque cada nodo interno de $P$ tiene una arista entrante y una saliente en $P$, y en los cuatro casos posibles (adelante/adelante, adelante/atrás, etc.) el cambio neto en ese nodo es $0$. El valor crece porque la primera arista de $P$ sale de $s$ y tiene que ser hacia adelante, ya que $s$ no tiene aristas entrantes en $G$.

\subsection{El algoritmo}

\begin{algorithmic}[1]
\Function{FordFulkerson}{$G,s,t,c$}
  \State $f(e)\gets 0$ para toda $e\in E$
  \While{existe un camino simple $P$ de $s$ a $t$ en $\flGf$}
    \State $f\gets$ \Call{Aumentar}{$f,P$}
    \State actualizar $\flGf$
  \EndWhile
  \State \Return $f$
\EndFunction
\end{algorithmic}

\begin{resultado}[(7.2) Integralidad durante la ejecución]
Si todas las capacidades son enteras, en toda etapa del algoritmo los valores $f(e)$ y las capacidades residuales de $\flGf$ son enteros.
\end{resultado}
Se prueba por inducción en las iteraciones. Al principio $f\equiv 0$. Si en una etapa todo es entero, el cuello de botella es entero y \textsc{Aumentar} suma o resta un entero.

\begin{resultado}[(7.4) y (7.5)]
Sea $C=\sum_{e\text{ sale de }s}c_e$. Con capacidades enteras, Ford–Fulkerson termina en a lo sumo $C$ iteraciones del \textbf{mientras}, y puede implementarse en tiempo $\Oh(mC)$.
\end{resultado}

En el examen de febrero de 2026 preguntaron qué cota es esa y qué significa cada símbolo. En el de julio de 2020 pidieron justificar que el algoritmo corre en $\Oh(mC)$ si el flujo máximo está acotado por $C$.\visto{Examen feb.~2026 Ej.~1(e): explicar $\Oh(mC)$}

\begin{proof}[Por qué $\Oh(mC)$]
\emph{$C$ acota las iteraciones.} Por (7.2) y (7.3) cada iteración aumenta $v(f)$ en al menos una unidad, y por (7.8) $v(f)\le c(\{s\},V\setminus\{s\})=C$. Por lo tanto hay a lo sumo $C$ iteraciones. Sirve también la capacidad de \emph{cualquier} otro corte, y en los ejercicios conviene elegir el de menor capacidad que se pueda acotar.

\emph{$m$ acota el costo de cada iteración.} Construir $\flGf$ cuesta $\Oh(m+n)$, porque tiene $n$ nodos y a lo sumo $2m$ aristas. Buscar el camino con BFS o DFS cuesta $\Oh(m+n)$ y aumentar cuesta $\Oh(n)$. Como $m\ge n/2$, todo esto es $\Oh(m)$.
\end{proof}

\begin{trampa}
\begin{itemize}
  \item $\Oh(mC)$ es \emph{pseudopolinomial}: si $C$ depende de números de la entrada (cupos, dosis, capacidades), la cota no es polinomial en el tamaño. En los ejercicios donde se pide ``polinomial en $|A|$ y $|E|$'' hay que encontrar un corte cuya capacidad dependa solo de \emph{cantidades}, típicamente uno formado por aristas de capacidad $1$.
  \item Decir que el algoritmo termina ``porque el flujo aumenta''. Hace falta que aumente al menos $1$, y eso es integralidad.
  \item Con capacidades reales (irracionales), Ford–Fulkerson puede no terminar. La cota es solo para capacidades enteras.
\end{itemize}
\end{trampa}

% =====================================================================
\section{Ford–Fulkerson a mano}
% =====================================================================

Se pide seguido: ``ejecute paso a paso, dibujando la red $G$ y el grafo residual $\flGf$''. En cada paso hay que mostrar el par $(G\text{ con } f,\ \flGf)$, marcar el camino de aumento elegido con su cuello de botella y, al final, justificar que el flujo es máximo con un corte.\visto{Examen jul.~2023 Ej.~3(a): FF paso a paso}

\subsection{Ejemplo completo}

Red del examen de julio de 2023: $c(s,a)=2$, $c(s,b)=4$, $c(a,b)=3$, $c(a,t)=1$, $c(b,t)=5$. A la izquierda va $G$ con la notación $f(e)/c_e$; a la derecha $\flGf$, con las aristas hacia atrás punteadas y el camino de aumento en verde.

\begin{center}
\begin{tikzpicture}[scale=0.95]
  \flnodosej
  \draw[fla] (s)--node[fll,above left]{0/2}(a);
  \draw[fla] (s)--node[fll,below left]{0/4}(b);
  \draw[fla] (a)--node[fll,right]{0/3}(b);
  \draw[fla] (a)--node[fll,above right]{0/1}(t);
  \draw[fla] (b)--node[fll,below right]{0/5}(t);
  \node[flt] at (1.55,-1.75) {$G$, $v(f)=0$};
\end{tikzpicture}\hspace{1.3cm}%
\begin{tikzpicture}[scale=0.95]
  \flnodosej
  \draw[fla] (s)--node[fll,above left]{2}(a);
  \draw[flp] (s)--node[fll,below left]{4}(b);
  \draw[fla] (a)--node[fll,right]{3}(b);
  \draw[fla] (a)--node[fll,above right]{1}(t);
  \draw[flp] (b)--node[fll,below right]{5}(t);
  \node[flt] at (1.55,-1.75) {$\flGf$: $P_1=s\,b\,t$, cuello $4$};
\end{tikzpicture}
\end{center}

\begin{center}
\begin{tikzpicture}[scale=0.95]
  \flnodosej
  \draw[fla] (s)--node[fll,above left]{0/2}(a);
  \draw[flp] (s)--node[fll,below left]{4/4}(b);
  \draw[fla] (a)--node[fll,right]{0/3}(b);
  \draw[fla] (a)--node[fll,above right]{0/1}(t);
  \draw[flp] (b)--node[fll,below right]{4/5}(t);
  \node[flt] at (1.55,-1.75) {$G$, $v(f)=4$};
\end{tikzpicture}\hspace{1.3cm}%
\begin{tikzpicture}[scale=0.95]
  \flnodosej
  \draw[flp] (s)--node[fll,above left]{2}(a);
  \draw[flr] (b)--node[fll,below left]{4}(s);
  \draw[fla] (a)--node[fll,right]{3}(b);
  \draw[flp] (a)--node[fll,above right]{1}(t);
  \draw[fla] (b) to[bend left=25] node[fll]{1}(t);
  \draw[flr] (t) to[bend left=25] node[fll]{4}(b);
  \node[flt] at (1.55,-1.75) {$\flGf$: $P_2=s\,a\,t$, cuello $1$};
\end{tikzpicture}
\end{center}

\begin{center}
\begin{tikzpicture}[scale=0.95]
  \flnodosej
  \draw[flp] (s)--node[fll,above left]{1/2}(a);
  \draw[flp] (s)--node[fll,below left]{4/4}(b);
  \draw[fla] (a)--node[fll,right]{0/3}(b);
  \draw[flp] (a)--node[fll,above right]{1/1}(t);
  \draw[flp] (b)--node[fll,below right]{4/5}(t);
  \node[flt] at (1.55,-1.75) {$G$, $v(f)=5$};
\end{tikzpicture}\hspace{1.3cm}%
\begin{tikzpicture}[scale=0.95]
  \flnodosej
  \draw[flp] (s) to[bend left=25] node[fll]{1}(a);
  \draw[flr] (a) to[bend left=25] node[fll]{1}(s);
  \draw[flr] (b)--node[fll,below left]{4}(s);
  \draw[flp] (a)--node[fll,right]{3}(b);
  \draw[flr] (t)--node[fll,above right]{1}(a);
  \draw[flp] (b) to[bend left=25] node[fll]{1}(t);
  \draw[flr] (t) to[bend left=25] node[fll]{4}(b);
  \node[flt] at (1.55,-1.75) {$\flGf$: $P_3=s\,a\,b\,t$, cuello $1$};
\end{tikzpicture}
\end{center}

\begin{center}
\begin{tikzpicture}[scale=0.95]
  \flnodosej
  \draw[flp] (s)--node[fll,above left]{2/2}(a);
  \draw[flp] (s)--node[fll,below left]{4/4}(b);
  \draw[flp] (a)--node[fll,right]{1/3}(b);
  \draw[flp] (a)--node[fll,above right]{1/1}(t);
  \draw[flp] (b)--node[fll,below right]{5/5}(t);
  \node[flt] at (1.55,-1.75) {$G$, $v(f)=6$};
\end{tikzpicture}\hspace{1.3cm}%
\begin{tikzpicture}[scale=0.95]
  \flnodosej
  \draw[flr] (a)--node[fll,above left]{2}(s);
  \draw[flr] (b)--node[fll,below left]{4}(s);
  \draw[fla] (a) to[bend left=25] node[fll]{2}(b);
  \draw[flr] (b) to[bend left=25] node[fll]{1}(a);
  \draw[flr] (t)--node[fll,above right]{1}(a);
  \draw[flr] (t)--node[fll,below right]{5}(b);
  \node[flt] at (1.55,-1.75) {$\flGf$: no hay camino $s$-$t$};
\end{tikzpicture}
\end{center}

En el último residual $s$ no tiene aristas salientes, así que el conjunto de nodos alcanzables desde $s$ es $A^*=\{s\}$. El corte $(\{s\},\{a,b,t\})$ tiene capacidad $2+4=6=v(f)$, y por (7.8) el flujo es máximo. La solución oficial elige los caminos en este mismo orden, pero cualquier orden lleva al mismo valor, $6$.

\subsection{Cuando hay que deshacer flujo}

En el ejemplo anterior nunca se usó una arista hacia atrás. Los dos ejercicios que siguen están pensados para que sí haga falta.\visto{2P 2017 Ej.~3(a): FF desde un flujo dado}

\begin{problema}{2P 2017, Ej.~3(a): partir de un flujo dado}
Capacidades todas $1$. El flujo inicial manda una unidad por $s\,u\,x\,t$ y otra por $s\,w\,z\,t$, así que $v(f)=2$. Encontrar el flujo máximo con Ford–Fulkerson, mostrando grafos residuales y grafos con flujo.

\smallskip
\textbf{Solución.} En $\flGf$ las aristas $(u,x)$ y $(w,z)$ quedan invertidas. El único camino de aumento es $s\,y\,z\,w\,x\,u\,v\,t$: usa $(z,w)$ y $(x,u)$ hacia atrás y tiene cuello $1$. Al aumentar, se ``deshacen'' $(u,x)$ y $(w,z)$ y quedan los tres caminos $s\,u\,v\,t$, $s\,w\,x\,t$ y $s\,y\,z\,t$. En el nuevo residual no hay camino $s$-$t$, así que el flujo máximo es $3$ (además, el corte $(\{s\},V\setminus\{s\})$ tiene capacidad $3$).
\end{problema}

\begin{center}
\begin{tikzpicture}[scale=0.92]
  \node[flv] (s) at (0,0) {$s$};
  \node[flv] (u) at (1.4,1.2) {$u$}; \node[flv] (v) at (2.8,1.2) {$v$};
  \node[flv] (w) at (1.4,0) {$w$};   \node[flv] (x) at (2.8,0) {$x$};
  \node[flv] (y) at (1.4,-1.2) {$y$}; \node[flv] (z) at (2.8,-1.2) {$z$};
  \node[flv] (t) at (4.2,0) {$t$};
  \draw[fla] (s)--(u); \draw[fla] (s)--(w); \draw[fla] (s)--(y);
  \draw[fla] (u)--(v); \draw[fla] (u)--(x); \draw[fla] (w)--(x);
  \draw[fla] (w)--(z); \draw[fla] (y)--(z);
  \draw[fla] (v)--(t); \draw[fla] (x)--(t); \draw[fla] (z)--(t);
  \draw[flp] (s)--(u); \draw[flp] (u)--(x); \draw[flp] (x)--(t);
  \draw[flp] (s)--(w); \draw[flp] (w)--(z); \draw[flp] (z)--(t);
  \node[flt] at (2.1,-1.95) {$G$ con el flujo inicial ($v=2$)};
\end{tikzpicture}\hspace{0.9cm}%
\begin{tikzpicture}[scale=0.92]
  \node[flv] (s) at (0,0) {$s$};
  \node[flv] (u) at (1.4,1.2) {$u$}; \node[flv] (v) at (2.8,1.2) {$v$};
  \node[flv] (w) at (1.4,0) {$w$};   \node[flv] (x) at (2.8,0) {$x$};
  \node[flv] (y) at (1.4,-1.2) {$y$}; \node[flv] (z) at (2.8,-1.2) {$z$};
  \node[flv] (t) at (4.2,0) {$t$};
  \draw[flr] (u)--(s); \draw[flr] (w)--(s); \draw[flp] (s)--(y);
  \draw[flp] (u)--(v); \draw[flrp] (x)--(u); \draw[flp] (w)--(x);
  \draw[flrp] (z)--(w); \draw[flp] (y)--(z);
  \draw[flp] (v)--(t); \draw[flr] (t)--(x); \draw[flr] (t)--(z);
  \node[flt] at (2.1,-1.95) {$\flGf$: $s\,y\,z\,w\,x\,u\,v\,t$};
\end{tikzpicture}
\end{center}

El segundo ejercicio es el del 2P de 2025: dada la red con un flujo de valor $7$, dibujar el residual y una red con flujo máximo.\visto{2P 2025 Ej.~1(d): residual y flujo máximo}

\begin{problema}{2P 2025, Ej.~1(d): residual y flujo máximo}
Flujo/capacidad: $s\to u$ $5/5$, $u\to t$ $0/4$, $u\to w$ $5/8$, $w\to t$ $7/7$, $s\to v$ $2/15$, $v\to w$ $2/10$. Dibujar $\flGf$ y una red con flujo de valor máximo.

\smallskip
\textbf{Solución.} En $\flGf$: $u\to s$ (5), $u\to t$ (4), $u\to w$ (3) y $w\to u$ (5), $t\to w$ (7), $s\to v$ (13) y $v\to s$ (2), $v\to w$ (8) y $w\to v$ (2). El camino $s\,v\,w\,u\,t$ usa $(w,u)$ hacia atrás y tiene cuello $\min\{13,8,5,4\}=4$. El flujo resultante es $s\to u$ $5/5$, $u\to t$ $4/4$, $u\to w$ $1/8$, $w\to t$ $7/7$, $s\to v$ $6/15$, $v\to w$ $6/10$, de valor $11$. Es máximo porque el corte $(\{s,u,v,w\},\{t\})$ tiene capacidad $4+7=11$.
\end{problema}

\begin{receta}[ejecutar Ford–Fulkerson]
\begin{enumerate}
  \item Dibujá $G$ con $f(e)/c_e$ y al lado $\flGf$. Poné las aristas hacia atrás punteadas y no dibujes las de capacidad residual $0$.
  \item Nombrá el camino elegido (por ejemplo $P_2=s\,a\,t$), su cuello de botella y el nuevo $v(f)$.
  \item Repetí hasta que no haya camino. En el último residual marcá el conjunto $A^*$ de nodos alcanzables desde $s$.
  \item Cerrá con el certificado: ``el corte $(A^*,V\setminus A^*)$ tiene capacidad $c=v(f)$; por (7.8) ningún flujo supera $c$, luego $f$ es máximo''.
\end{enumerate}
\end{receta}

\begin{trampa}
\begin{itemize}
  \item Olvidar las aristas hacia atrás, o dibujarlas con capacidad $c_e-f(e)$ en lugar de $f(e)$.
  \item Calcular el cuello de botella con las capacidades de $G$ en vez de las residuales.
  \item Terminar sin justificar que el flujo es máximo. ``No encuentro más caminos'' no alcanza: hay que exhibir el residual sin camino o el corte.
  \item Usar un camino que repite nodos. El camino de aumento tiene que ser simple.
\end{itemize}
\end{trampa}

% =====================================================================
\section{Flujo máximo y corte mínimo}
% =====================================================================

\begin{resultado}[(7.9)]
Si $f$ es un flujo $s$-$t$ tal que no hay camino $s$-$t$ en $\flGf$, entonces existe un corte $(A^*,B^*)$ con $v(f)=c(A^*,B^*)$. En consecuencia, $f$ tiene valor máximo entre todos los flujos y $(A^*,B^*)$ tiene capacidad mínima entre todos los cortes.
\end{resultado}

En el examen de julio de 2025 se pidió demostrar esto, y la pauta lista lo que se espera: identificar el corte, probar que es un corte, probar que las aristas salientes están saturadas y las entrantes vacías, concluir con (7.6) y deducir la consecuencia con (7.8).\visto{Examen jul.~2025 Ej.~2: demostrar (7.9)}

\begin{proof}
Sea $A^*$ el conjunto de nodos $v$ para los que hay un camino $s$-$v$ en $\flGf$, y sea $B^*=V\setminus A^*$.

\emph{Es un corte.} Es una partición; $s\in A^*$ por el camino trivial y $t\notin A^*$ por hipótesis.

\emph{Aristas que salen de $A^*$.} Sea $e=(u,v)\in E$ con $u\in A^*$, $v\in B^*$. Si fuera $f(e)<c_e$, $e$ sería arista hacia adelante de $\flGf$. Pegándola al camino $s$-$u$ obtendríamos un camino $s$-$v$, contra $v\in B^*$. Entonces $f(e)=c_e$.

\emph{Aristas que entran a $A^*$.} Sea $e=(u,v)\in E$ con $u\in B^*$, $v\in A^*$. Si fuera $f(e)>0$, $(v,u)$ sería arista hacia atrás de $\flGf$ y tendríamos un camino $s$-$u$, contra $u\in B^*$. Entonces $f(e)=0$.

\emph{Conclusión.} Por (7.6),
$v(f)=\flout(A^*)-\flin(A^*)=\sum_{e\text{ sale de }A^*}c_e-0=c(A^*,B^*)$.
Por (7.8), cualquier flujo $f'$ cumple $v(f')\le c(A^*,B^*)=v(f)$ y cualquier corte $(A,B)$ cumple $c(A,B)\ge v(f)=c(A^*,B^*)$.
\end{proof}

\begin{resultado}[(7.10)--(7.14)]
\begin{itemize}
  \item (7.10) El flujo que devuelve Ford–Fulkerson es un flujo máximo.
  \item (7.11) Dado un flujo de valor máximo, se puede calcular un corte de capacidad mínima en tiempo $\Oh(m)$.
  \item (7.12)--(7.13) En toda red, el valor máximo de un flujo $s$-$t$ es igual a la capacidad mínima de un corte $s$-$t$.
  \item (7.14) Si todas las capacidades son enteras, existe un flujo máximo con $f(e)$ entero en toda arista.
\end{itemize}
\end{resultado}

(7.14) es la \emph{propiedad de integralidad}, y todos los ejercicios de modelado la usan. Garantiza que Ford–Fulkerson devuelve un flujo entero, así que una arista de capacidad $1$ lleva $0$ o $1$ y se la puede leer como ``asignado / no asignado''.

\subsection{Encontrar el corte mínimo}

Se preguntó en el examen de diciembre de 2018 (``¿cómo calcularía un corte de capacidad mínima en tiempo $\Oh(m)$ a partir de $\flGf$?'') y en el 2P de 2025. La respuesta es la construcción de la demostración de (7.9).\visto{Examen dic.~2018 Ej.~4: corte mínimo en $\Oh(m)$}

\begin{algorithmic}[1]
\Function{CorteMínimo}{$\flGf$} \Comment{$f$ flujo máximo}
  \State $A\gets$ nodos alcanzados por un BFS (o DFS) desde $s$ en $\flGf$
  \State \Return $(A,\,V\setminus A)$
\EndFunction
\end{algorithmic}

Es correcto por (7.9): como $f$ es máximo, no hay camino $s$-$t$ en $\flGf$, y el conjunto de alcanzables da un corte de capacidad $v(f)$, que es mínima. El BFS cuesta $\Oh(m+n)=\Oh(m)$ porque $\flGf$ tiene a lo sumo $2m$ aristas y $m\ge n/2$.\visto{2P 2025 Ej.~1(c): corte mínimo desde $G_r$}

\subsection{Verdadero o falso sobre cortes y flujos}

En las preguntas de verdadero o falso casi siempre la respuesta es ``falso'' y se justifica con un contraejemplo chico con capacidades $1$. Si es verdadero, se argumenta con flujo máximo = corte mínimo o con integralidad.\visto{Examen feb.~2018 P.~1: unicidad de cortes y flujos}

\begin{problema}{Examen feb.~2018, P.~1: ¿unicidad?}
(a) Dado un flujo máximo, ¿existe un único corte mínimo? (b) ¿Un corte mínimo determina un único flujo máximo?

\smallskip
\textbf{Solución.} (a) Falso. En $s\xrightarrow{1}u\xrightarrow{1}t$, los cortes $(\{s\},\{u,t\})$ y $(\{s,u\},\{t\})$ tienen capacidad $1$. (b) Falso. Con $s\to u\to w$, $s\to v\to w$ y $w\to t$, todas de capacidad $1$, se puede mandar la unidad por $u$ o por $v$. Son dos flujos máximos distintos y en ambos el corte mínimo es $(\{s,u,v,w\},\{t\})$.
\end{problema}

\begin{problema}{Examen jul.~2023, Ej.~3(b): sumar 1 a cada capacidad}
Si $(A,B)$ es un corte mínimo con capacidades $c_e$, ¿sigue siéndolo con capacidades $c_e+1$?

\smallskip
\textbf{Solución.} Falso. Con $s\xrightarrow{3}a$, $s\xrightarrow{4}b$, $a\xrightarrow{1}b$, $a\xrightarrow{1}t$, $b\xrightarrow{5}t$, el corte $(\{s,a\},\{b,t\})$ tiene capacidad $4+1+1=6$ y es mínimo. Al sumar $1$ pasa a valer $5+2+2=9$, mientras que $(\{s,a,b\},\{t\})$ vale $2+6=8$. Sumar $1$ penaliza a los cortes con muchas aristas. Multiplicar todas las capacidades por $\lambda>0$, en cambio, sí conserva los cortes mínimos, porque multiplica todas las capacidades de corte por $\lambda$.
\end{problema}

El ejercicio 1 del 2P de 2020 juega con la paridad: hay que construir un ejemplo y después demostrar una propiedad con flujo máximo = corte mínimo.\visto{2P 2020 Ej.~1: capacidades impares, pares}

\begin{problema}{2P 2020, Ej.~1: paridad}
(a) Red con todas las capacidades impares y flujo máximo par. (b) Si todas las capacidades son pares, el flujo máximo es par.

\smallskip
\textbf{Solución.} (a) $s\to u\to t$ y $s\to v\to t$, todas de capacidad $1$. El flujo que manda $1$ por cada arista cumple capacidad y conservación (en $u$ y en $v$ entra $1$ y sale $1$), y vale $2$. Es máximo porque $c(\{s\},\{u,v,t\})=2$ y vale (7.8). (b) El flujo máximo es igual a la capacidad de un corte mínimo (7.13), que es una suma de capacidades pares y por lo tanto es par. Lo mismo vale para múltiplos de cualquier $k$.
\end{problema}

Otra forma de usar cortes aparece en el examen de febrero de 2023: para probar que \emph{no} se puede satisfacer una demanda alcanza con exhibir un corte de capacidad menor que la demanda, sin correr ningún algoritmo.\visto{Examen feb.~2023 Ej.~1(b): corte que prueba insuficiencia}

\begin{trampa}
\begin{itemize}
  \item Confundir ``el flujo máximo es único'' con ``el valor del flujo máximo es único''. El valor es único; el flujo y el corte no tienen por qué serlo.
  \item Para el corte mínimo, tomar los nodos alcanzables en $G$ en lugar de en $\flGf$.
  \item Decir ``integralidad'' sin aclarar qué da: existe un flujo máximo entero, y Ford–Fulkerson lo encuentra.
\end{itemize}
\end{trampa}

% =====================================================================
\section{Matching bipartito}
% =====================================================================

Sea $G=(X\cup Y,E)$ bipartito. Un \emph{matching} es un conjunto de aristas en el que ningún nodo aparece dos veces. La red $G'$ se arma así: todas las aristas se orientan de $X$ a $Y$, se agrega $s$ con aristas $(s,x)$ para todo $x\in X$ y $t$ con aristas $(y,t)$ para todo $y\in Y$, y todas las capacidades valen $1$.

\begin{center}
\begin{tikzpicture}[scale=0.95]
  \node[flv] (s) at (0,0) {$s$};
  \foreach \i/\y in {1/1.1,2/0,3/-1.1}{
    \node[flv] (x\i) at (1.7,\y) {$x_\i$};
    \node[flv] (y\i) at (3.5,\y) {$y_\i$};}
  \node[flv] (t) at (5.2,0) {$t$};
  \foreach \i in {1,2,3}{\draw[fla] (s)--(x\i); \draw[fla] (y\i)--(t);}
  \draw[fla] (x1)--(y2); \draw[fla] (x2)--(y1); \draw[fla] (x3)--(y3);
  \draw[flp] (x1)--(y1); \draw[flp] (x2)--(y3); \draw[fla] (x3)--(y2);
  \draw[flp] (s)--(x1); \draw[flp] (s)--(x2); \draw[flp] (y1)--(t); \draw[flp] (y3)--(t);
  \node[flt] at (2.6,-1.85) {todas las capacidades valen $1$};
\end{tikzpicture}
\end{center}

\begin{resultado}[(7.37) y (7.38)]
El tamaño de un matching máximo en $G$ es igual al valor de un flujo máximo en $G'$, y las aristas del matching son las aristas de $X$ a $Y$ que llevan flujo. Ford–Fulkerson encuentra un matching máximo en tiempo $\Oh(mn)$.
\end{resultado}

\begin{proof}[Idea]
\emph{De matching a flujo:} un matching de $k$ aristas $(x_i,y_i)$ induce el flujo que manda $1$ por cada camino $s\,x_i\,y_i\,t$, de valor $k$. \emph{De flujo a matching:} por (7.14) hay un flujo máximo entero, de valor $k$. Cada arista lleva $0$ o $1$. Cada $x$ recibe a lo sumo $1$ (una sola arista entrante de capacidad $1$), así que por conservación manda flujo por a lo sumo una arista; lo mismo vale para cada $y$. Las aristas $X\to Y$ con flujo $1$ forman entonces un matching, y hay $k$ de ellas porque el corte $(\{s\}\cup X,\,Y\cup\{t\})$ lo atraviesa todo el flujo. \emph{Tiempo:} $C\le|X|\le n$, así que $\Oh(mC)=\Oh(mn)$.
\end{proof}

En el 2P de 2017 se preguntó cuánto puede medir un camino de aumento en esta red. La respuesta sale de mirar cómo alterna el camino entre $X$ e $Y$.\visto{2P 2017 Ej.~3(b): largo máximo de camino de aumento}

\begin{problema}{2P 2017, Ej.~3(b): largo máximo de un camino de aumento}
Con $|U|=m$ y $|W|=n$, acotar el largo de los caminos de aumento que puede encontrar Ford–Fulkerson.

\smallskip
\textbf{Solución.} El camino va $s\to U\to W\to\cdots\to t$. Para pasar por más de dos nodos del bipartito tiene que volver de $W$ a $U$, y eso solo se puede por aristas hacia atrás. Por lo tanto alterna $U,W,U,W,\dots$ sin repetir nodos, y visita a lo sumo $\min\{m,n\}$ nodos de cada lado. Tiene como máximo $2\min\{m,n\}+2$ nodos, es decir, largo $2\min\{m,n\}+1$. Si $m=n$ puede recorrer todos los nodos ($2m+1$ aristas); el ejemplo de la parte (a) lo alcanza con $m=3$.
\end{problema}

\begin{trampa}
\begin{itemize}
  \item Dejar las aristas del bipartito sin orientar, o sin capacidad. Van de $X$ a $Y$ con capacidad $1$; también sirve $\infty$, pero hay que decirlo.
  \item Leer el matching del flujo sin invocar integralidad. Un flujo fraccionario de $1/2$ en dos aristas no es un matching.
\end{itemize}
\end{trampa}

% =====================================================================
\section{Modelado: asignaciones con capacidades}
% =====================================================================

Este es \emph{el} ejercicio de flujo de los parciales recientes. Aparece con cualquier historia (estudiantes y asignaturas, personas y tareas, laboratorios y países, actividades y días, turnos y departamentos, empresas y trabajos, tipos de sangre), pero la estructura es siempre la misma: matching bipartito generalizado con capacidades en los extremos.\visto{2P 2018 Ej.~1: estudiantes y asignaturas con cupo}

\begin{center}
\begin{tikzpicture}[scale=0.95]
  \node[flv] (s) at (0,0) {$s$};
  \foreach \i/\y in {1/1.2,2/0,3/-1.2}{\node[flv] (p\i) at (2,\y) {$p_\i$};}
  \foreach \j/\y in {1/0.6,2/-0.6}{\node[flv] (q\j) at (4.2,\y) {$q_\j$};}
  \node[flv] (t) at (6.2,0) {$t$};
  \draw[flp] (s)--node[fll,above left]{$a_1$}(p1);
  \draw[flp] (s)--node[fll]{$a_2$}(p2);
  \draw[flp] (s)--node[fll,below left]{$a_3$}(p3);
  \draw[fla] (p1)--node[fll]{1}(q1); \draw[fla] (p2)--node[fll,pos=0.3]{1}(q1);
  \draw[fla] (p2)--node[fll,pos=0.3]{1}(q2); \draw[fla] (p3)--node[fll]{1}(q2);
  \draw[flp] (q1)--node[fll,above right]{$b_1$}(t);
  \draw[flp] (q2)--node[fll,below right]{$b_2$}(t);
  \node[flt,align=center] at (2,-1.95) {elementos de un lado\\(capacidad propia)};
  \node[flt,align=center] at (4.2,-1.95) {elementos del otro\\(cupo)};
\end{tikzpicture}
\end{center}

Hay una arista $(p_i,q_j)$ de capacidad $1$ si y solo si el par $(i,j)$ está permitido. La capacidad $1$ impide que el mismo par se use dos veces. Cada unidad de flujo por $s\,p_i\,q_j\,t$ representa ``se asigna $i$ a $j$''. La capacidad $a_i$ de $(s,p_i)$ es cuánto puede tomar $i$, y la capacidad $b_j$ de $(q_j,t)$ es cuánto admite $j$.

\begin{receta}[modelado con flujo (estilo de las soluciones oficiales)]
\begin{enumerate}
  \item \textbf{Red.} Definí $V$: $s$, $t$ y un nodo por cada elemento. Definí $E$ con su capacidad, separando por tipo de arista: ``para cada $e\in E$ agregamos $(s,v_e)$ con capacidad $k$\ldots''. Decí qué representa una unidad de flujo.
  \item \textbf{Algoritmo.} (1) Construir $G$. (2) Obtener con Ford–Fulkerson un flujo máximo \emph{entero} $f$. (3) Si se pregunta factibilidad, comparar $v(f)$ con la cota $D$ (suma de demandas). (4) Leer la solución: asignar $i$ a $j$ si $f(p_i,q_j)=1$ (o $>0$).
  \item \textbf{Flujo $\Rightarrow$ solución.} Tomá el flujo entero de valor $D$ y verificá cada restricción del enunciado, citando lo que usás. Integralidad (7.2)/(7.14): las aristas de capacidad $1$ llevan $0$ o $1$. Capacidad: no se supera ningún cupo. Conservación: lo que entra a $p_i$ sale por sus aristas. Saturación: $v(f)=D$ y $D$ es la capacidad del corte junto a $t$ (o junto a $s$), así que esas aristas están llenas.
  \item \textbf{Solución $\Rightarrow$ flujo.} Dada una solución válida, definí $f$ explícitamente ($1$ en cada par asignado; en $(s,p_i)$, lo asignado a $i$; en $(q_j,t)$, lo asignado a $j$) y verificá capacidad y conservación. Como Ford–Fulkerson da un flujo máximo (7.10) y $v\le c(\cdot)=D$ (7.8), el algoritmo encuentra valor $D$ y responde que sí.
  \item \textbf{Complejidad.} Contá $|V|$ y $|E|$. Elegí un corte de capacidad $C$ \emph{acotada por cantidades}, por ejemplo $(\{s\}\cup P,\,Q\cup\{t\})$ con aristas de capacidad $1$, que da $C\le|P||Q|$. Aplicá (7.5): $\Oh(|E|\,C)$. Construir la red y leer la solución cuesta $\Oh(|E|)$. Como son una cantidad constante de pasos, domina el más caro.
\end{enumerate}
\end{receta}

Los enunciados siguientes están resumidos. Para cada uno, la solución dice la red, cómo se lee la respuesta y la cota de tiempo.\visto{2P 2019 Ej.~1: personas y tareas, $\Oh(k^2q^2)$}

\begin{problema}{2P 2018, Ej.~1: estudiantes y asignaturas}
Cada asignatura $a$ tiene cupo $q_a$; cada estudiante $e$ se inscribió en $S_e$ y puede cursar a lo sumo $k$. Maximizar el total de asignaciones en tiempo polinomial en $|A|$ y $|E|$.

\smallskip
\textbf{Solución.} $(s,v_e)$ con capacidad $k$; $(v_e,w_a)$ con capacidad $1$ si $a\in S_e$; $(w_a,t)$ con capacidad $q_a$. Se asigna $e$ a $a$ si $f(v_e,w_a)=1$. Hay $x$ asignaciones factibles si y solo si hay un flujo de valor $x$. Cota: $C=k|E|$ por el corte $(\{s\},\dots)$, $m+n=\Oh(|E||A|)$, total $\Oh(|E|^2|A|\,k)$, que es polinomial porque se puede suponer $k\le|A|$.
\end{problema}

\begin{problema}{2P 2019, Ej.~1: personas y tareas con mínimos}
Persona $p_i$ hace tareas de $T_i$, a lo sumo $C_i\le k$. La tarea $t_j$ necesita \emph{al menos} $R_j$ personas. Asignar o informar que no hay solución, en $\Oh(k^2q^2)$.

\smallskip
\textbf{Solución.} $(s,p_i)$ con capacidad $C_i$; $(p_i,t_j)$ con capacidad $1$ si $t_j\in T_i$; $(t_j,t)$ con capacidad $R_j$. Sea $R=\sum_j R_j$. Hay solución si y solo si $v(f)=R$. La clave está en la vuelta: si una solución asigna más de $R_j$ personas a $t_j$, se liberan las sobrantes y queda una con \emph{exactamente} $R_j$, que da un flujo de valor $R$. Cota: $m=\Oh(qk)$ y el corte $(\{s\}\cup P,\,T\cup\{t\})$ tiene capacidad $\le qk$, así que el total es $\Oh(q^2k^2)$. (También sirve $(\{s\},\dots)$ con $\sum C_i\le qk$.)
\end{problema}

\begin{problema}{Examen dic.~2020, Ej.~2: vacunas}
El laboratorio $\ell_i$ produce $c(\ell_i)$ dosis; el país $p_j$ demanda $d(p_j)$; solo algunos pares $(\ell_i,p_j)$ están permitidos. Decidir si se cubre toda la demanda y dar la matriz $A_{ij}$ de dosis.

\smallskip
\textbf{Solución.} $(s,\ell_i)$ con capacidad $c(\ell_i)$; $(\ell_i,p_j)$ con capacidad $d(p_j)$ (o $\infty$) si el par está permitido; $(p_j,t)$ con capacidad $d(p_j)$. Se responde que sí si $f(p_j,t)=d(p_j)$ para todo $j$, y entonces $A_{ij}=f(\ell_i,p_j)$. Corrección: por conservación, a cada país llegan $d(p_j)$ dosis y de cada laboratorio salen a lo sumo $c(\ell_i)$; solo se usan pares permitidos. Acá las aristas del medio \emph{no} tienen capacidad $1$, porque se mandan muchas dosis por el mismo par.
\end{problema}

En este último caso la cota $\Oh(mC)$ con $C=\sum_j d(p_j)$ es pseudopolinomial. La letra no pedía tiempo polinomial, así que no hacía falta más.\visto{Examen dic.~2020 Ej.~2: laboratorios y países}

\begin{problema}{Examen jul.~2021, Ej.~1: actividades y días}
La actividad $i$ dura $d_i$ días, que tienen que caer en $[s_i,f_i]$; cada día se trabaja en una sola actividad. Asignar días o decir que no se puede, en $\Oh(kn^2)$.

\smallskip
\textbf{Solución.} $(s,i)$ con capacidad $d_i$; $(i,j)$ con capacidad $1$ para cada día $j\in[s_i,f_i]$; $(j,t)$ con capacidad $1$. Sea $D=\sum d_i$. Hay solución si y solo si $v(f)=D$, y a la actividad $i$ le tocan los días $j$ con $f(i,j)=1$. Cota: $|E|=\Oh(kn)$ y el corte $(V\setminus\{t\},\{t\})$ tiene capacidad $n$, así que el total es $\Oh(kn^2)$.
\end{problema}

El ejercicio de julio de 2021 muestra bien la elección del corte para la complejidad. El corte junto a $s$ tiene capacidad $D=\sum d_i$, que depende de números de la entrada; el corte junto a $t$ tiene capacidad $n$ y da la cota exacta que se pedía.\visto{Examen jul.~2021 Ej.~1: actividades y días}

\begin{problema}{Examen dic.~2021, Ej.~3(b): enfermeros por turno y departamento}
El turno $T_i$ tiene $E_i$ enfermeros; el departamento $D_j$ recibe a lo sumo $R_j$ por día; en el turno $i$, el departamento $j$ recibe a lo sumo $M_{ij}$. Maximizar el personal asignado y dar cuántos de cada turno van a cada departamento.

\smallskip
\textbf{Solución.} $(s,T_i)$ con capacidad $E_i$; $(T_i,D_j)$ con capacidad $M_{ij}$; $(D_j,t)$ con capacidad $R_j$. Con Ford–Fulkerson, la cantidad del turno $i$ asignada a $D_j$ es $f(T_i,D_j)$. Con tres turnos y tres departamentos la red tiene tamaño constante.
\end{problema}

\begin{problema}{2P 2023, Ej.~2: directivos y trabajos comunitarios}
La empresa $e_i$ tiene $d_i$ directivos, que tienen que ir a trabajos \emph{distintos}; el trabajo $t_j$ tiene cupo $c_j$. ¿Existe una asignación válida? Tiempo polinomial en $h$ y $k$.

\smallskip
\textbf{Solución.} $(s,e_i)$ con capacidad $d_i$; $(e_i,t_j)$ con capacidad $1$ para todo par (la capacidad $1$ obliga a que los trabajos sean distintos); $(t_j,z)$ con capacidad $c_j$. Hay solución si y solo si $v(f)=n=\sum_i d_i$. La solución oficial prueba la ida con la saturación de $(s,e_i)$ más integralidad, capacidad y conservación, y la vuelta construyendo el flujo y usando (7.10) y (7.8). Cota: el corte $(\{s,e_1,\dots,e_h\},\{t_1,\dots,t_k,z\})$ tiene capacidad $hk$ y $|E|=\Oh(hk)$, así que el total es $\Oh(h^2k^2)$.
\end{problema}

La solución del 2P de 2023 es la más completa de las oficiales. Sigue paso a paso la receta: (a.1) saturación de las aristas de $s$, (a.2) flujo $0$ o $1$ en el medio por integralidad, (a.3) cupos respetados por capacidad, y después la verificación de cada condición de la asignación.\visto{2P 2023 Ej.~2: empresas, trabajos distintos}

\begin{problema}{Examen feb.~2023, Ej.~1: transfusiones de sangre}
Hay oferta $s_O,s_A,s_B,s_{AB}$ y demanda $d_O,d_A,d_B,d_{AB}$. $O$ dona a todos, $A$ dona a $A$ y $AB$, $B$ dona a $B$ y $AB$, y $AB$ solo a $AB$. (a) Modelar y decidir si alcanza. (b) Con oferta $50,36,11,8$ y demanda $45,42,10,3$, mostrar con un corte que no alcanza y dar una asignación máxima.

\smallskip
\textbf{Solución.} (a) $(s,\text{oferta}_X)$ con capacidad $s_X$; $(\text{oferta}_X,\text{demanda}_Y)$ con capacidad $\infty$ si $Y$ puede recibir de $X$; $(\text{demanda}_Y,t)$ con capacidad $d_Y$. Alcanza si y solo si todas las aristas hacia $t$ quedan saturadas. (b) Sea $A=\{s,\text{of}_B,\text{of}_{AB},\text{dem}_B,\text{dem}_{AB}\}$. Las aristas que salen de $A$ son $(s,\text{of}_O)$, $(s,\text{of}_A)$, $(\text{dem}_B,t)$ y $(\text{dem}_{AB},t)$, de capacidad $50+36+10+3=99<100$. Una asignación de $99$: $O\to O$ $45$, $O\to A$ $5$, $A\to A$ $36$, $B\to B$ $10$, $B\to AB$ $1$, $AB\to AB$ $2$.
\end{problema}

La parte (b) del ejercicio de sangre es el uso típico de un corte como certificado. Las capacidades $\infty$ obligan a que el corte no deje una oferta del lado de $s$ si la demanda que esa oferta puede cubrir queda del lado de $t$. Por eso se elige $A$ cerrado ``hacia adelante'': $B$ y $AB$ solo donan a $B$ y $AB$.\visto{Examen feb.~2023 Ej.~1: tipos de sangre}

El ejercicio de febrero de 2024 agrega una restricción que no se puede expresar solo con capacidades en los extremos. Se resuelve con una capa intermedia de nodos.\visto{Examen feb.~2024 Ej.~3: cursos con horarios superpuestos}

\begin{problema}{Examen feb.~2024, Ej.~3: cursos con superposición horaria}
Como en 2P 2018, pero el estudiante $e$ puede cursar a lo sumo $p_e$ cursos y hay una partición $\mathcal S$ de los cursos en franjas horarias: a cada estudiante se le asigna a lo sumo un curso por franja. Sugerencia: $V=\{s,t\}\cup\mathcal E\cup W\cup C$ con $W=\mathcal E\times\mathcal S$.

\smallskip
\textbf{Solución.} $(s,e)$ con capacidad $p_e$; $(e,(e,S))$ con capacidad $1$ si $S\cap C_e\ne\emptyset$; $((e,S_c),c)$ con capacidad $1$ para cada $c\in C_e$; $(c,t)$ con capacidad $q_c$. Como por cada nodo $(e,S)$ entra a lo sumo $1$, por conservación sale a lo sumo $1$, y eso da ``a lo sumo un curso por franja''. Cota: $|W|\le|\mathcal E||C|$ y el corte $(\{s\}\cup\mathcal E,\dots)$ tiene capacidad $\le|\mathcal E||\mathcal S|\le|\mathcal E||C|$, así que es polinomial.
\end{problema}

\begin{center}
\begin{tikzpicture}[scale=0.95]
  \node[flv] (s) at (0,0) {$s$};
  \node[flv] (e) at (1.6,0) {$e$};
  \node[flv,minimum size=7mm,font=\scriptsize] (w1) at (3.4,0.8) {$e,S_1$};
  \node[flv,minimum size=7mm,font=\scriptsize] (w2) at (3.4,-0.8) {$e,S_2$};
  \node[flv] (c1) at (5.2,1.3) {$c_1$};
  \node[flv] (c2) at (5.2,0.3) {$c_2$};
  \node[flv] (c3) at (5.2,-0.8) {$c_3$};
  \node[flv] (t) at (6.8,0.2) {$t$};
  \draw[flp] (s)--node[fll,above]{$p_e$}(e);
  \draw[flp] (e)--node[fll,above left]{1}(w1);
  \draw[flp] (e)--node[fll,below left]{1}(w2);
  \draw[fla] (w1)--node[fll]{1}(c1); \draw[fla] (w1)--node[fll]{1}(c2);
  \draw[fla] (w2)--node[fll]{1}(c3);
  \draw[fla] (c1)--node[fll,above right]{$q_{c_1}$}(t);
  \draw[fla] (c2)--node[fll]{$q_{c_2}$}(t);
  \draw[fla] (c3)--node[fll,below right]{$q_{c_3}$}(t);
  \node[flt] at (3.4,-1.6) {franja $S_1=\{c_1,c_2\}$, franja $S_2=\{c_3\}$};
\end{tikzpicture}
\end{center}

Este truco (un nodo por par ``elemento $\times$ grupo'' con una arista de capacidad $1$ que entra) sirve para cualquier restricción del tipo ``a lo sumo $r$ por grupo'': la arista que entra a $(e,S)$ lleva capacidad $r$.

\begin{trampa}
\begin{itemize}
  \item Probar una sola dirección. La equivalencia ``hay solución $\iff$ $v(f)=D$'' se demuestra en los dos sentidos, salvo que la letra diga lo contrario (como en dic.~2020, que pedía solo una).
  \item En la vuelta, decir ``la solución induce un flujo'' sin definirlo arista por arista ni verificar capacidad y conservación.
  \item No invocar integralidad al leer la asignación de un flujo.
  \item Usar capacidad $\infty$ en las aristas del medio cuando el par no se puede repetir (directivos, estudiantes). La capacidad $1$ es la que impide repetir.
  \item Cota de tiempo con $C=\sum$ (cupos) cuando se pide polinomial en las cantidades: no es polinomial. Buscá el corte con aristas de capacidad $1$ o el corte junto a $t$.
  \item ``Al menos $R_j$'' modelado con capacidad $\infty$: así no se fuerza nada. Se exige exactamente $R_j$ y se argumenta que alcanza.
\end{itemize}
\end{trampa}

% =====================================================================
\section{Mapa de lo que se preguntó}
% =====================================================================

Relevamiento de los segundos parciales 2008--2025 y de los exámenes 2009--2026. En los parciales y exámenes anteriores a noviembre de 2017 no aparece flujo, y en los segundos parciales de 2021 y 2024 tampoco hubo.

\begin{center}
\small
\begin{tabularx}{\textwidth}{@{}>{\raggedright\arraybackslash}p{2.75cm}>{\raggedright\arraybackslash}p{2.5cm}>{\raggedright\arraybackslash}X@{}}
\toprule
Fuente & Tema & Qué se pedía \\
\midrule
2P 2017 Ej.~3(a) & FF a mano & Flujo máximo desde un flujo dado, con residuales; necesita aristas hacia atrás. \\
2P 2017 Ej.~3(b) & Matching bipartito & Largo máximo de un camino de aumento: $2\min\{m,n\}+1$. \\
Ex.\ feb.~2018 P.~1 & Cortes (V/F) & ¿Corte mínimo único? ¿Flujo máximo único dado el corte? Contraejemplos. \\
2P 2018 Ej.~1 & Modelado & Estudiantes y asignaturas con cupo y máximo $k$; complejidad polinomial. \\
Ex.\ dic.~2018 Ej.~4 & Corte mínimo & Calcular un corte mínimo desde $\flGf$ en $\Oh(m)$, es (7.11). \\
2P 2019 Ej.~1 & Modelado & Personas y tareas, mínimo $R_j$ por tarea, máximo $C_i$; $\Oh(k^2q^2)$. \\
2P 2020 Ej.~1 & Max-flujo/min-corte & Capacidades impares con máximo par; capacidades pares implican máximo par. \\
Ex.\ jul.~2020 Ej.~2 & Teoría & Condición de capacidad; probar $v(f)\le c(A,B)$; FF en $\Oh(mC)$. \\
Ex.\ dic.~2020 Ej.~2 & Modelado & Laboratorios y países (vacunas); probar el sentido flujo $\Rightarrow$ solución. \\
Ex.\ jul.~2021 Ej.~1 & Modelado & Actividades en ventanas de días; $\Oh(kn^2)$ con el corte junto a $t$. \\
Ex.\ dic.~2021 Ej.~3 & Residual y modelado & Definir grafo residual; enfermeros por turno y departamento. \\
\bottomrule
\end{tabularx}
\end{center}

\begin{center}
\small
\begin{tabularx}{\textwidth}{@{}>{\raggedright\arraybackslash}p{2.75cm}>{\raggedright\arraybackslash}p{2.5cm}>{\raggedright\arraybackslash}X@{}}
\toprule
Fuente & Tema & Qué se pedía \\
\midrule
2P 2022 Ej.~2 & Teoría & Definir red, flujo, valor, corte y capacidad; demostrar (7.8). \\
Ex.\ feb.~2023 Ej.~1 & Modelado y cortes & Tipos de sangre; corte de capacidad $99<100$ como certificado; asignación máxima. \\
Ex.\ jul.~2023 Ej.~3 & FF a mano y V/F & Ejecutar FF con $G$ y $\flGf$; ¿sumar $1$ a las capacidades conserva el corte mínimo? \\
2P 2023 Ej.~2 & Modelado & Directivos y trabajos distintos con cupos; ida, vuelta y $\Oh(h^2k^2)$. \\
Ex.\ feb.~2024 Ej.~3 & Modelado & Cursos con franjas horarias: capa intermedia $\mathcal E\times\mathcal S$. \\
Ex.\ jul.~2025 Ej.~2 & Teoría & Demostrar (7.9): sin camino en $\flGf$ hay un corte con $v(f)=c(A,B)$. \\
2P 2025 Ej.~1(b--d) & Teoría y FF & Definir corte; corte mínimo desde $G_r$; residual y flujo máximo de una red. \\
Ex.\ dic.~2025 Ej.~1(d) & Teoría & Definición de flujo $s$-$t$. \\
Ex.\ feb.~2026 Ej.~1(e) & Complejidad & Cota $\Oh(mC)$ de FF: qué es $m$, qué es $C$ y por qué aparecen. \\
\bottomrule
\end{tabularx}
\end{center}

\textbf{Lectura rápida.} En los segundos parciales el flujo vale entre 10 y 17 puntos. Casi siempre es un ejercicio de modelado (2018, 2019, 2023) o un bloque de teoría y ejecución (2017, 2020, 2022, 2025). En los exámenes alternan las mismas dos cosas. Lo que más se repite: armar la red de asignación, probar los dos sentidos con integralidad, acotar el tiempo con un corte bien elegido, demostrar (7.8) o (7.9) y obtener el corte mínimo desde el residual.

% =====================================================================
%  Parte 4 — Reducciones y NP-completitud
%  Fragmento: se incluye con \input desde el documento maestro.
% =====================================================================

% ---------- macros locales de esta parte ----------
\newcommand{\npdato}[1]{{\sans\bfseries\color{acento}#1.}\enspace}
\newcommand{\npsi}{\emph{sí}}
\newcommand{\npno}{\emph{no}}
\newcommand{\npIS}{\prob{Independent Set}}
\newcommand{\npVC}{\prob{Vertex Cover}}
\newcommand{\npSC}{\prob{Set Cover}}
\newcommand{\npSP}{\prob{Set Packing}}
\newcommand{\npCL}{\prob{Clique}}
\newcommand{\npSAT}{\prob{sat}}
\newcommand{\npTSAT}{\prob{3-sat}}
\newcommand{\npSS}{\prob{Subset Sum}}
\newcommand{\npTDM}{\prob{3d-Matching}}
\newcommand{\npTSP}{\prob{tsp}}
\newcommand{\npneg}[1]{\overline{#1}}

\parte{np}{Reducciones y \texorpdfstring{$\NP$}{NP}-completitud}%
{\begin{gathered}X\ \NP\text{-completo},\ \ X\redp Y,\ \ Y\in\NP\\ \Longrightarrow\ \ Y\ \NP\text{-completo}\end{gathered}}%
{Para probar que un problema nuevo es difícil no se lo resuelve: se muestra que, si se lo pudiera resolver rápido, también se resolvería rápido uno que ya sabemos que es difícil. Todo el capítulo es aprender a hacer esa transformación y a escribirla bien.}

% =====================================================================
\section{Conceptos: qué es reducir y qué es \texorpdfstring{$\NP$}{NP}}
% =====================================================================

\subsection{Problemas de decisión y de optimización}

En toda esta parte los problemas son \emph{de decisión}: la respuesta es \npsi\ o \npno. Formalmente, un problema $X$ es el conjunto de tiras (instancias) cuya respuesta es \npsi, y un algoritmo $A$ resuelve $X$ si para toda entrada $s$ vale $A(s)=\text{\npsi}\iff s\in X$. El algoritmo es polinomial si su tiempo es $\Oh(p(|s|))$ para algún polinomio $p$, donde $|s|$ es el largo en bits de la entrada.

Los problemas de optimización se convierten en problemas de decisión agregando un umbral $k$: si se maximiza, la pregunta es <<¿existe una solución de valor \emph{al menos} $k$?>>; si se minimiza, <<¿existe una de valor \emph{a lo sumo} $k$?>>. La versión de decisión no es más difícil que la de optimización (se resuelve la de optimización y se compara con $k$), y tampoco es mucho más fácil: con búsqueda binaria sobre $k$ alcanzan $\Oh(\log n)$ llamadas a la versión de decisión para obtener el valor óptimo. Por eso, para hablar de dificultad, alcanza con mirar la versión de decisión.

\subsection{Reducciones polinomiales}

\begin{definicion}[Kleinberg--Tardos, sección 8.1]
$Y\redp X$ (<<$Y$ se reduce polinomialmente a $X$>>) si cualquier instancia de $Y$ se puede resolver con una cantidad polinomial de pasos de cómputo más una cantidad polinomial de llamadas a una caja negra que resuelve $X$.
\end{definicion}

En la práctica, y en todas las soluciones oficiales, la reducción tiene una forma más restringida: una \emph{transformación} $f$ que toma una instancia $s$ de $Y$, corre en tiempo polinomial en $|s|$, y cumple
\[
  s \text{ es instancia \npsi\ de } Y \iff f(s) \text{ es instancia \npsi\ de } X .
\]
Después se llama una sola vez a la caja negra de $X$ con $f(s)$ y se devuelve su respuesta. Notá que $|f(s)|$ queda polinomial en $|s|$ automáticamente, porque $f$ no tiene tiempo de escribir algo más largo.

\begin{resultado}[Propiedades que se citan en las soluciones]
\begin{enumerate}[label=(\roman*)]
  \item Si $Y\redp X$ y $X$ se resuelve en tiempo polinomial, entonces $Y$ también.
  \item Si $Y\redp X$ y $Y$ no se puede resolver en tiempo polinomial, entonces $X$ tampoco.
  \item Transitividad: si $Z\redp Y$ e $Y\redp X$, entonces $Z\redp X$ (se componen las transformaciones; la composición de polinomios es un polinomio).
  \item Con la definición del libro, si $Y\in\Pclass$ entonces $Y\redp X$ para cualquier $X$: se resuelve $Y$ directamente y se llama cero veces a la caja negra.
\end{enumerate}
\end{resultado}

La lectura que hay que tener automatizada: $Y\redp X$ quiere decir <<$X$ es al menos tan difícil como $Y$>>, o <<si tuviera un algoritmo para $X$, resolvería $Y$>>. Para probar que un problema \emph{nuevo} es difícil, el problema \emph{conocido} va a la izquierda: se transforman instancias del conocido en instancias del nuevo.

\begin{problema}{Ejemplo base: Independent Set $\redp$ Vertex Cover}
En un grafo $G=(V,E)$, $S$ es un conjunto independiente si y solo si $V\setminus S$ es un cubrimiento de vértices: una arista con los dos extremos en $S$ es exactamente una arista que $V\setminus S$ no cubre.

\npdato{Transformación} dada $(G,k)$ de \npIS\ con $n$ vértices, se construye $(G,\,n-k)$ de \npVC. Es $\Oh(1)$ si $n$ está en la representación, $\Oh(n+m)$ si hay que copiar el grafo.

\npdato{Equivalencia} si $S$ es independiente con $|S|\ge k$, $V\setminus S$ es un cubrimiento con $|V\setminus S|\le n-k$; recíprocamente, si $C$ es un cubrimiento con $|C|\le n-k$, $V\setminus C$ es independiente de tamaño $\ge k$. La misma idea da $\npVC\redp\npIS$: los dos problemas son equivalentes.
\end{problema}

Esta transformación se pidió tal cual, sin demostración, como pregunta corta.\visto{2P 2025 Ej. 1(e): IS $\redp$ VC}

\subsection{Certificadores, \texorpdfstring{$\Pclass$}{P} y \texorpdfstring{$\NP$}{NP}}

$\Pclass$ es la clase de los problemas de decisión que tienen un algoritmo polinomial. $\NP$ es la clase de los problemas para los que una respuesta \npsi\ se puede \emph{verificar} rápido si alguien nos da una prueba. La definición que se pide de memoria es la siguiente.

\begin{definicion}[certificador eficiente]
Un algoritmo $B$ es un certificador eficiente para el problema $X$ si
\begin{enumerate}[label=(\arabic*)]
  \item $B$ recibe dos argumentos, una instancia $s$ y un certificado $t$;
  \item el tiempo de ejecución de $B$ es polinomial en $|s|+|t|$;
  \item existe un polinomio $p$ tal que, para toda tira $s$: $s\in X$ si y solo si existe $t$ con $|t|\le p(|s|)$ y $B(s,t)=\text{\npsi}$.
\end{enumerate}
\end{definicion}

\begin{definicion}
$\NP$ es la clase de los problemas de decisión que tienen un certificador eficiente.
\end{definicion}

La condición (3) tiene dos direcciones y en el parcial hay que probar las dos: si $s$ es \npsi\ \emph{existe} un certificado corto que $B$ acepta, y si $B$ acepta \emph{algún} $t$, entonces $s$ es \npsi. Se pidieron las tres definiciones (certificador, $\NP$, $\NP$-completo) como pregunta teórica pura.\visto{Examen feb. 2019 Ej. 3: definiciones}

\begin{itemize}
  \item $\Pclass\subseteq\NP$: si $A$ resuelve $X$ en tiempo polinomial, el certificador $B(s,t)=A(s)$ ignora $t$ (se toma $t$ vacío).
  \item Todo problema de $\NP$ se resuelve en tiempo exponencial: se prueban los $2^{p(|s|)}$ certificados posibles.
  \item $\NP$ \emph{no} quiere decir <<no polinomial>>; quiere decir <<polinomial no determinístico>>, o sea, <<verificable en tiempo polinomial>>. Lo que no se sabe es si $\Pclass=\NP$.
\end{itemize}

El certificado canónico es el de \npSAT: una asignación de valores de verdad a las variables. El certificador evalúa cada cláusula con esa asignación y responde \npsi\ si todas se satisfacen; es lineal en el tamaño de la fórmula. Esta descripción se pidió en dos exámenes recientes, sin demostración de corrección.\visto{Examen dic. 2025 Ej. 1(e): certificado de SAT}

Otros certificados típicos: para \npIS, \npVC\ o \npSC, el conjunto elegido (como arreglo de bits o lista de índices); para ciclo hamiltoniano, la lista de vértices en el orden del ciclo; para 3-coloreo, el color de cada vértice; para \npSS, el subconjunto.\visto{Examen jul. 2026 Ej. 1(f): certificado de SAT}

\subsection{\texorpdfstring{$\NP$}{NP}-completo, \texorpdfstring{$\NP$}{NP}-difícil y Cook--Levin}

\begin{definicion}
Un problema $X$ es $\NP$-completo si (a) $X\in\NP$ y (b) para todo $Y\in\NP$ vale $Y\redp X$. Si se cumple (b) pero no necesariamente (a), $X$ es $\NP$-difícil.
\end{definicion}

\begin{resultado}[Los dos hechos que se usan siempre]
\begin{enumerate}[label=(\roman*)]
  \item Si $X$ es $\NP$-completo, entonces $X$ se resuelve en tiempo polinomial si y solo si $\Pclass=\NP$. (Si $X\in\Pclass$, todo $Y\in\NP$ cumple $Y\redp X$ y por la propiedad (i) $Y\in\Pclass$. Si $\Pclass=\NP$, como $X\in\NP$ entonces $X\in\Pclass$.)
  \item Si $Z$ es $\NP$-completo, $X\in\NP$ y $Z\redp X$, entonces $X$ es $\NP$-completo. (Para $Y\in\NP$: $Y\redp Z$ por ser $Z$ completo, $Z\redp X$ por hipótesis, y por transitividad $Y\redp X$.)
\end{enumerate}
\end{resultado}

El hecho (ii) es toda la estrategia: nunca se prueba (b) a mano para todos los $Y$; se elige un $Z$ conocido y se prueba una sola reducción $Z\redp X$. La definición de $\NP$-completo y la explicación de cómo usarla se pidieron como pregunta corta.\visto{Examen dic. 2025 Ej. 1(f): definir $\NP$-completo}

Hace falta un primer problema $\NP$-completo para arrancar la cadena. Ese es el teorema de Cook--Levin. El libro lo prueba para \prob{Circuit Satisfiability}: dado un circuito booleano con compuertas \textsc{and}, \textsc{or}, \textsc{not} y algunas entradas libres, ¿hay valores para las entradas que hagan que la salida sea 1? La idea de la prueba: si $X\in\NP$ tiene certificador $B$, para cada tamaño de entrada se puede <<desenrollar>> el cómputo polinomial de $B$ en un circuito de tamaño polinomial; se fijan las entradas correspondientes a $s$ y se dejan libres las de $t$. El circuito es satisfacible si y solo si existe $t$ con $B(s,t)=\text{\npsi}$. Después, $\prob{Circuit-sat}\redp\npTSAT$ con una variable por compuerta y unas pocas cláusulas por compuerta que fuerzan su valor; con eso \npTSAT\ queda $\NP$-completo y de ahí sale el resto del árbol.

En el curso, \npSAT, \npTSAT\ y todos los problemas del catálogo de la sección~\ref{sec:np-catalogo} se dan por conocidos $\NP$-completos: se pueden usar sin demostrarlo, citando el libro.

\subsection{Más allá de \texorpdfstring{$\NP$}{NP}: co-NP y PSPACE}

En los parciales relevados no se preguntó nada de esto, pero conviene tener la idea para no confundirse en un verdadero o falso.

\begin{itemize}
  \item $X\in\coNP$ si su complemento $\overline{X}$ (las instancias \npno) está en $\NP$. Ejemplos: <<¿la fórmula es insatisfacible?>>, <<¿el grafo \emph{no} tiene ciclo hamiltoniano?>>. No se sabe si $\NP=\coNP$. Como $\Pclass$ es cerrada por complemento, si $\NP\ne\coNP$ entonces $\Pclass\ne\NP$.
  \item $\Pclass\subseteq\NP\cap\coNP$. Un problema en $\NP\cap\coNP$ tiene una <<buena caracterización>>: hay certificado corto tanto para el \npsi\ como para el \npno. Ejemplo del curso: <<¿hay un flujo de valor al menos $v$?>>; un flujo certifica el \npsi\ y un corte de capacidad menor a $v$ certifica el \npno.
  \item $\mathsf{PSPACE}$ es la clase de problemas que se resuelven con memoria polinomial. Vale $\Pclass\subseteq\NP\subseteq\mathsf{PSPACE}$ (se prueban todos los certificados reutilizando el mismo espacio). El problema completo típico es \prob{qsat} (fórmulas con cuantificadores alternados), que modela juegos de dos jugadores y planificación.
\end{itemize}

\subsection{Preguntas conceptuales de verdadero o falso}

Las preguntas de este tipo se contestan con una tabla mental de implicaciones. Lo único que importa es de qué lado de $\redp$ está cada problema.

\begin{resultado}[Qué se puede deducir de $Y\redp X$]
\begin{itemize}
  \item $X\in\Pclass\Rightarrow Y\in\Pclass$. \quad $Y\notin\Pclass\Rightarrow X\notin\Pclass$.
  \item $Y\in\Pclass$: no dice nada de $X$. \quad $X$ $\NP$-completo: no dice nada de $Y$ (salvo si $Y\in\NP$, que ya se sabía).
  \item $Y$ $\NP$-completo y $X\in\NP$: $X$ es $\NP$-completo. \quad $Y$ $\NP$-completo y $X\in\Pclass$: $\Pclass=\NP$.
  \item Para todo $Y\in\NP$ y todo $Z$ $\NP$-completo: $Y\redp Z$, siempre. En particular todo problema de $\Pclass$ se reduce a cualquier $\NP$-completo.
  \item Si $\Pclass\ne\NP$: ningún $\NP$-completo está en $\Pclass$, y ningún $\NP$-completo se reduce a un problema de $\Pclass$.
\end{itemize}
\end{resultado}

\begin{problema}{Sí, no o no se sabe (2P 2017, Ej. 4)}
Sea \prob{Interval Scheduling} en versión de decisión: dada una colección de intervalos y una cota $k$, ¿hay al menos $k$ intervalos que no se superponen? Para cada afirmación, responder <<sí>>, <<no>> o <<no se sabe; si se supiera, se resolvería $\Pclass=\NP$>>.

\npdato{(a) ¿Interval Scheduling $\redp$ Vertex Cover?} Sí. Interval Scheduling está en $\Pclass$ (greedy por tiempo de finalización), así que se resuelve en tiempo polinomial con una caja negra de \npVC\ llamándola cero veces. Otra justificación: Interval Scheduling está en $\NP$ y todo problema de $\NP$ se reduce a \npVC, que es $\NP$-completo.

\npdato{(b) ¿Independent Set $\redp$ Interval Scheduling?} No se sabe; es equivalente a $\Pclass=\NP$. Si fuera cierto, como Interval Scheduling está en $\Pclass$, \npIS\ estaría en $\Pclass$ y por ser $\NP$-completo $\Pclass=\NP$. Recíprocamente, si $\Pclass=\NP$, \npIS\ se resuelve en tiempo polinomial y se reduce a cualquier cosa.
\end{problema}

Esta pregunta (la misma, con las mismas respuestas) aparece en la práctica de segundo parcial del curso.\visto{2P 2017 Ej. 4: sí / no / no se sabe}

\begin{problema}{Cuatro afirmaciones (examen jul. 2018, Pregunta 2)}
Sean $X\in\Pclass$, $Y\in\NP$ y $Z$ $\NP$-completo. Decir si cada afirmación es siempre verdadera, siempre falsa o no se puede afirmar; en el último caso, qué tendría que cumplirse para que sea verdadera.

\npdato{1. $Y\redp Z$} Siempre verdadera: es la definición de $\NP$-completo.

\npdato{2. $X\redp Z$} Siempre verdadera: $X\in\Pclass\subseteq\NP$ y se aplica el punto 1.

\npdato{3. $Z\redp Y$} No se puede afirmar. Es verdadera exactamente cuando $Y$ es $\NP$-completo: si vale, $Y\in\NP$ y es al menos tan difícil como un $\NP$-completo.

\npdato{4. $Z\redp X$} No se puede afirmar. Si fuera verdadera, $Z$ se resolvería en tiempo polinomial y entonces $\Pclass=\NP$; recíprocamente, si $\Pclass=\NP$ vale (propiedad (iv)). Es verdadera si y solo si $\Pclass=\NP$.
\end{problema}

Ninguna de las dos preguntas pide demostraciones: piden la justificación de una o dos oraciones basada en las propiedades de arriba.\visto{Examen jul. 2018 Preg. 2: X, Y, Z}

Para practicar, otras afirmaciones del mismo estilo (no salieron tal cual, pero se contestan con la misma tabla):

\begin{enumerate}[label=\alph*)]
  \item <<Si $Y\redp X$ e $Y\in\Pclass$, entonces $X\in\Pclass$>>. Falso en general: $Y$ fácil no dice nada de $X$ (todo problema de $\Pclass$ se reduce a \npSAT).
  \item <<Si $X\redp\npTSAT$, entonces $X$ es $\NP$-completo>>. Falso: prueba que $X$ es \emph{a lo sumo} tan difícil como \npTSAT. La dirección correcta es $\npTSAT\redp X$ (y además $X\in\NP$).
  \item <<\npSS\ tiene un algoritmo $\Oh(nW)$, por lo tanto $\Pclass=\NP$>>. Falso: $\Oh(nW)$ es exponencial en la cantidad de bits de $W$.
  \item <<Si $\Pclass\ne\NP$, todo problema de $\NP$ que no está en $\Pclass$ es $\NP$-completo>>. Falso: si $\Pclass\ne\NP$ existen problemas intermedios (teorema de Ladner). No hace falta saberlo para el parcial, pero sí no afirmar lo contrario.
  \item <<Dos problemas $\NP$-completos se reducen uno al otro>>. Verdadero: ambos están en $\NP$ y cada uno es completo.
  \item <<Si un problema $\NP$-completo está en $\coNP$, entonces $\NP=\coNP$>>. Verdadero (sección 8.9 del libro), aunque no se usa en los parciales.
\end{enumerate}

% =====================================================================
\section{Catálogo de problemas \texorpdfstring{$\NP$}{NP}-completos del libro}
\label{sec:np-catalogo}
% =====================================================================

Estos son los problemas que se pueden usar como <<problema de referencia>> sin demostrar que son $\NP$-completos. Para cada uno: la versión de decisión, el certificado y la reducción con la que el libro lo prueba. En los parciales nunca se pidió reproducir una reducción del libro con gadgets complicados, pero sí reconocer los problemas y saber de cuál viene cada uno.

\subsection{Empaquetamiento y cobertura}

\npdato{\npIS} dado un grafo $G=(V,E)$ y un natural $k$, ¿existe $S\subseteq V$ con $|S|\ge k$ y sin aristas entre vértices de $S$? Certificado: $S$. Se prueba desde \npTSAT\ (ver abajo).

\npdato{\npVC} dado $G$ y $k$, ¿existe $C\subseteq V$ con $|C|\le k$ tal que toda arista tiene al menos un extremo en $C$? Certificado: $C$. Se prueba desde \npIS\ con $(G,k)\mapsto(G,n-k)$.

\npdato{\npSC} dado un conjunto $U$ de $n$ elementos, subconjuntos $S_1,\dots,S_m\subseteq U$ y $k$, ¿hay a lo sumo $k$ de los $S_i$ cuya unión es $U$? Certificado: los índices elegidos. Se prueba desde \npVC: $U=E$ y, para cada vértice $v$, $S_v=\{\text{aristas incidentes a }v\}$, con el mismo $k$. Un conjunto de vértices cubre todas las aristas si y solo si los $S_v$ correspondientes cubren $U$.

\npdato{\npSP} dado $U$, subconjuntos $S_1,\dots,S_m$ y $k$, ¿hay al menos $k$ de los $S_i$ disjuntos dos a dos? Se prueba desde \npIS\ con la misma construcción: $U=E$, $S_v$ las aristas incidentes a $v$. Dos $S_v$ son disjuntos si y solo si $u$ y $v$ no son adyacentes.

\npdato{\npCL} dado $G$ y $k$, ¿hay $k$ vértices adyacentes dos a dos? Se prueba desde \npIS\ con el grafo complemento $\overline{G}$ (misma $V$, arista donde $G$ no tiene) y el mismo $k$: $S$ es independiente en $G$ si y solo si es clique en $\overline{G}$. Construir $\overline G$ es $\Oh(n^2)$.

\npdato{\prob{Hitting Set}} dado $A$, subconjuntos $B_1,\dots,B_m$ de $A$ y $k$, ¿hay $H\subseteq A$ con $|H|\le k$ que toca a todos los $B_i$? Es el <<transpuesto>> de \npSC\ (se eligen elementos en lugar de conjuntos). Desde \npVC: $A=V$ y un $B_i=\{u,v\}$ por cada arista.

\subsection{Satisfacibilidad}

\npdato{\npSAT} dadas variables $x_1,\dots,x_n$ y cláusulas $C_1,\dots,C_k$ (cada una disyunción de términos $x_i$ o $\npneg{x_i}$), ¿hay una asignación que satisface todas las cláusulas? \npTSAT\ es el caso en que cada cláusula tiene exactamente tres términos sobre variables distintas. Certificado: la asignación. \npTSAT\ se prueba desde \prob{Circuit-sat} (Cook--Levin).

Una reducción que viene bien conocer porque es la base de los ejercicios de <<$k$-\prob{sat}>>: $\npSAT\redp\npTSAT$. Una cláusula larga $(t_1\vee t_2\vee\dots\vee t_r)$ se parte con variables nuevas $y_j$:
\[
(t_1\vee t_2\vee y_1)\wedge(\npneg{y_1}\vee t_3\vee y_2)\wedge\dots\wedge(\npneg{y_{r-3}}\vee t_{r-1}\vee t_r).
\]
Si algún $t_i$ es verdadero, se eligen los $y_j$ para satisfacer el resto de la cadena; si todos los $t_i$ son falsos, la cadena obliga $y_1=1$, luego $y_2=1$, \dots, y la última cláusula falla. Las cláusulas cortas se completan con variables nuevas forzadas.

\begin{problema}{3-SAT $\redp$ Independent Set (libro, sección 8.2)}
\npdato{Construcción} dada una instancia de \npTSAT\ con $k$ cláusulas, se arma un grafo con $3k$ vértices: un triángulo por cláusula, con un vértice por cada término. Además se une con una arista cada par de vértices que representan términos opuestos ($x_i$ en una cláusula y $\npneg{x_i}$ en otra). Se pregunta por un independiente de tamaño $k$.

\npdato{Ida} si hay asignación que satisface, se elige en cada triángulo un término verdadero. Son $k$ vértices, no hay dos del mismo triángulo y no hay dos opuestos (no pueden ser ambos verdaderos).

\npdato{Vuelta} un independiente de tamaño $k$ tiene exactamente un vértice por triángulo y no contiene opuestos. Se pone verdadero cada término elegido (no hay conflicto) y el resto de las variables arbitrarias: cada cláusula tiene un término verdadero.
\end{problema}

\subsection{Secuenciación}

\npdato{Ciclo hamiltoniano} dado un grafo (dirigido o no), ¿hay un ciclo simple que pasa por todos los vértices? Certificado: los vértices en orden. El libro lo prueba para grafos dirigidos desde \npTSAT; la versión no dirigida sale de la dirigida partiendo cada vértice $v$ en un camino $v_{\text{in}}-v_{\text{mid}}-v_{\text{out}}$.

\npdato{Camino hamiltoniano} igual, pero un camino simple (no hace falta volver). Desde ciclo hamiltoniano dirigido: se elige un vértice $v$ y se lo parte en $v'$, que se queda con las aristas salientes de $v$, y $v''$, que se queda con las entrantes. Como $v'$ no tiene entrantes y $v''$ no tiene salientes, todo camino hamiltoniano va de $v'$ a $v''$, y eso es lo mismo que un ciclo hamiltoniano por $v$ en el grafo original.

\npdato{\npTSP} dadas $n$ ciudades, distancias $d(u,v)$ y una cota $D$, ¿hay un recorrido que visita cada ciudad una vez, vuelve al inicio y tiene largo total $\le D$? Desde ciclo hamiltoniano: $d(u,v)=1$ si $(u,v)$ es arista y $2$ si no, con $D=n$. Un recorrido de largo $n$ usa solo aristas del grafo.

\begin{problema}{3-SAT $\redp$ ciclo hamiltoniano dirigido (libro, sección 8.5)}
\npdato{Gadget de variable} por cada variable $x_i$ una fila de nodos $v_{i,1},\dots,v_{i,b}$ ($b=3k+3$) con aristas en ambos sentidos entre vecinos. Recorrer la fila de izquierda a derecha codifica $x_i=1$; de derecha a izquierda, $x_i=0$. Los extremos de cada fila se conectan con los dos extremos de la fila siguiente, y se agregan un nodo inicial $s$, uno final $t$ y la arista $(t,s)$. Sin cláusulas, hay exactamente $2^n$ ciclos hamiltonianos, uno por asignación.

\npdato{Gadget de cláusula} por cada cláusula $C_j$ un nodo $c_j$. Si $x_i$ aparece en $C_j$ se agregan $v_{i,3j}\to c_j\to v_{i,3j+1}$ (un desvío que solo se puede tomar yendo de izquierda a derecha); si aparece $\npneg{x_i}$, se agregan $v_{i,3j+1}\to c_j\to v_{i,3j}$ (solo yendo de derecha a izquierda).

\npdato{Equivalencia} una asignación que satisface da un ciclo: se recorren las filas según la asignación y cada $c_j$ se visita por el desvío de un término verdadero. Recíprocamente, en un ciclo hamiltoniano cada $c_j$ se visita por un desvío que vuelve a la misma fila, y el sentido de esa fila hace verdadero al término correspondiente.
\end{problema}

\subsection{Partición: 3D-Matching}

\npdato{\npTDM} dados conjuntos disjuntos $X,Y,Z$ de tamaño $n$ cada uno y un conjunto $T\subseteq X\times Y\times Z$ de ternas, ¿hay $n$ ternas de $T$ que cubren cada elemento exactamente una vez? Es la generalización a tres dimensiones del matching bipartito perfecto (que sí está en $\Pclass$, por flujo). Certificado: las $n$ ternas.

\begin{problema}{3-SAT $\redp$ 3D-Matching (libro, sección 8.6), en una línea por gadget}
\npdato{Variable} un <<anillo>> de $2k$ elementos núcleo y $2k$ puntas, con ternas (núcleo$_j$, núcleo$_{j+1}$, punta$_j$). Para cubrir el núcleo hay que tomar todas las ternas pares o todas las impares: eso deja libres las puntas impares ($x_i=1$) o las pares ($x_i=0$).

\npdato{Cláusula} dos elementos núcleo nuevos y una terna con cada punta libre que corresponde a un término de la cláusula: solo se puede cubrir si algún término es verdadero.

\npdato{Limpieza} gadgets extra, cada uno con ternas con todas las puntas, para cubrir las puntas que sobran. Se reparten los elementos en $X$, $Y$, $Z$ de modo que cada terna tenga uno de cada uno.
\end{problema}

\subsection{Coloreo}

\npdato{$k$-coloreo} dado $G$, ¿se pueden asignar $k$ colores a los vértices de modo que vecinos tengan colores distintos? Para $k=2$ está en $\Pclass$ (es ver si el grafo es bipartito, con BFS). Para $k=3$ ya es $\NP$-completo. Certificado: el color de cada vértice. De $3$-coloreo a $k$-coloreo, $k>3$: se agregan $k-3$ vértices nuevos adyacentes a todos (y entre sí); se llevan $k-3$ colores propios y quedan 3 para el grafo original.

\begin{problema}{3-SAT $\redp$ 3-coloreo (libro, sección 8.7)}
\npdato{Paleta} un triángulo con nodos especiales $T$, $F$ y $B$ (base); sus colores representan verdadero, falso y <<otro>>.

\npdato{Variables} por cada $x_i$, dos nodos $v_i$ y $\npneg{v_i}$ unidos entre sí y con $B$. Así cada uno toma color de $T$ o de $F$, y son opuestos.

\npdato{Cláusulas} por cada cláusula, un subgrafo de seis nodos conectado a sus tres términos y a $T$ y $F$, armado para que se pueda colorear si y solo si alguno de los tres términos tiene el color de $T$. Si los tres son $F$, el nodo de salida del gadget queda forzado a $F$ y choca con una arista a $F$.
\end{problema}

\subsection{Problemas numéricos}

\npdato{\npSS} dados naturales $w_1,\dots,w_n$ y una meta $W$, ¿hay un subconjunto que suma exactamente $W$? Certificado: el subconjunto. Tiene un algoritmo $\Oh(nW)$ que no es polinomial (ver arriba).

\npdato{Mochila (decisión)} dados pesos $w_i$, valores $v_i$, capacidad $W$ y meta $V$, ¿hay un subconjunto de peso $\le W$ y valor $\ge V$? Desde \npSS: $v_i=w_i$ y $V=W$; peso $\le W$ y valor $\ge W$ obligan a sumar exactamente $W$.

\npdato{Secuenciación con release times y deadlines} $n$ trabajos en una sola máquina, el trabajo $i$ dura $t_i$, puede empezar desde $r_i$ y debe terminar antes de $d_i$; ¿hay una planificación sin solapamientos que cumple todo? Desde \npSS: un trabajo por cada $w_i$, con $r_i=0$, $t_i=w_i$, $d_i=S+1$ (donde $S=\sum w_i$), y un trabajo extra con $r=W$, $t=1$, $d=W+1$. El extra ocupa exactamente $[W,W+1]$, así que antes de él tiene que haber trabajos que llenan exactamente $W$ unidades.

\begin{problema}{3D-Matching $\redp$ Subset Sum (libro, sección 8.8)}
\npdato{Construcción} con $|X|=|Y|=|Z|=n$ y $m$ ternas, a cada terna $(x_i,y_j,z_k)$ se le asocia el número que en base $b=m+1$ tiene un 1 en las posiciones $i$, $n+j$ y $2n+k$ y 0 en las demás. La meta es $W=\sum_{r=0}^{3n-1}b^{r}$ (todos los dígitos en 1).

\npdato{Por qué funciona} como hay a lo sumo $m$ números y cada dígito vale 0 o 1, nunca hay acarreo en base $m+1$. Sumar $W$ es tener exactamente un 1 en cada posición, es decir, cubrir cada elemento exactamente una vez. Cada número tiene $3n$ dígitos en base $m+1$, o sea $\Oh(n\log m)$ bits: la transformación es polinomial.
\end{problema}

\subsection{El árbol de reducciones}

Cada flecha $A\to B$ significa $A\redp B$ en la prueba del libro: se transforman instancias de $A$ en instancias de $B$, y $B$ hereda la dificultad.

\begin{center}
\begin{tikzpicture}[
    nodo/.style={draw=acento,fill=acento!7,rounded corners=2pt,font=\sans\footnotesize,
                 align=center,inner sep=3pt,minimum height=0.62cm},
    raiz/.style={nodo,fill=acento!22},
    fl/.style={-{Stealth[length=2mm]},draw=gris,thick},
    x=1cm,y=1cm]
  \node[raiz] (csat) at (6.2,0)    {Circuit-SAT\\[-1pt]{\scriptsize Cook--Levin}};
  \node[raiz] (sat)  at (6.2,-1.5) {SAT $\equiv$ 3-SAT};
  \node[nodo] (is)   at (2.4,-3.1) {Independent Set};
  \node[nodo] (hc)   at (6.1,-3.1) {ciclo\\[-1pt]hamiltoniano};
  \node[nodo] (tdm)  at (8.9,-3.1) {3D-Matching};
  \node[nodo] (col)  at (11.6,-3.1){3-coloreo};
  \node[nodo] (vc)   at (0.75,-4.7){Vertex\\[-1pt]Cover};
  \node[nodo] (sp)   at (2.4,-4.7) {Set\\[-1pt]Packing};
  \node[nodo] (cl)   at (3.95,-4.7){Clique};
  \node[nodo] (hp)   at (5.45,-4.7){camino\\[-1pt]ham.};
  \node[nodo] (tsp)  at (6.9,-4.7) {TSP};
  \node[nodo] (ss)   at (8.9,-4.7) {Subset Sum};
  \node[nodo] (kcol) at (11.6,-4.7){$k$-coloreo};
  \node[nodo] (sc)   at (0.75,-6.3){Set\\[-1pt]Cover};
  \node[nodo] (hs)   at (2.4,-6.3) {Hitting\\[-1pt]Set};
  \node[nodo] (sch)  at (7.95,-6.3){secuenciación con\\[-1pt]release times};
  \node[nodo] (mo)   at (10.25,-6.3){mochila};
  \draw[fl] (csat) -- (sat);
  \draw[fl] (sat) -- (is);  \draw[fl] (sat) -- (hc);
  \draw[fl] (sat) -- (tdm); \draw[fl] (sat) -- (col);
  \draw[fl] (is) -- (vc); \draw[fl] (is) -- (sp); \draw[fl] (is) -- (cl);
  \draw[fl] (vc) -- (sc); \draw[fl] (vc) -- (hs);
  \draw[fl] (hc) -- (hp); \draw[fl] (hc) -- (tsp);
  \draw[fl] (tdm) -- (ss);
  \draw[fl] (ss) -- (sch); \draw[fl] (ss) -- (mo);
  \draw[fl] (col) -- (kcol);
\end{tikzpicture}
\end{center}

\npVC\ e \npIS\ se reducen en ambos sentidos; también \npSAT\ y \npTSAT. \prob{Hitting Set} y \npCL\ no están entre las reducciones que el libro desarrolla en el texto (son ejercicios estándar), pero son tan directas que conviene tenerlas en el árbol: de ahí cuelgan varios ejercicios de examen.

% =====================================================================
\section{Cómo se escribe una demostración de \texorpdfstring{$\NP$}{NP}-completitud}
% =====================================================================

Todas las soluciones oficiales tienen la misma estructura, y desde 2024 la pauta de corrección la explicita punto por punto. Si se sigue la plantilla, el puntaje está casi asegurado aunque la reducción sea la obvia.

\begin{receta}[demostrar que $X$ es $\NP$-completo]
\textbf{0. Plantear.} <<Para probar que $X$ es $\NP$-completo hay que probar (I) $X\in\NP$ y (II) $Y\redp X$ para todo $Y\in\NP$. Para (II) alcanza con probar $Z\redp X$ para un $Z$ $\NP$-completo conocido y usar la transitividad de $\redp$.>>

\textbf{1. $X\in\NP$.}
\begin{itemize}
  \item Describir la instancia $s$ y su representación (matriz de adyacencia, lista de cláusulas, arreglo\dots).
  \item Definir el certificado $t$ (casi siempre: la solución candidata, como arreglo de bits o lista de índices) y decir por qué $|t|\le p(|s|)$.
  \item Escribir el certificador $B(s,t)$ en pseudocódigo: una serie de chequeos que responden \npno\ si alguno falla y \npsi\ al final.
  \item Tiempo de $B$: polinomial en $|s|+|t|$ (contar los ciclos).
  \item Si $s$ es \npsi, existe $t$ (la solución misma) con $B(s,t)=\text{\npsi}$.
  \item Si $B(s,t)=\text{\npsi}$ para algún $t$, cada chequeo garantiza una condición del problema, así que $s$ es \npsi.
\end{itemize}

\textbf{2. $Z\redp X$.}
\begin{itemize}
  \item Tomar una instancia \emph{arbitraria} $s_Z$ del problema conocido $Z$ y describir la instancia $s_X=f(s_Z)$, componente por componente (<<para cada vértice $v_i$ se define\dots>>).
  \item Justificar que $f$ corre en tiempo polinomial en $|s_Z|$ (y por lo tanto $|s_X|$ es polinomial).
  \item Ida: si $s_Z$ es \npsi, construir explícitamente una solución de $s_X$ a partir de la de $s_Z$.
  \item Vuelta: si $s_X$ es \npsi, construir una solución de $s_Z$ a partir de la de $s_X$. (Equivalente: si $s_Z$ es \npno, $s_X$ es \npno.)
  \item Cerrar: <<se resuelve $Z$ con la transformación, que es polinomial, y una sola llamada a un algoritmo que resuelve $X$; por lo tanto $Z\redp X$>>.
\end{itemize}

\textbf{3. Concluir.} <<Como $Z$ es $\NP$-completo, todo $Y\in\NP$ cumple $Y\redp Z$; por transitividad $Y\redp X$. Junto con $X\in\NP$, $X$ es $\NP$-completo.>>
\end{receta}

Cuando la letra da una hipótesis, se usa y no se demuestra de nuevo: <<asuma que existe un certificador eficiente>> resuelve el paso 1 en una oración; <<asuma que si la original es \npsi\ la transformada es \npsi>> elimina la ida. La pauta de 2025 pide explícitamente \emph{no} describir el certificador en ese caso, sino decir que se asume.\visto{2P 2025 Ej. 2: usar las hipótesis dadas}

\begin{trampa}
\begin{itemize}
  \item \textbf{Reducir al revés.} Probar $X\redp Z$ (el nuevo al conocido) no prueba nada sobre la dificultad de $X$. La pauta de 2024 lo califica de <<error grave>>. Chequeo: la instancia que se toma \emph{arbitraria} es la del problema conocido.
  \item \textbf{Olvidar $X\in\NP$.} Sin eso se probó $\NP$-difícil, no $\NP$-completo.
  \item \textbf{Probar una sola dirección.} La ida suele ser obvia; la vuelta es donde está el ejercicio (\prob{Cometa}, \prob{nti4sat}, \prob{casi-sat}).
  \item \textbf{Vuelta sobre instancias arbitrarias de $X$.} La vuelta se prueba solo para instancias de la forma $f(s_Z)$; se puede (y se debe) usar cómo fueron construidas.
  \item \textbf{Transformación que usa la respuesta.} $f$ no puede depender de conocer la solución de $s_Z$ (por ejemplo, <<se toma $k$ igual al tamaño del máximo independiente>>), porque calcularla no es polinomial.
  \item \textbf{Construcción no polinomial.} Escribir números en unario, enumerar todos los subconjuntos o crear un objeto por cada asignación de verdad.
  \item \textbf{Instancias inválidas.} La instancia construida tiene que respetar la definición de $X$: en \prob{4-sat} los cuatro términos son sobre variables distintas, en \prob{Red-Social} la amistad es simétrica, en \prob{Cometa} $d\ge 3$.
  \item \textbf{Certificador que busca.} El certificador verifica la solución que recibe; si para responder tiene que buscar entre muchas opciones, no es polinomial.
  \item \textbf{Certificado sin tamaño.} Hay que decir por qué $|t|$ es polinomial en $|s|$, no solo describirlo.
  \item \textbf{Olvidar el tiempo de la transformación.} Aunque sea una copia, hay que decir que es lineal o polinomial.
\end{itemize}
\end{trampa}

\subsection{Cómo elegir el problema de origen}

La sugerencia de la letra suele decir de dónde reducir. Cuando no la hay, se busca el problema del catálogo que se parece más <<estructuralmente>> al nuevo.

\begin{center}\small
\begin{tabularx}{\textwidth}{@{}>{\raggedright\arraybackslash}X >{\raggedright\arraybackslash}p{3.6cm}@{}}
\toprule
Si el problema nuevo\dots & conviene reducir desde \\
\midrule
elige elementos de un grafo (o de una matriz de <<conflictos>>) que no pueden convivir, y pide al menos $k$ & \npIS \\
pide al menos $k$ elementos todos compatibles entre sí, o un subgrafo completo & \npCL\ (o \npIS\ con complemento) \\
pide a lo sumo $k$ objetos que <<tocan>> todas las restricciones de a pares & \npVC \\
pide a lo sumo $k$ grupos cuya unión cubre todo (o $k$ elementos que tocan todos los grupos) & \npSC\ o \npVC \\
tiene variables booleanas, cláusulas, apuestas o tableros con dos estados por fila & \npSAT\ o \npTSAT \\
pide recorrer todo una vez, un camino o ciclo largo, un orden & camino o ciclo hamiltoniano \\
pide sumar exactamente un número, repartir en dos mitades, cargar una capacidad & \npSS \\
pide asignar etiquetas a vértices con restricciones de vecinos & 3-coloreo \\
pide particionar en ternas & \npTDM \\
\bottomrule
\end{tabularx}
\end{center}

Una observación que ahorra tiempo: casi todos los ejercicios de examen son \emph{disfraces} del problema de referencia, y la transformación es un renombre (vértices $\to$ ubicaciones, aristas $\to$ tareas, cláusulas $\to$ columnas). Si la transformación que se está armando es complicada, probablemente se eligió mal el origen.

\subsection{Los problemas de examen, de un vistazo}

La figura muestra de qué problema conocido se redujo cada ejercicio en la solución oficial.

\begin{center}
\begin{tikzpicture}[
    src/.style={draw=acento,fill=acento!14,rounded corners=2pt,font=\sans\footnotesize\bfseries,
                align=center,inner sep=3pt,minimum height=0.6cm,minimum width=2.6cm},
    hoja/.style={draw=grisclaro,fill=white,font=\sans\scriptsize,align=left,
                 text width=2.65cm,inner sep=3pt,anchor=north},
    nada/.style={hoja,draw=gris,dashed,text=gris},
    fl/.style={-{Stealth[length=2mm]},draw=gris,thick},
    x=1cm,y=1cm]
  \node[src] (is) at (1.6,0)  {Independent Set};
  \node[src] (cl) at (4.8,0)  {Clique};
  \node[src] (vc) at (8.0,0)  {Vertex Cover};
  \node[src] (sc) at (11.2,0) {Set Cover};
  \node[hoja] (his) at (1.6,-0.9) {UOPI \hfill 2P 2018\\ Red-Social \hfill feb.\,2021\\ Cocinero \hfill jul.\,2022\\ RF \hfill dic.\,2022};
  \node[hoja] (hcl) at (4.8,-0.9) {Cometa \hfill dic.\,2023\\ Subgrafo isom. \hfill jul.\,2024};
  \node[hoja] (hvc) at (8.0,-0.9) {PCG \hfill feb.\,2020\\ Test \hfill jul.\,2021\\ Contratar \hfill 2P 2022};
  \node[hoja] (hsc) at (11.2,-0.9){Acceso \hfill 2P 2023\\ UV \hfill 2P 2025};
  \node[src] (sat) at (1.6,-3.6)  {SAT};
  \node[src] (tsat) at (4.8,-3.6) {3-SAT};
  \node[src] (hp) at (8.0,-3.6)   {camino ham.};
  \node[src,draw=gris,fill=grisclaro!60] (otros) at (11.2,-3.6) {sin ejercicios};
  \node[hoja] (hsat) at (1.6,-4.5) {casiSAT \hfill 2P 2020\\ Falsable \hfill feb.\,2022\\ Al menos $k$ \hfill jul.\,2025};
  \node[hoja] (htsat) at (4.8,-4.5){Solitaire \hfill 2P 2019\\ 4-SAT \hfill 2P 2024\\ NTI4Sat \hfill jul.\,2026};
  \node[hoja] (hhp) at (8.0,-4.5)  {camino simple\\ de largo $\ge k$ \hfill 2P 2021};
  \node[nada] (hot) at (11.2,-4.5) {Subset Sum, 3-coloreo, 3D-Matching, TSP};
  \foreach \a/\b in {is/his,cl/hcl,vc/hvc,sc/hsc,sat/hsat,tsat/htsat,hp/hhp}
    \draw[fl] (\a) -- (\b);
\end{tikzpicture}
\end{center}

% =====================================================================
\section{Problemas de examen, agrupados por problema de origen}
% =====================================================================

Para cada ejercicio: enunciado resumido, cómo se prueba que está en $\NP$, la transformación y los argumentos de ida y vuelta. Los certificadores de los ejercicios son todos del mismo tipo (recorrer la solución candidata y chequear las restricciones); se detallan en los primeros y después se resumen.

\subsection{Desde Independent Set}

La señal es <<elegir al menos $k$ cosas sin conflictos de a pares>>. La transformación casi siempre deja el grafo igual (o lo codifica en una matriz) y deja $k$ igual.

El primero es el más directo de todos: \npIS\ con pesos. La solución oficial es un buen modelo de redacción del certificador, con las dos direcciones explícitas.\visto{2P 2018 Ej. 3: UOPI, desde IS}

\begin{problema}{Ubicación óptima de plantas industriales (UOPI)}
\npdato{Problema} grafo $G$ cuyos vértices son ubicaciones con valor entero $w_i$; una arista indica que dos ubicaciones son excluyentes. Dado $t$, ¿hay ubicaciones no excluyentes entre sí cuyo valor total es al menos $t$?

\npdato{En $\NP$} certificado: tira de $n$ bits $c$ (ubicaciones elegidas). $B$ suma los $w_i$ con $c_i=1$ y rechaza si la suma es menor a $t$; después, para cada $i$ con $c_i=1$, recorre la lista de adyacencia de $i$ y rechaza si encuentra un $j$ con $c_j=1$. Es $\Oh(n+m)$. Si $B$ acepta, la selección suma al menos $t$ y no tiene pares excluyentes; si la instancia es \npsi, la tira de esa selección hace aceptar a $B$.

\npdato{Reducción $\npIS\redp\prob{uopi}$} $(G,k)\mapsto(G,\,w_i=1\ \forall i,\,t=k)$, en $\Oh(n)$. Con pesos 1, la suma de un conjunto es su tamaño; un conjunto es independiente si y solo si no tiene exclusiones. Así, hay independiente de tamaño $\ge k$ si y solo si hay selección válida de valor $\ge t$.
\end{problema}

El segundo es \npCL\ con otro nombre, y la reducción desde \npIS\ es el grafo complemento: <<amigos>> son los no adyacentes.\visto{Examen feb. 2021 Ej. 3: Red-Social, desde IS}

\begin{problema}{Red-Social}
\npdato{Problema} personas $P=\{p_1,\dots,p_n\}$ con conjuntos de amigos $A_i$ (relación simétrica) y un natural $k$: ¿hay al menos $k$ personas amigas todas entre sí?

\npdato{Reducción $\npIS\redp\prob{Red-Social}$} dada $(G,k)$: $P=V$; $A_i=\{v_j : j\ne i,\ v_j \text{ no adyacente a } v_i\}$; $k'=k$. La relación es simétrica porque la adyacencia lo es. Tiempo $\Oh(n^3)$ con una búsqueda ingenua (o $\Oh(n^2)$ con matriz).

\npdato{Ida y vuelta} $S$ independiente en $G$ $\iff$ cada par de $S$ es no adyacente $\iff$ cada par de $S$ es amigo en $R$. El tamaño se conserva porque $k'=k$.

\npdato{En $\NP$} certificado: arreglo booleano de $n$ posiciones. $B$ rechaza si hay menos de $k$ elegidos y, para cada par $(p,q)$ de elegidos con $p\ne q$, rechaza si $q\notin A_p$. Tiempo $\Oh(n^3)$.
\end{problema}

En \prob{Cocinero} el problema nuevo tiene un parámetro más (la disconformidad tolerada $d$); se fija $d=0$ y queda exactamente \npIS. Es el truco de <<restringir>>: si un caso particular de $X$ ya es $\NP$-completo, $X$ también lo es.\visto{Examen jul. 2022 Ej. 3: Cocinero, desde IS}

\begin{problema}{Cocinero}
\npdato{Problema} matriz simétrica $M$ de $n\times n$ con $M[i,j]=1$ si los ingredientes $i\ne j$ son incompatibles. La disconformidad de un plato $P$ es $\sum_{i,j\in P,\,i<j}M[i,j]$. Dados $k$ y $d$, ¿hay un plato con al menos $k$ ingredientes y disconformidad a lo sumo $d$?

\npdato{En $\NP$} certificado: el conjunto de ingredientes $t$, $|t|=\Oh(|s|)$. $B$ chequea $|t|\ge k$ y calcula la suma sobre los pares de $t$ en $\Oh(|t|^2)$; acepta si es $\le d$.

\npdato{Reducción $\npIS\redp\prob{Cocinero}$} $(G,k)\mapsto(M,k,0)$ con $M$ la matriz de adyacencia de $G$, en $\Oh(n^2)$. Disconformidad 0 significa que no hay pares incompatibles en el plato, es decir, un conjunto independiente.
\end{problema}

En \prob{RF} la dificultad está en traducir: cada arista pasa a ser un elemento de $X$ relacionado con sus dos extremos, y <<función parcial>> pasa a significar <<cada arista tiene a lo sumo un extremo elegido>>.\visto{Examen dic. 2022 Ej. 3: RF, desde IS}

\begin{problema}{Relación a función (RF)}
\npdato{Problema} una relación $R\subseteq X\times Y$ dada como matriz de ceros y unos (filas $Y$, columnas $X$) y un entero $k$: ¿hay un subconjunto $U\subseteq Y$ de al menos $k$ elementos tal que la restricción de $R$ a $U$ es una función parcial, o sea, cada $x\in X$ se relaciona con a lo sumo un elemento de $U$?

\npdato{En $\NP$} certificado: tira de $|Y|$ bits. $B$ chequea que tenga al menos $k$ unos y, columna por columna, que haya a lo sumo un 1 entre las filas elegidas. Tiempo $\Oh(|X|\,|Y|)$.

\npdato{Reducción $\npIS\redp\prob{rf}$} $Y=V$, $X=E$, y por cada arista $e=(u,w)$ se ponen en $R$ los pares $(e,u)$ y $(e,w)$; $k'=k$. Polinomial: la matriz tiene $n\cdot m$ celdas.

\npdato{Ida} si $U$ es independiente, cada arista tiene a lo sumo un extremo en $U$, así que cada $x=e$ se relaciona con a lo sumo un elemento de $U$.

\npdato{Vuelta (contrarrecíproco)} si $(G,k)$ es \npno, todo $U$ de $k$ vértices contiene una arista $(u,w)$, y entonces $x=(u,w)$ se relaciona con dos elementos de $U$: ningún $U$ sirve y $(R,k)$ es \npno.
\end{problema}

\subsection{Desde Clique}

Cuando lo que se busca es un subgrafo completo (o algo que contiene uno), conviene arrancar directamente de \npCL. La letra lo daba como sugerencia en los dos casos.

La reducción de \prob{Subgrafo Isomorfo} es la más corta del relevamiento: el grafo patrón es $K_k$. Lo interesante es el certificador, que recibe la biyección.\visto{Examen jul. 2024 Ej. 3: subgrafo isomorfo}

\begin{problema}{Subgrafo isomorfo}
\npdato{Problema} dados dos grafos $A$ y $G$ por matrices de adyacencia, ¿hay un subgrafo $B$ de $G$ isomorfo a $A$? ($A$ y $B$ son isomorfos si hay una biyección $f$ que preserva la adyacencia en ambos sentidos.)

\npdato{En $\NP$} certificado: lista de pares $(a,v)\in V_A\times V_G$ que define $f$. $B$ chequea que la lista tenga $|V_A|$ pares, que no se repitan vértices de $A$ ni de $G$ (así $f$ es función e inyectiva), y para cada par $(a,a')$ que $M_A[a,a']=M_G[f(a),f(a')]$. Tiempo $\Oh(|V_A|^2)$ más las búsquedas en la lista.

\npdato{Reducción $\npCL\redp\prob{Subgrafo Isomorfo}$} dada $(G,k)$ con $k\le|V|$, se construye $K$, el grafo completo de $k$ vértices, y se pregunta por $(K,G)$. Es $\Oh(k^2)=\Oh(|V|^2)$. Un clique de tamaño $k$ es un subgrafo isomorfo a $K$ y viceversa.
\end{problema}

\prob{Cometa} es el ejercicio con más trampa del relevamiento: la construcción natural (un solo camino colgado) no funciona, y la solución oficial lo dice explícitamente. Además hay que tratar aparte el caso $k>n$.\visto{Examen dic. 2023 Ej. 3: Cometa, desde Clique}

\begin{problema}{Cometa}
\npdato{Problema} un grafo de $2d$ vértices ($d\ge3$) es una cometa si $d$ vértices forman un clique y los otros $d$ forman un camino simple con un extremo conectado a un único vértice del clique. Dados $G$ (matriz de adyacencia) y $d\ge 3$, ¿hay un subgrafo de $G$ que es una cometa con al menos $2d$ vértices? Se sugiere \prob{Clique3}: \npCL\ con $k\ge 3$.

\npdato{En $\NP$} certificado: lista de $2d$ vértices, con el clique primero (el vértice de unión en la posición $d$) y el camino después. $B$ chequea que sean distintos, que los primeros $d$ sean adyacentes dos a dos y que cada par consecutivo desde la posición $d$ hasta la $2d$ sea adyacente. Son $\Oh(d^2)$ consultas a la matriz.

\npdato{Reducción $\prob{Clique3}\redp\prob{Cometa}$} si $k>n$, se responde \npno\ (o se devuelve una instancia \npno\ fija). Si no, $d=k$ y $G'$ es $G$ más, \emph{por cada vértice $v$}, un camino nuevo de $d$ vértices con un extremo unido a $v$ (caminos disjuntos). $G'$ tiene $n(k+1)$ vértices: polinomial porque $k\le n$.

\npdato{Ida} el clique de tamaño $k$ junto con el camino colgado de cualquiera de sus vértices es una cometa de $2d$ vértices.

\npdato{Vuelta} un vértice de un camino agregado tiene grado $\le 2$ y sus dos vecinos no son adyacentes, así que no está en ningún clique de tamaño $\ge3$. Como $d\ge3$, el clique de la cometa está entero dentro de $G$: $G$ tiene un clique de tamaño $k$.

\npdato{Lo que no sirve} colgar un único camino con un extremo adyacente a \emph{todos} los vértices de $G$: ese extremo junto con un clique de tamaño $d-1$ de $G$ forma un clique de tamaño $d$, y una instancia \npno\ de \npCL\ podría dar una cometa.
\end{problema}

\subsection{Desde Vertex Cover y Set Cover}

La señal es <<a lo sumo $k$ cosas que cubren todo>>. Hay dos formas del mismo problema y conviene distinguirlas para armar la transformación sin dudar:

\begin{itemize}
  \item \emph{forma \npSC}: se eligen \emph{grupos} y se quiere cubrir todos los \emph{elementos} (\prob{Test}: se eligen tests, se cubren líneas; \prob{Acceso}: se eligen lugares, se cubren dispositivos; \prob{uv}: se eligen caminos, se cubren vértices);
  \item \emph{forma \prob{Hitting Set}}: se eligen \emph{elementos} y se quiere tocar todos los \emph{grupos} (\prob{pcg}; \prob{Contratar}: se eligen aplicantes, cada tarea tiene su grupo de calificados).
\end{itemize}

Desde \npVC\ las dos son inmediatas: en la primera, cada vértice es un grupo (sus aristas incidentes) y cada arista es un elemento; en la segunda, cada vértice es un elemento y cada arista es un grupo de dos.

\prob{pcg} es literalmente \prob{Hitting Set}. La parte (c) solo pedía encadenar: $\npVC$ es $\NP$-completo, $\npVC\redp\prob{pcg}$ y se asume $\prob{pcg}\in\NP$.\visto{Examen feb. 2020 Ej. 1: PCG, desde VC}

\begin{problema}{Conjunto golpeador (PCG)}
\npdato{Problema} conjunto $A$, subconjuntos $B_1,\dots,B_m$ y $k$: ¿hay $S\subseteq A$ con $|S|\le k$ que contiene al menos un elemento de cada $B_i$?

\npdato{Reducción $\npVC\redp\prob{pcg}$} $A=V$; por cada arista $\{u,v\}$ un conjunto $B_i=\{u,v\}$; $k'=k$. Con listas de adyacencia, una pasada: para cada $v$ y cada vecino $w>v$ se agrega $\{v,w\}$ (la condición $w>v$ evita duplicados). Tiempo $\Oh(n+m)$.

\npdato{Ida y vuelta} $S$ golpea a todos los $B_i$ $\iff$ toda arista tiene un extremo en $S$ $\iff$ $S$ es un cubrimiento. El tamaño se conserva.
\end{problema}

En \prob{Test} los tests son vértices y las líneas son aristas; el test $v_i$ ejecuta las líneas de las aristas incidentes a $v_i$. Es la reducción del libro $\npVC\redp\npSC$ con otro vocabulario.\visto{Examen jul. 2021 Ej. 3: Test, desde VC}

\begin{problema}{Test}
\npdato{Problema} un programa de $N$ líneas y $M$ tests; el test $i$ ejecuta las líneas $L_i\subseteq\{1,\dots,N\}$. Dado $k$, ¿hay a lo sumo $k$ tests que ejecutan cada línea al menos una vez?

\npdato{Reducción $\npVC\redp\prob{Test}$} $N=m$ (una línea por arista), $M=n$ (un test por vértice), $L_i=\{j : e_j \text{ incide en } v_i\}$, $k'=k$. Tiempo $\Oh(nm)$.

\npdato{Ida} si $S$ cubre, los tests de $S$ ejecutan todas las líneas: la línea $j$ es la arista $e_j$, que tiene un extremo $v_{i_q}\in S$, y $j\in L_{i_q}$. \npdato{Vuelta} si $T$ ejecuta todas las líneas, para cada arista $e_j$ hay un test $i_q\in T$ con $j\in L_{i_q}$, o sea $e_j$ incide en $v_{i_q}$.

\npdato{En $\NP$} certificado: arreglo booleano de $M$ posiciones. $B$ chequea que haya a lo sumo $k$ elegidos y marca las líneas ejecutadas por los elegidos; acepta si quedan todas marcadas.
\end{problema}

\prob{Contratar} es la forma \prob{Hitting Set}: aplicantes $=$ vértices, tareas $=$ aristas, y los calificados para una tarea son los dos extremos de su arista. La letra sugería \npVC.\visto{2P 2022 Ej. 3: Contratar, desde VC}

\begin{problema}{Contratar}
\npdato{Problema} aplicantes $A=\{a_1,\dots,a_n\}$, tareas $t_1,\dots,t_m$ y, para cada tarea, el conjunto $C_i\subseteq A$ de calificados. Dado $k$, ¿se pueden contratar a lo sumo $k$ aplicantes de modo que cada tarea tenga al menos un calificado contratado?

\npdato{En $\NP$} certificado: arreglo $Z$ de $n$ bits. $B$ rechaza si $Z$ tiene más de $k$ unos; usa un arreglo $T$ de $m$ booleanos en \textit{false} y, para cada tarea $t_j$ y cada $a_i\in C_j$ con $Z[i]=1$, marca $T[j]$; acepta si todo $T$ quedó en \textit{true}. Tiempo $\Oh(nm)$.

\npdato{Reducción $\npVC\redp\prob{Contratar}$} un aplicante $a_i$ por vértice $v_i$, una tarea $t_i$ por arista $e_i=(v_j,v_k)$ con $C_i=\{a_j,a_k\}$, y $k'=k$. Tiempo $\Oh(n+m)$.

\npdato{Ida y vuelta} $S$ cubre la arista $e_j$ $\iff$ algún extremo de $e_j$ está en $S$ $\iff$ algún calificado para $t_j$ está contratado.
\end{problema}

\prob{Acceso} se reduce desde \npSC\ con la matriz de anchos de banda como matriz de pertenencia: $B(d,a_i)=1$ si $d\in S_i$ y 0 si no.\visto{2P 2023 Ej. 3: Acceso, desde Set Cover}

\begin{problema}{Acceso}
\npdato{Problema} dispositivos $D=\{d_1,\dots,d_n\}$, lugares $A=\{a_1,\dots,a_\ell\}$, una matriz $B$ de $n\times\ell$ con anchos de banda no negativos y un natural $k$. Un conjunto de lugares $S$ cubre a $d$ si $B(d,a)>0$ para algún $a\in S$. ¿Se puede cubrir todo $D$ con a lo sumo $k$ lugares?

\npdato{En $\NP$} certificado: vector de a lo sumo $k$ índices de lugares (más chico que la matriz). $B$ recorre las filas de la matriz y chequea que cada dispositivo tenga $B(d,a_i)>0$ para algún $a_i$ del certificado. Tiempo $\Oh(nk)$.

\npdato{Reducción $\npSC\redp\prob{Acceso}$} dada $(U,S_1,\dots,S_m,r)$: $D=U$, $\ell=m$, $k=r$, $B(d,a_i)=1$ si $d\in S_i$ y $0$ si no. Tiempo $\Oh(|U|\cdot m)$.

\npdato{Ida y vuelta} una subcolección de a lo sumo $r$ conjuntos cuya unión es $U$ corresponde exactamente a a lo sumo $k$ lugares que cubren todos los dispositivos.
\end{problema}

En \prob{uv} la letra guiaba paso a paso y pedía explícitamente identificar cuál problema es $X$ y cuál es $Y$ en $Y\redp X$: el origen $Y$ es \npSC, el destino $X$ es \prob{uv}. También daba como hipótesis la ida y la pertenencia a $\NP$.\visto{2P 2025 Ej. 2: UV, desde Set Cover}

\begin{problema}{Unión es V (UV)}
\npdato{Problema} un grafo $G=(V,E)$, un conjunto $\mathcal P=\{P_1,\dots,P_\ell\}$ de caminos simples de $G$ y un entero $k$: ¿hay $k$ caminos de $\mathcal P$ cuya unión de vértices es todo $V$?

\npdato{Reducción $\npSC\redp\prob{uv}$} dada $(U,\{S_1,\dots,S_\ell\},k)$: $V$ tiene un vértice $v_i$ por cada elemento $e_i\in U$; $G$ es el grafo completo sobre $V$ (así cualquier orden de vértices es un camino); por cada $S_i=\{e_{i_1},\dots,e_{i_h}\}$ se crea el camino $P_i=(v_{i_1},\dots,v_{i_h})$; el $k$ es el mismo. Si no se quiere el grafo completo, hay que agregar a mano las aristas de cada $P_i$. Es una copia: tiempo lineal en la instancia (más $\Oh(|U|^2)$ si se arma el completo).

\npdato{Vuelta (lo que se pedía)} si $k$ caminos $P_{i_1},\dots,P_{i_k}$ cubren $V$, para cada $e\in U$ su vértice $v$ está en algún $P_{i_j}$ y por construcción $e\in S_{i_j}$: los $S_{i_j}$ cubren $U$.

\npdato{Detalle} si \npSC\ pide <<a lo sumo $k$>> y \prob{uv} pide exactamente $k$, se supone $k\le\ell$ y, si hiciera falta, se agregan caminos de relleno; la solución oficial lo menciona en una nota al pie.
\end{problema}

\subsection{Desde SAT y 3-SAT}

La señal es que haya algo con dos estados por <<variable>> (fila, evento, ficha) y restricciones del tipo <<al menos uno de estos>>. La transformación pone una variable por cada objeto binario y una cláusula por cada restricción. Cuando el problema nuevo permite cláusulas de cualquier largo, se reduce desde \npSAT; cuando exige largo fijo, desde \npTSAT.

\prob{Solitaire} venía con la transformación dada: había que probar la vuelta, armar la reducción y concluir.\visto{2P 2019 Ej. 3: Solitaire, desde 3-SAT}

\begin{problema}{Solitaire}
\npdato{Problema} un tablero $n\times m$ con celdas vacías o con una ficha $X$ u $O$. Se gana si se puede quitar fichas de modo que en cada fila las que quedan sean todas del mismo tipo y cada columna tenga al menos una ficha.

\npdato{Transformación (dada)} desde una instancia de \npTSAT\ con variables $x_1..x_n$ y cláusulas $C_1..C_m$: tablero $n\times m$; si $x_i$ aparece en $C_j$ se pone $X$ en $(i,j)$, si aparece $\npneg{x_i}$ se pone $O$. Se recorren las cláusulas llenando una columna por cláusula: lineal.

\npdato{Vuelta} dada una configuración ganadora, $v(x_i)=1$ si la fila $i$ quedó con fichas $X$, $v(x_i)=0$ si quedó con $O$, arbitrario si quedó vacía (no hay ambigüedad porque cada fila es de un solo tipo). Cada columna $C_j$ conserva alguna ficha: si es $X$ en la fila $i$, el término $x_i$ es verdadero; si es $O$, el término $\npneg{x_i}$ es verdadero. Toda cláusula se satisface.

\npdato{Ida (se asumía)} con una asignación que satisface, en la fila $i$ se dejan solo las fichas del tipo que coincide con $v(x_i)$; cada columna conserva la ficha de su término verdadero.

\npdato{En $\NP$} certificado: lista de posiciones que se conservan. $B$ chequea que cada posición tenga ficha, que cada columna aparezca y que no haya dos fichas de distinto tipo en la misma fila. Con $R$ fichas, $\Oh(R^2)$ más los accesos al tablero.
\end{problema}

En \prob{casi-sat} se pide satisfacer todas menos una. El truco es agregar dos cláusulas contradictorias: exactamente una de ellas falla siempre, y por lo tanto las originales tienen que satisfacerse todas.\visto{2P 2020 Ej. 3: casiSAT, desde SAT}

\begin{problema}{casiSAT}
\npdato{Problema} variables $X$ y $k$ cláusulas: ¿hay una asignación que satisface al menos $k-1$ cláusulas?

\npdato{En $\NP$} certificado: la asignación. $B$ cuenta las cláusulas satisfechas recorriendo cada término una vez; acepta si son al menos $k-1$.

\npdato{Reducción $\npSAT\redp\prob{casi-sat}$} dada una instancia con variables $x_1..x_m$ y $r$ cláusulas, se agrega una variable nueva $z$ y dos cláusulas de un término: $(z)$ y $(\npneg z)$. Quedan $k=r+2$ cláusulas.

\npdato{Ida} la asignación que satisface las $r$ originales, con $z=1$, satisface $r+1=k-1$ cláusulas. \npdato{Vuelta} cualquier asignación falla en $(z)$ o en $(\npneg z)$; si satisface $k-1$, falla exactamente en una de esas dos y satisface las $r$ originales.
\end{problema}

\prob{Falsable} es el dual de \npSAT: una disyunción de conjunciones (forma normal disyuntiva) es falsa cuando todas las conjunciones son falsas, o sea, cuando cada una tiene un término falso.\visto{Examen feb. 2022 Ej. 2: Falsable, desde SAT}

\begin{problema}{Falsable}
\npdato{Problema} conjunciones $C_1,\dots,C_m$ (cada una $t_1\wedge\dots\wedge t_k$ con términos $x_i$ o $\npneg{x_i}$): ¿hay una asignación que hace falsa a $F=\bigvee_{i}C_i$?

\npdato{En $\NP$} certificado: la asignación. $B$ chequea que cada conjunción tenga algún término falso. Lineal.

\npdato{Reducción $\npSAT\redp\prob{Falsable}$} se niega la fórmula de \npSAT\ y se aplica De Morgan: cada cláusula $(t_1\vee\dots\vee t_r)$ pasa a la conjunción $(\npneg{t_1}\wedge\dots\wedge\npneg{t_r})$. Es cambiar el signo de cada término: lineal.

\npdato{Equivalencia} con la misma asignación, una cláusula original es verdadera si y solo si su conjunción negada es falsa. Todas las cláusulas son verdaderas si y solo si todas las conjunciones son falsas, es decir, $F$ es falsa.
\end{problema}

\prob{4-sat} parece trivial y no lo es: no se puede repetir una variable en una cláusula, así que <<agregar un cuarto término cualquiera>> no funciona si no se controla su valor. La solución oficial da dos reducciones.\visto{2P 2024 Ej. 2: 4-SAT, desde 3-SAT}

\begin{problema}{4-SAT}
\npdato{Problema} cláusulas de exactamente 4 términos sobre 4 variables distintas, representadas con matrices $V$ (índices) y $N$ (negaciones) de $k\times4$: ¿hay asignación que las satisface todas?

\npdato{En $\NP$} certificado: $n$ booleanos. $B$ evalúa las $k$ cláusulas en $\Oh(k)$. Como cada variable aparece en alguna cláusula, $n\le 4k$ y el certificado no es más largo que $N$.

\npdato{Reducción 1} se agregan variables nuevas $z,u,v,w$. Cada cláusula $C_i$ de \npTSAT\ pasa a $C_i\vee z$. Se agregan 8 cláusulas $(\ell_u\vee\ell_v\vee\ell_w\vee\npneg z)$, una por cada combinación de signos de $\ell_u\in\{u,\npneg u\}$, $\ell_v$, $\ell_w$. Para cualquier valor de $u,v,w$, una de esas 8 tiene sus tres primeros términos falsos, así que fuerzan $z=0$. Ida: misma asignación con $z=0$ y $u,v,w$ arbitrarios. Vuelta: $z=0$ forzado, luego cada $C_i\vee z$ se satisface por un término original.

\npdato{Reducción 2} cada cláusula $C_i$ se duplica: $C_i\vee x_{n+1}$ y $C_i\vee\npneg{x_{n+1}}$, con $x_{n+1}$ nueva ($2k$ cláusulas). Ida: misma asignación y $x_{n+1}$ arbitrario. Vuelta: cualquiera sea $x_{n+1}$, en una de las dos copias el cuarto término es falso y esa copia se satisface por un término de $C_i$.

\npdato{Tiempo} en las dos, agregar una columna a $V$ y $N$ y unas filas: lineal.
\end{problema}

En \prob{Al menos $k$} hay $k+1$ apostadores y se pide que acierten al menos $k$: otra vez <<todos menos uno>>, y otra vez el truco es agregar dos apuestas opuestas de las que exactamente una acierta.\visto{Examen jul. 2025 Ej. 3: Al menos k, desde SAT}

\begin{problema}{Al menos $k$}
\npdato{Problema} $n$ eventos (ganar $G_i$ o perder $P_i$) y $k+1$ apostadores, cada uno con apuestas sobre algunos eventos. ¿Hay resultados con los que al menos $k$ apostadores aciertan alguna apuesta? Se asume que hay un algoritmo polinomial que, dados la instancia y unos resultados, verifica si al menos $k$ aciertan.

\npdato{En $\NP$} el algoritmo dado es el certificador; el certificado es el conjunto de resultados, polinomial porque cada resultado corresponde a un evento con alguna apuesta. Hay que decir explícitamente las dos direcciones.

\npdato{Reducción $\npSAT\redp\prob{Al menos }k$} dada una instancia con $n$ variables y $c$ cláusulas: un evento por variable; un apostador por cláusula, con $G_i$ por cada término $x_i$ y $P_j$ por cada $\npneg{x_j}$; dos apostadores extra con $(G_n)$ y $(P_n)$. Hay $c+2$ apostadores y el umbral es $c+1$.

\npdato{Equivalencia} exactamente uno de los dos extra acierta. Por lo tanto aciertan al menos $c+1$ si y solo si aciertan los $c$ de las cláusulas. Con la correspondencia $x_i=1\iff$ se gana el evento $i$, el apostador de $C_h$ acierta si y solo si $C_h$ se satisface.

\npdato{Lo que no sirve} un apostador por cláusula sin los dos extra (umbral $c-1$): una instancia \npno\ de \npSAT\ que satisface $c-1$ cláusulas daría una instancia \npsi.
\end{problema}

\prob{nti4sat} (<<no todos iguales>>) pide en cada cláusula un término verdadero y uno falso. La reducción es agregar la misma variable nueva $s$ a todas las cláusulas; la clave de la vuelta es la simetría por complemento, que la letra daba como nota.\visto{Examen jul. 2026 Ej. 3: NTI4Sat, desde 3-SAT}

\begin{problema}{No todos iguales 4-SAT (NTI4Sat)}
\npdato{Problema} cláusulas de exactamente 4 términos sobre variables distintas. Una asignación satisface una cláusula si le da valor verdadero a algún término y falso a algún término. ¿Hay una asignación que satisface todas? Se asume que está en $\NP$.

\npdato{Reducción $\npTSAT\redp\prob{nti4sat}$} con una variable nueva $s$, cada cláusula $t_1\vee t_2\vee t_3$ pasa a $t_1\vee t_2\vee t_3\vee s$. Lineal.

\npdato{Ida} la asignación que satisface la instancia de \npTSAT, con $s=0$: cada cláusula tiene un término verdadero (de los originales) y uno falso ($s$).

\npdato{Vuelta} si una asignación satisface todas en el sentido <<no todos iguales>>, su complemento también (los verdaderos y los falsos se intercambian). Entre las dos, una tiene $s=0$; en esa, cada cláusula necesita un término verdadero, que tiene que ser uno de los tres originales. Restringida a las variables originales, satisface la instancia de \npTSAT.

\npdato{Lo que no sirve} agregar como cuarto término el opuesto de uno de los tres (repite variable: cláusula inválida, y además toda instancia sería \npsi), o agregar una fórmula en lugar de un término.
\end{problema}

\subsection{Desde camino o ciclo hamiltoniano}

La señal es <<recorrer todo>> o <<un camino o ciclo de largo al menos $k$>>. La reducción típica fija el parámetro para que <<largo $\ge k$>> obligue a pasar por todos los vértices.

El único ejercicio de este tipo daba camino hamiltoniano como $\NP$-completo y pedía la reducción y la pertenencia a $\NP$ por separado.\visto{2P 2021 Ej. 2: camino simple, desde CH}

\begin{problema}{Camino simple de largo al menos $k$ (CS)}
\npdato{Problema} dado un grafo $G$ y un natural $k$, ¿hay un camino simple (sin vértices repetidos) con al menos $k$ aristas?

\npdato{Reducción $\prob{ch}\redp\prob{cs}$} $G\mapsto(G,\,k=n-1)$. Copiar $G$ es lineal.

\npdato{Ida} un camino hamiltoniano no repite vértices y visita los $n$: es simple de largo $n-1$. \npdato{Vuelta} un camino simple de largo $n-1$ visita $n$ vértices distintos, o sea todos: es hamiltoniano.

\npdato{En $\NP$} certificado: lista de vértices $\pi$. $B$ rechaza si $\pi$ repite un vértice, si tiene $k$ o menos vértices, o si dos consecutivos no son adyacentes. Tiempo $\Oh(|\pi|^2+|\pi|m)$, y $|\pi|\le n$ porque no repite.
\end{problema}

Variantes que se resuelven igual: <<¿hay un ciclo de largo al menos $k$?>> (desde ciclo hamiltoniano con $k=n$), <<¿hay un camino hamiltoniano entre $s$ y $t$ dados?>>, o un \npTSP\ con restricciones extra (desde ciclo hamiltoniano con distancias 1 y 2).

\subsection{Desde Subset Sum, 3-coloreo y 3D-Matching}

En los parciales y exámenes relevados no se tomó ninguno con estos orígenes, pero están en el programa y son la elección natural para problemas numéricos o de etiquetado. Tres ejemplos cortos para practicar:

\begin{problema}{Partición (desde Subset Sum)}
\npdato{Problema} dados naturales $w_1,\dots,w_n$, ¿se pueden repartir en dos grupos de igual suma?

\npdato{Reducción} sea $S=\sum w_i$. Dada $(w,W)$ de \npSS, se agrega el número $\delta=|S-2W|$. Si $W\le S/2$, el total pasa a $2S-2W$ y cada mitad debe sumar $S-W$: un subconjunto que suma $W$ junto con $\delta$ suma $S-W$, y recíprocamente el lado que contiene a $\delta$, sin $\delta$, suma $W$. Si $W>S/2$, el total es $2W$ y el lado que no contiene a $\delta$ suma $W$. El número agregado tiene $\Oh(\log S)$ bits.
\end{problema}

\begin{problema}{4-coloreo (desde 3-coloreo)}
Se agrega un vértice nuevo adyacente a todos. En un 4-coloreo ese vértice tiene un color que no usa nadie más, así que el resto usa 3 colores; recíprocamente, un 3-coloreo se extiende dándole el cuarto color al vértice nuevo. En $\NP$: certificado, el color de cada vértice; se chequea cada arista.
\end{problema}

\begin{problema}{Mochila con valor exacto (desde Subset Sum)}
<<¿Hay un subconjunto de peso $\le W$ y valor exactamente $V$?>> Desde \npSS\ con $v_i=w_i$ y $V=W$. La lección general: si el problema nuevo tiene números y pide <<exactamente>>, \npSS\ casi siempre es el origen, y los números se copian (no se escriben en unario).
\end{problema}

% =====================================================================
\section{Mapa de lo que se preguntó}
% =====================================================================

Relevamiento de los segundos parciales (2008--2025) y exámenes (2009--2026) disponibles. Hasta 2016 los parciales y exámenes no incluían este tema; desde 2018 aparece en todos los segundos parciales y en buena parte de los exámenes, en general como el último ejercicio y con 17 a 40 puntos.

\begin{center}\small
\begin{tabularx}{\textwidth}{@{}>{\raggedright\arraybackslash}p{2.1cm} >{\raggedright\arraybackslash}X >{\raggedright\arraybackslash}p{3.3cm}@{}}
\toprule
\multicolumn{3}{@{}l}{\sans\bfseries Segundos parciales}\\
\midrule
Fuente & Qué se pidió & Reducido desde \\
\midrule
2P 2017 Ej. 4 & sí, no o no se sabe: Interval Scheduling y VC, IS e Interval Scheduling & conceptual \\
2P 2018 Ej. 3 & UOPI (independiente con pesos, valor $\ge t$) en $\NP$ y $\NP$-completo & \npIS \\
2P 2019 Ej. 3 & Solitaire: vuelta de una transformación dada, reducción y $\NP$-completo & \npTSAT \\
2P 2020 Ej. 3 & casiSAT (satisfacer al menos $k-1$ cláusulas) & \npSAT \\
2P 2021 Ej. 2 & camino simple de largo $\ge k$ & camino hamiltoniano \\
2P 2022 Ej. 3 & Contratar (aplicantes que cubren tareas) & \npVC \\
2P 2023 Ej. 3 & Acceso (puntos de acceso que cubren dispositivos) & \npSC \\
2P 2024 Ej. 2 & 4-SAT, con pauta de corrección detallada & \npTSAT \\
2P 2025 Ej. 1 & (e) transformación IS a VC; (f) definición de certificador eficiente & conceptual \\
2P 2025 Ej. 2 & UV (caminos cuya unión es $V$), guiado paso a paso & \npSC \\
\bottomrule
\end{tabularx}
\end{center}

\begin{center}\small
\begin{tabularx}{\textwidth}{@{}>{\raggedright\arraybackslash}p{2.1cm} >{\raggedright\arraybackslash}X >{\raggedright\arraybackslash}p{3.3cm}@{}}
\toprule
\multicolumn{3}{@{}l}{\sans\bfseries Exámenes}\\
\midrule
Fuente & Qué se pidió & Reducido desde \\
\midrule
Jul. 2018 P. 2 & $X\in\Pclass$, $Y\in\NP$, $Z$ $\NP$-completo: cuatro reducciones, ¿siempre, nunca o depende? & conceptual \\
Feb. 2019 Ej. 3 & definiciones de certificador eficiente, $\NP$ y $\NP$-completo & conceptual \\
Feb. 2020 Ej. 1 & PCG (conjunto golpeador) y definición de $\NP$ & \npVC \\
Feb. 2021 Ej. 3 & Red-Social (personas todas amigas entre sí) & \npIS\ (complemento) \\
Jul. 2021 Ej. 3 & Test (tests que ejecutan todas las líneas) & \npVC \\
Feb. 2022 Ej. 2 & Falsable (fórmula en forma disyuntiva que se puede hacer falsa) & \npSAT \\
Jul. 2022 Ej. 3 & Cocinero (ingredientes con disconformidad $\le d$) & \npIS\ ($d=0$) \\
Dic. 2022 Ej. 3 & RF (relación cuya restricción es función parcial) & \npIS \\
Dic. 2023 Ej. 3 & Cometa (clique de $d$ más camino de $d$) & \prob{Clique3} \\
Jul. 2024 Ej. 3 & subgrafo isomorfo & \npCL \\
Jul. 2025 Ej. 3 & Al menos $k$ (apostadores que aciertan) & \npSAT \\
Dic. 2025 Ej. 1 & (e) certificado de SAT; (f) definición de $\NP$-completo y cómo usarla & conceptual \\
Jul. 2026 Ej. 1 & (f) certificado de SAT & conceptual \\
Jul. 2026 Ej. 3 & NTI4Sat (no todos iguales, 4 términos) & \npTSAT \\
\bottomrule
\end{tabularx}
\end{center}

Lo que se repite: el origen es \npIS, \npVC/\npSC\ o \npSAT/\npTSAT\ en 15 de los 18 ejercicios de reducción (los otros tres salen de \npCL\ o de camino hamiltoniano); la transformación es casi siempre una copia con otro nombre; y lo que más puntaje cuesta es la redacción completa del certificador y la vuelta de la equivalencia.


\end{document}
